\documentclass[11pt]{amsart}

\usepackage[margin=1.15in]{geometry}
\usepackage{amsmath,amssymb,amsthm}
\usepackage{mathtools}
\usepackage{booktabs}
\usepackage[hidelinks]{hyperref}
\hypersetup{pdftitle={Bounds for the Length of Finite Pierce Expansions},pdfauthor={Jason Davies}}
\usepackage{xcolor}
\emergencystretch=2em
\raggedbottom

\newtheorem{theorem}{Theorem}[section]
\newtheorem{lemma}[theorem]{Lemma}
\newtheorem{proposition}[theorem]{Proposition}
\newtheorem{corollary}[theorem]{Corollary}
\theoremstyle{definition}
\newtheorem{definition}[theorem]{Definition}
\theoremstyle{remark}
\newtheorem{remark}[theorem]{Remark}

\newcommand{\lean}[1]{\texttt{#1}}

\title[The length of finite Pierce expansions]{Bounds for the Length of Finite Pierce Expansions}

\author{Jason Davies}
\email{jason@jasondavies.com}

\date{October 2026}

\begin{document}

\begin{abstract}
For integers $a>b>0$, let $P(a,b)$ count the positive terms of the iteration
$x\mapsto a\bmod x$, starting at $b$. We prove
\[
 P(a,b)\ll_\varepsilon a^{\beta_*+\varepsilon},\qquad
 \beta_*=\frac{41301-3\sqrt{82849}}{137480}=0.2941336553\ldots.
\]
Positive integer Hankel determinants and related minors group short orbit
segments into a small number of patterns. The fixed-dividend equations place the quotient increments
for each pattern on affine curves. We improve the resulting integer-point
count by adjoining further integer invariants, reducing their lattice
coordinates, and combining an effective Hilbert irreducibility sieve with an
exact polar identity. We also prove an elementary $3/10+\varepsilon$ bound and
a bound for orbits with few rises in their quotient increments.
The proofs, including the curve-counting input and the required effective
Hilbert irreducibility theorem, are formalized in Lean~4 with Mathlib without
proof holes or additional mathematical hypotheses. Shallit's conjectured
$O((\log a)^2)$ bound remains open.
\end{abstract}

\maketitle
\enlargethispage{1.1\baselineskip}

\section{Introduction}\label{sec:intro}

Let $a>b>0$ be integers. Put $b_0=b$ and $b_{i+1}=a\bmod b_i$ while $b_i>0$.
The sequence decreases until it reaches $0$; if $b_n=0$ we write $P(a,b)=n$.
For example, $P(35,22)=7$, from $22,13,9,8,3,2,1,0$. This is the length of
the Pierce expansion of $b/a$~\cite{Pierce,Shallit86}, since
\[
 \frac ba=\frac1{q_0}-\frac1{q_0q_1}+\frac1{q_0q_1q_2}-\cdots,
 \qquad q_i=\lfloor a/b_i\rfloor.
\]
Shallit~\cite{OPG} asked whether
\begin{equation}\label{eq:conj}
 P(a,b)=O\bigl((\log a)^2\bigr).
\end{equation}

Shallit~\cite{Shallit86} proved $P(a,b)\ll a^{1/2}$ by ``Archimedean'' arguments.
Erd\H{o}s and Shallit~\cite{ES} combined these with a divisor argument to obtain
$P(a,b)\ll_\varepsilon a^{1/3+\varepsilon}$, and showed that
$\max_bP(a,b)/\log a$ does not tend to $0$. Chase and Pandey~\cite{CP} improved
the upper exponent to $1/3-2/177$, using exponential sum estimates, and proved
$\max_bP(a,b)\gg\log a/\log\log a$ for all large $a$.
Baraskar and Vukusic~\cite{BV} study the related sets of iterated remainders
at fixed iteration depth. Shallit's talk pages~\cite{ShallitTalks} report a
bound $n^{3/10}$, described as ``LLM, unpublished so far''. We have not seen
that argument; our elementary bound below has the same exponent up to
$\varepsilon$, and we make no claim about how the arguments are related.

Our main result is the following uniform bound.
\begin{theorem}\label{thm:polar}
For every $\varepsilon>0$ there is $C_\varepsilon>0$, independent of $a,b$, such that
\[
 P(a,b)\le C_\varepsilon a^{\beta_*+\varepsilon}\qquad(0<b<a),\qquad
 \beta_*=\frac{41301-3\sqrt{82849}}{137480}.
\]
The fixed powers $(m,n)=(3,2)$ in the proof give the rational exponent
$4317/14677$.
\end{theorem}
Here $\beta_*=0.2941336553\ldots$, while $4317/14677=0.2941336785\ldots$.
The argument combines small integer invariants of the orbit with
coefficient-uniform counting on curves. Its final improvement comes from a
specialization sieve and a linear description of the curvature direction.

\subsection*{From integer invariants to curves}
Write $k_i=q_{i+1}-q_i$ for the quotient increments and
$d_i=b_i-b_{i+1}$ for the drops. Four consecutive terms $x,y,z,w$ have drops
$d,e,f$ and two consecutive quotient increments $u,v$. Their curvature
$D=dv-eu$ satisfies
\begin{equation}\label{eq:intro-curv}
 Dz=df+ude+(u-1)e^2.
\end{equation}
Thus $D$ is a positive integer, although its defining terms largely cancel,
and it is small when the drops are small. A six-term identity supplies a
second positive integer $L$ and controls the quotient by the small
parameters of a single block (Section~\ref{sec:six}).

These integers belong to one hierarchy. Each tail carries a finite positive
measure $\nu_i$ of mass $b_i/a$, and
\[
 I_r(i)=a\det\left(\int x^{j+l}\,d\nu_i\right)_{j,l<r}
\]
is a positive integer when the tail has at least $r$ positive terms.
The first four invariants are $b_i,d_i,D_i,L_i$, in the nonterminal ranges
of Theorem~\ref{thm:hankel}.

We count the $T$ orbit terms in a remainder window $[A,2A)$, then sum over
dyadic $A$. A fixed proportion of starts have their next few increments
bounded by $H\asymp a/(AT)$. The ten maximal minors of a $3\times5$ matrix
of drops and increments, including $I_4(i),I_4(i+1),I_4(i+2)$, then have
absolute value at most $X\asymp A^4H^6/a^3$. Their vector is the
\emph{pattern} of the start. The Pl\"ucker relations reduce the number of
patterns to $a^{o(1)}X^5$ (Section~\ref{sec:patterns}).

For a fixed pattern, three nodes $s_j=q_{i+j}-q_i$, $j=1,2,3$, determine the
increments. The nodes satisfy a quadratic equation and, on retaining the
fixed dividend $a$, a second equation of degree at most twelve. They
therefore lie on a union of curves of total degree at most $24$.
We adjoin three integer curvature minors $W=(W_0,W_1,W_2)$ and the fifth
Hankel invariant $J=I_5(i)$, whose sizes satisfy
\[
 |W_j|\ll S:=A^3H^3/a^2,\qquad |J|\ll Y:=A^5H^{10}/a^4.
\]
The zeros and poles of these functions give lower bounds for the degree of
the graph obtained by adjoining $W_0^m,W_1^m,W_2^m,J^n$, for fixed positive
integers $m,n$. This makes the coefficient-uniform curve bound of Castryck,
Cluckers, Dittmann and Nguyen~\cite{CCDN} effective: an integral curve over
$\mathbb Q$ of fixed degree $\delta$ in a box of size $B$ has at most
$B^{1/\delta+o(1)}$ integer points.
Conics and components on which the reconstruction fails are counted separately.

The pattern also determines a rank-three lattice. If $\lambda$ is the
smallest sup norm of a nonzero lattice vector, a reduced basis gives integer
gap and curvature coordinates of sizes $O(H/\lambda)$ and
$O(S/\lambda^2)$. Averaged bounds for short lattice vectors improve the
sum over patterns without changing the graph degree
(Section~\ref{sec:lattice}).

\subsection*{The polar improvement}
Two further facts give Theorem~\ref{thm:polar}. First, an effective Hilbert
irreducibility sieve forces the relevant curve component to have degree
$24$ outside a sparse set of patterns. Second, the exact polar identity
identifies the projective direction $[W]$ with three affine-linear forms
in the gaps. On the normalized curve, the common zeros of these
forms, counted with multiplicity and residue degree, contribute at most four, outside a fixed proper algebraic
exceptional locus. For the powered graph this gives the lower bound
\[
 \deg\Gamma\ge\frac{5n+3m}{2}\,24+m(24-4)=60n+56m.
\]
The sieve and the geometric exceptional locus together discard at most
$O_\eta(a^\eta X^{14/3})$ patterns, which are handled by the earlier curve
estimate. Applying the curve bound in the reduced lattice coordinates and
optimizing the fixed ratio $n/m$ gives $\beta_*$
(Section~\ref{sec:polar}). This completes the route from the orbit to the
main exponent.

\subsection*{Intermediate bounds}
The two curve-counting stages before the polar improvement give useful
bounds without the specialization sieve. Put
\[
 \gamma_* = \frac{471}{1559}-\frac{6\sqrt{105}}{7795}
          =0.2942294\ldots,\qquad
 \alpha_* = \frac{1310-\sqrt{655}}{4365}=0.2942513\ldots.
\]
\begin{theorem}[Exterior lattice bound]\label{thm:lattice}
For every $\varepsilon>0$, $P(a,b)\ll_\varepsilon a^{\gamma_*+\varepsilon}$
uniformly for $0<b<a$.
\end{theorem}
\begin{theorem}[Joint-invariant bound]\label{thm:joint}
For every $\varepsilon>0$, $P(a,b)\ll_\varepsilon a^{\alpha_*+\varepsilon}$
uniformly for $0<b<a$.
\end{theorem}
These stages fit into the progression below. The first four bounds are
formalized in the repository, but their proofs are not included here.
\begin{center}\small
\begin{tabular}{@{}ll@{}}
\toprule
Exponent & Main input\\
\midrule
$23/77$ & $I_4$ as an integer label\\
$11/37$ & a ternary quadratic and norm equations on slices\\
$8/27$ & the $X^{5+o(1)}$ count of minor vectors\\
$45/152$ & lattices of pivot gaps and their short vectors\\
$\alpha_*$ & adjoining curvature and rank-five invariants\\
$\gamma_*$ & reducing the joint coordinates in a lattice basis\\
\bottomrule
\end{tabular}
\end{center}
Fixed powers give the rational instances $128/435$ and $2201/7480$ of the
joint-invariant bound, and $153/520$ of the exterior lattice bound; see
Corollaries~\ref{cor:jointexponent} and~\ref{cor:latticeexponent}.
The repository also proves $P(a,b)\ll a^{3/10}\sqrt{\log\log3a}$ conditional
on the short-sum estimate of Nair and Tenenbaum~\cite{NT} and the mean-value
estimate of de la Bret\`eche and Tenenbaum~\cite{BT} for the Erd\H{o}s--Hooley
function. These analytic estimates are explicit hypotheses; their proofs
are not part of this formalization.

\subsection*{Elementary and structured bounds}
The initial identities already yield an elementary improvement over the
exponent $1/3$.
\begin{theorem}\label{thm:main}
For every $\varepsilon>0$, $P(a,b)\ll_\varepsilon a^{3/10+\varepsilon}$
uniformly for $0<b<a$.
\end{theorem}
The proof uses only division equations and the divisor estimate
$\tau(n)\ll_\eta n^\eta$. For a window of $T$ remainders in $[A,2A)$,
\begin{align}
 T(T-1)&\le4A\Delta,\label{eq:intro1}\\
 T^6A^2&\ll a^3\Delta,\label{eq:intro2}
\end{align}
where $\Delta=\max_{1\le n\le32a^7}\tau(n)$. The first is the
Erd\H{o}s--Shallit drop argument. For the second, there are few small triples
$(u,v,D)$, and a divisor argument bounds how often a triple can occur.
Multiplying the second bound by the square of the first eliminates $A$
and gives $T^{10}\ll a^3\Delta^3$.
The six-step identity further gives
\begin{equation}\label{eq:intro3}
 T^6A^2\ll a^3,
\end{equation}
without a divisor loss. This sharper estimate is used in the later arguments.

There is also a bound sensitive to the order of the increments. Let
\[
 R=1+\#\{i:0\le i\le P(a,b)-3,\ k_i<k_{i+1}\}.
\]
\begin{theorem}\label{thm:rises}
$P(a,b)\ll a^{1/4}\sqrt R$ uniformly for $0<b<a$. In particular,
nonincreasing quotient increments give $P(a,b)\ll a^{1/4}$.
\end{theorem}
Curvature positivity makes the drops decrease between rises. The six-step
identity limits repeated equal increments, and summing the resulting
window bounds proves the theorem. This argument is independent of the
curve-counting refinements.

\subsection*{Reading guide and formal verification}
Sections~\ref{sec:curv}--\ref{sec:six} establish the elementary identities and
bound. Sections~\ref{sec:hankel}--\ref{sec:patterns} construct the integer
invariants and their patterns. Section~\ref{sec:joint} builds the enlarged
curves, Section~\ref{sec:lattice} reduces their lattice coordinates, and
Section~\ref{sec:polar} completes the main theorem. The independent rise
bound is proved in Section~\ref{sec:rises}. The conic count, graph-degree
upper bound, and three-dimensional lattice reduction are proved in
Appendices~\ref{app:conic}, \ref{app:degree} and~\ref{app:reduction}.

All the unconditional results above, including their curve-counting and
specialization inputs, are formalized in Lean~4~\cite{Lean} with
Mathlib~\cite{Mathlib}. Section~\ref{sec:lean} identifies the proof artifact
and its verification scope. Section~\ref{sec:limits} explains the present
$5/17$ barrier; Shallit's logarithmic conjecture remains open.

\section{Preliminaries}

\subsection{Notation}
Throughout, $a>b>0$ are integers, $b_0=b$, $b_{i+1}=a\bmod b_i$, and $P=P(a,b)$, so that
$a>b_0>b_1>\dots>b_{P-1}>0=b_P$. Write $q_i=\lfloor a/b_i\rfloor$, so
\[
a=q_ib_i+b_{i+1},\qquad 0\le b_{i+1}<b_i\qquad(0\le i<P).
\]
An index $i$ of the orbit is also called a \emph{state}; quantities ``at state $i$'' are
computed from $b_i,b_{i+1},\dots$ and $q_i,q_{i+1},\dots$, and are defined whenever the
orbit terms they use exist.
Let $\tau(n)$ be the number of positive divisors of $n$, and let
\[
\Delta=\Delta(a)=\max_{1\le n\le 32a^7}\tau(n).
\]
We write $F\ll_\eta G$ when $F\le C_\eta G$, with the constant uniform in the
dividend and orbit. A factor denoted $a^{o(1)}$ is bounded by
$C_\eta a^\eta$ for every fixed $\eta>0$; such factors include the divisor
losses below. We write $F\asymp G$ when both $F\ll G$ and $G\ll F$ hold.
Many statements below carry explicit numerical constants; these are the constants of the
Lean proof. Where a constant is large we state the bound with $\ll$ and record the constant
in the proof.

\begin{lemma}[Quotients increase]\label{lem:quot}
If $b_{i+1}>0$ then $q_{i+1}\ge q_i+1$.
\end{lemma}

\begin{proof}
Since $b_{i+1}<b_i$ we have $(q_i+1)b_{i+1}\le q_ib_i+b_{i+1}=a$, so
$q_{i+1}=\lfloor a/b_{i+1}\rfloor\ge q_i+1$.
\end{proof}

\begin{lemma}[Divisor bound]\label{lem:div}
For every integer $k\ge1$ there is $C_k$ with $\tau(n)^k\le C_k\,n$ for all $n\ge1$.
\end{lemma}

\begin{proof}
This is standard: $\tau(n)^k/n=\prod_{p^\alpha\| n}(\alpha+1)^k/p^\alpha$. Each factor with
$p\ge 2^k$ is at most $(\alpha+1)^k/2^{k\alpha}\le1$. Each of the fewer than $2^k$ factors with
$p<2^k$ is at most $\sup_\alpha(\alpha+1)^k/2^\alpha<\infty$.
\end{proof}

\subsection{Window runs}

\begin{definition}
Let $A\ge1$ be an integer. A \emph{window run} of length $T$ is a sequence $x_0,x_1,\dots,x_{T-1}$ of
consecutive terms of the sequence $(b_i)$, so $x_{j+1}=a\bmod x_j$, all lying in $[A,2A)$.
\end{definition}

In a window run, $x_0>x_1>\dots>x_{T-1}$ and $x_j\le a$.

\section{The curvature of three consecutive steps}\label{sec:curv}

Fix four consecutive terms
\[
x=b_i,\quad y=b_{i+1}=a\bmod x,\quad z=b_{i+2}=a\bmod y,\quad w=b_{i+3}=a\bmod z,
\]
with $x,y,z>0$ (and $w\ge0$). Write
\[
q=\left\lfloor \frac ax\right\rfloor,\quad
u=\left\lfloor \frac ay\right\rfloor-q,\quad
v=\left\lfloor \frac az\right\rfloor-\left\lfloor \frac ay\right\rfloor,\qquad
d=x-y,\quad e=y-z,\quad f=z-w .
\]
By Lemma~\ref{lem:quot}, $u,v\ge1$, and clearly $d,e,f\ge1$. The division equations read
\begin{equation}\label{eq:div}
xq+y=a,\qquad y(q+u)+z=a,\qquad z(q+u+v)+w=a .
\end{equation}

\begin{definition}
The \emph{curvature} of the block is $D=dv-eu$.
\end{definition}

\begin{lemma}[Curvature identity]\label{lem:curv}
\[
D\,z=df+ude+(u-1)e^2 .
\]
In particular $D\ge1$.
\end{lemma}

\begin{proof}
Subtracting the first two equations of~\eqref{eq:div} gives $uy=qd+e$, hence
$uz=uy-ue=qd+(1-u)e$. Subtracting the last two gives $vz=(q+u)e+f$. Therefore
\[
(dv-eu)z=d\bigl((q+u)e+f\bigr)-e\bigl(qd+(1-u)e\bigr)=df+ude+(u-1)e^2 .
\]
The right-hand side is at least $df\ge1$, and $z>0$, so $D>0$.
\end{proof}

\begin{lemma}[The fiber equation]\label{lem:fiber-eq}
Put $M(q)=uq^2+(u^2+u+v)q+v$ and $N(q)=uq+u^2+v$. Then
\[
x\,M(q)=a\,N(q)+D .
\]
\end{lemma}

\begin{proof}
From~\eqref{eq:div}, $d=(q+1)x-a$ and $e=y-z=(q+u+1)y-a=a(q+u)-q(q+u+1)x$. Substituting
these into $D=dv-eu$ gives
$D=x\bigl(uq(q+u+1)+v(q+1)\bigr)-a\bigl(u(q+u)+v\bigr)$, which is the claim.
\end{proof}

\begin{lemma}[Size of the parameters]\label{lem:size}
If $x\le a$ then
\[
u\,yx\le a(d+e),\qquad v\,zy\le a(e+f),\qquad D\,z\le (u+1)(d+e+f)^2 .
\]
\end{lemma}

\begin{proof}
From $uy=qd+e$ we get $uyx=(qx)d+ex\le ad+ae$, since $qx\le a$ and $x\le a$. Similarly
$vz=(q+u)e+f$ with $(q+u)y\le a$ and $y\le a$ gives the second bound. For the third, by
Lemma~\ref{lem:curv},
$Dz=df+ude+(u-1)e^2\le df+u\,e(d+e)\le (1+u)(d+e+f)^2$.
\end{proof}

\section{Proof of Theorem~\ref{thm:main}}\label{sec:proof310}

\subsection{Few blocks share a triple \texorpdfstring{$(u,v,D)$}{(u,v,D)}}

\begin{lemma}[Fiber count]\label{lem:fiber}
Let $a\ge1$, $1\le u,v\le a$ and $1\le D\le 2a$. Let $M(q)=uq^2+pq+v$ with $p=u^2+u+v$, and
let $\Lambda(q)=cq+\ell$ with $c=au$ and $\ell=a(u^2+v)+D$. Then the set
\[
Q=\{q\in\mathbb Z_{\ge0}:\ xM(q)=\Lambda(q)\text{ for some integer }x\ge1\}
\]
has at most $\Delta=\max_{1\le n\le 32a^7}\tau(n)$ elements.
\end{lemma}

\begin{proof}
Put $U(q)=ucq+cp-u\ell$ and $R=u\ell^2-cp\ell+c^2v$. Expanding shows
\begin{equation}\label{eq:resultant}
c^2M(q)-U(q)\Lambda(q)=R
\end{equation}
identically in $q$. If $q\in Q$ with witness $x$, substituting $\Lambda(q)=xM(q)$
into~\eqref{eq:resultant} gives
\begin{equation}\label{eq:R}
R=M(q)\bigl(c^2-U(q)x\bigr).
\end{equation}
We bound the sizes first. From $u,v\le a$ and $D\le 2a$ we get $\ell\le 4a^3$, $p\le 3a^2$
and $c\le a^2$, so
\[
|R|\le u\ell^2+cp\ell+c^2v\le 16a^7+12a^7+a^5\le 32a^7,\qquad c^2\le a^4\le 32a^7 .
\]

\emph{Case $R\ne0$.} By~\eqref{eq:R}, $M(q)$ is a positive divisor of $|R|$. Since $u,p\ge1$,
the map $q\mapsto M(q)$ is strictly increasing on $q\ge0$, hence injective. Therefore
$|Q|\le\tau(|R|)\le\Delta$.

\emph{Case $R=0$.} Since $M(q)>0$, \eqref{eq:R} gives $U(q)x=c^2$. Here $c^2>0$ and $x\ge1$,
so $U(q)$ is a positive divisor of $c^2$. Since $U$ is linear with positive leading
coefficient $uc$, it is injective. Therefore $|Q|\le\tau(c^2)\le\Delta$.
\end{proof}

The case $R=0$ cannot be dropped: it is the case in which $M$ and $\Lambda$ might share a factor.

\subsection{Counting visits to a window}

In this subsection $x_0,\dots,x_{T-1}$ is a window run in $[A,2A)$, and
$q_j=\lfloor a/x_j\rfloor$ are its quotients. We continue to write
$\Delta=\max_{1\le n\le 32a^7}\tau(n)$.

\begin{proposition}[Drop bound; after Erd\H{o}s--Shallit]\label{prop:drop}
If $T\ge2$ then $T(T-1)\le 4A\Delta$.
\end{proposition}

\begin{proof}
Let $h_j=x_j-x_{j+1}\ge1$ for $0\le j\le T-2$. The total drop is
$\sum_jh_j=x_0-x_{T-1}\le A-1$. Thus fewer than $(T-1)/2$ drops exceed
$2A/(T-1)$. The remaining set $\mathcal S$ has $2|\mathcal S|\ge T$.

For $j\in\mathcal S$ we have $1\le h_j\le H:=\lfloor 2A/(T-1)\rfloor$. Also
$a+h_j=q_jx_j+x_{j+1}+x_j-x_{j+1}=(q_j+1)x_j$, so $x_j\mid a+h_j$. Moreover
$a+h_j\le 2a\le 32a^7$. Since $j\mapsto x_j$ is injective, at most $\tau(a+h)\le\Delta$ indices
in $\mathcal S$ have $h_j=h$. Hence $|\mathcal S|\le H\Delta$, and
\[
T(T-1)\le 2|\mathcal S|(T-1)\le 2\Delta\,H(T-1)\le 4A\Delta. \qedhere
\]
\end{proof}

\begin{proposition}[Curvature bound]\label{prop:curv}
If $T\ge12$ then $T^6A^2\le 995328\,a^3\Delta$.
\end{proposition}

\begin{proof}
For $0\le j\le T-4$ let $s_j=x_j-x_{j+3}$ be the total drop of the block
$x_j,x_{j+1},x_{j+2},x_{j+3}$. By telescoping,
\[
\sum_{j=0}^{T-4}s_j=(x_0+x_1+x_2)-(x_{T-3}+x_{T-2}+x_{T-1})\le 3(2A-1)-3A=3A-3 .
\]
Call a block small if $s_j\le12A/T$, and let $\mathcal S$ be the set of
small blocks. Fewer than $T/4$ of the $T-3$ blocks fail this condition,
so $|\mathcal S|\ge T/2$ when $T\ge12$.

Fix $j\in\mathcal S$. Apply Section~\ref{sec:curv} to $x=x_j$, so that $y,z,w=x_{j+1},x_{j+2},x_{j+3}$, all
in $[A,2A)$. Write $s=s_j=d+e+f$. Lemma~\ref{lem:size} gives $uA^2\le uyx\le a(d+e)\le as$.
Multiplying by $T$ and using $sT\le 12A$ gives $uA^2T\le 12aA$, that is,
\[
u\,AT\le 12a,\qquad\text{and likewise}\qquad v\,AT\le 12a .
\]
Also $DA\le Dz\le(u+1)s^2\le 2us^2$. Multiplying by $T^2$ gives
$DAT^2\le 2u(sT)^2\le 288uA^2$, so $DT^2\le 288uA$, and then
\[
D\,T^3\le 288\,u\,AT\le 3456\,a .
\]
Thus the triple $(u,v,D)$ of a small block lies in $[1,U]^2\times[1,D^*]$, where
\[
U=\lfloor 12a/(AT)\rfloor\le a,\qquad D^*=\lfloor 3456a/T^3\rfloor\le 2a .
\]
The inequalities $U\le a$ and $D^*\le 2a$ hold because $AT\ge12$ and $T^3\ge1728$.

Now fix such a triple and consider the small blocks that have it. By Lemma~\ref{lem:fiber-eq},
each one gives a point $q=q_j=\lfloor a/x_j\rfloor$ of the set $Q$ of Lemma~\ref{lem:fiber}.
The map $j\mapsto q_j$ is injective on these blocks: $q_j$ determines $x_j=\Lambda(q_j)/M(q_j)$, and
$x_j$ determines $j$. By Lemma~\ref{lem:fiber}, at most $\Delta$ small blocks share a triple.
Consequently
\[
\frac T2\le |\mathcal S|\le U^2D^*\Delta
\le\left(\frac{12a}{AT}\right)^2\frac{3456a}{T^3}\Delta,
\]
which is the claimed bound.
\end{proof}

\begin{corollary}[Window bound]\label{cor:window}
If $T\ge12$ then $T^{10}\le 63700992\,a^3\Delta^3$.
\end{corollary}

\begin{proof}
The drop bound and $T\le2(T-1)$ give $T^4\le64A^2\Delta^2$.
Multiply by the curvature bound and cancel $A^2$.
\end{proof}

\begin{remark}
In terms of $B=a/A$, Propositions~\ref{prop:drop} and~\ref{prop:curv} say
$T\ll a^{1/2+o(1)}B^{-1/2}$ and $T\ll a^{1/6+o(1)}B^{1/3}$. The first is better when
$B>a^{2/5}$, and both equal $a^{3/10+o(1)}$ at $B=a^{2/5}$. Corollary~\ref{cor:window} takes
the product of the two bounds, so no case split is needed.
\end{remark}

\subsection{Conclusion of the proof}

Let $0<b<a$. For $0\le k\le\log_2 a$ let $T_k$ be the number of indices $i<P$ with
$2^k\le b_i<2^{k+1}$. Every $b_i$ with $i<P$ lies in $[1,a)$, so
$P=\sum_{0\le k\le \log_2 a}T_k$. Since $(b_i)$ is decreasing, the indices counted by $T_k$
are consecutive. So when $T_k>0$ they form a window run in $[A,2A)$ with $A=2^k$.

For every $\eta>0$, Lemma~\ref{lem:div} gives $\Delta\ll_\eta a^\eta$,
since its defining maximum is over $n\le32a^7$. Taking $\eta=5\varepsilon/3$
in Corollary~\ref{cor:window} yields
$T_k\ll_\varepsilon a^{3/10+\varepsilon/2}$; the windows with $T_k<12$
satisfy the same estimate. There are $O(\log a)\ll_\varepsilon a^{\varepsilon/2}$
windows, so summing proves Theorem~\ref{thm:main}. \hfill$\square$

\section{Six consecutive steps}\label{sec:six}

Fix six consecutive terms $x_0>x_1>\dots>x_5\ge0$ of the sequence $(b_i)$ with $x_4>0$, and
write their quotients as
\begin{gather*}
\lfloor a/x_0\rfloor=q,\quad \lfloor a/x_1\rfloor=q+u,\quad \lfloor a/x_2\rfloor=q+u+v,\\
\lfloor a/x_3\rfloor=q+u+v+w,\quad \lfloor a/x_4\rfloor=q+u+v+w+z .
\end{gather*}
(Unlike in Section~\ref{sec:curv}, here $w$ and $z$ denote quotient increments.) By
Lemma~\ref{lem:quot}, $u,v,w,z\ge1$. Let $f_k=x_k-x_{k+1}\ge1$ be the drops, and let
\[
D=f_0v-f_1u,\qquad E=f_1w-f_2v,\qquad F=f_2z-f_3w
\]
be the curvatures of the blocks starting at $x_0$, $x_1$ and $x_2$. Write $\mathcal E_k$ for the division
equation $x_kq_k+x_{k+1}-a=0$, where $q_k=\lfloor a/x_k\rfloor$. Put
\begin{align*}
\sigma(u,v,w)&=u^2v+(u-1)v^2+uw,\\
\chi(u,v,w,z)&=\bigl(u^2+(u-1)v\bigr)z+w\bigl(u^3+uv(2u-1)+(u-1)v^2\bigr)+w^2\bigl((u-1)(u+v)+1\bigr),\\
L&=(vw+w^2+z)\,D-\bigl(u^2+(u-1)v\bigr)E-w\,(wf_0-uf_2).
\end{align*}

\begin{lemma}[The six-step identity]\label{lem:six}
\begin{enumerate}
\item $\sigma(v,w,z)\,D-\sigma(u,v,w)\,E=v\,L$.
\item $q\,w\,L=\chi(u,v,w,z)\,E+\sigma(u,v,w)\,F$.
\end{enumerate}
\end{lemma}

\begin{proof}
Part~(1) is a polynomial identity: since $\sigma(v,w,z)\equiv-w^2$ and $\sigma(u,v,w)\equiv uw\pmod v$,
and $wD+uE=v(wf_0-uf_2)$, the difference is divisible by $v$, and expanding gives $L$. Part~(2)
holds modulo the division equations. Explicitly, $qwL-\chi E-\sigma(u,v,w)F$ equals
\begin{align*}
&w(v^2w+vw^2+vz-w^2)\,\mathcal E_0-w(u^2w+2uvw+uw^2+uz+v^2w+vw^2-vw+vz-w^2)\,\mathcal E_1\\
&+w(u^2v+u^2w+uv^2+2uvw+uw^2+uw+uz-v^2-vw)\,\mathcal E_2-w(u^2v+uv^2+uw-v^2)\,\mathcal E_3 .
\end{align*}
This is an identity of polynomials in $a$, $q$, the $x_k$ and the increments, checked by expansion.
\end{proof}

\begin{corollary}\label{cor:six}
Suppose $q\ge1$ and $u,v,w,z\le m$. Then $E,F\ge1$, $L\ge1$, and
\[
q\le 9m^3E+3m^3F .
\]
\end{corollary}

\begin{proof}
Lemma~\ref{lem:curv}, applied to the blocks starting at $x_1$ and $x_2$, gives $E,F\ge1$.
Clearly $\chi>0$ and $\sigma\ge0$, so by Lemma~\ref{lem:six}(2) $qwL>0$. Since $L$ is an integer,
$L\ge1$ and $qw\le \chi E+\sigma F$. Bounding each monomial gives $\chi\le 9m^3w$ and $\sigma(u,v,w)\le 3m^3w$.
Dividing by $w$ gives the claim.
\end{proof}

So the quotient $q$, which is about $a/A$, is controlled by the small parameters of a single block
of six terms. No counting over many blocks is needed.

\begin{proposition}[Six-step window bound]\label{prop:six}
For a window run of length $T\ge10$,
\[
T^6A^2\le 96\cdot10^6\,a^3 .
\]
\end{proposition}

\begin{proof}
For $0\le j\le T-6$ let $s_j=x_j-x_{j+5}$. By telescoping,
$\sum_{j\le T-6}s_j\le 5(2A-1)-5A=5A-5$. So some $j$ has $s_j(T-5)\le 5A-5$, and since
$T\le 2(T-5)$, we get $s:=s_j$ with $sT\le 10A$. Apply Corollary~\ref{cor:six} to
$x_j,\dots,x_{j+5}$, all of which lie in $[A,2A)$, with $m=\max(u,v,w,z)$.

By Lemma~\ref{lem:size} at $x_j$, $x_{j+1}$ and $x_{j+2}$, each increment satisfies
$(\text{increment})\cdot A^2\le a s$. Multiplying by $T$ gives $mAT\le 10a$. By
Lemma~\ref{lem:size} again, $EA\le Ex_{j+3}\le(v+1)s^2\le 2ms^2$, and likewise $FA\le 2ms^2$.
Hence $qA\le 24m^4s^2$. Also $a=qx_j+x_{j+1}<(q+1)x_j<2A(q+1)\le 4Aq$. Therefore
\[
qA\cdot A^4T^6\le 24\,(mAT)^4(sT)^2\le 24\cdot10^6\,a^4A^2,
\]
so $qA^3T^6\le 24\cdot 10^6a^4$, and $aA^2T^6\le 4qA^3T^6\le 96\cdot10^6\,a^4$.
\end{proof}

\begin{remark}
In terms of $B=a/A$, Proposition~\ref{prop:six} reads $T\ll a^{1/6}B^{1/3}$. This is
Proposition~\ref{prop:curv} without the factor $\Delta^{1/6}$. Combined with
Proposition~\ref{prop:drop} as in Corollary~\ref{cor:window}, it gives $T^{10}\ll a^3\Delta^2$
and another proof of Theorem~\ref{thm:main}. The Lean proof of Theorem~\ref{thm:main} uses
Corollary~\ref{cor:window}.
\end{remark}

\section{The Hankel hierarchy}\label{sec:hankel}

From now on we write $N=P(a,b)$, so that
$b_0>\dots>b_{N-1}>0=b_N$, and $k_i=q_{i+1}-q_i$. For each remaining tail we construct a probability measure
on its quotient differences. Reweighting by the reciprocal quotient will
make its total mass equal to $b_i/a$.

Put $p_i(j)=q_{i+j}-q_i$ for $0\le j<N-i$; the symbols $p_i,w_i$
are local to this construction. Define the weights backwards from
$w_{N-1}(0)=1$ by
\[
w_i(0)=1-\sum_{j=0}^{N-i-2}\frac{w_{i+1}(j)}{p_i(j+1)},\qquad
w_i(j+1)=\frac{w_{i+1}(j)}{p_i(j+1)}
\quad(0\le j<N-i-1).
\]
Let $\mu_i$ be the measure with mass $w_i(j)$ at $p_i(j)$, and let
\[
 d\nu_i(x)=\frac{d\mu_i(x)}{q_i+x},\qquad
 h_i(s)=\int x^s\,d\nu_i=\sum_{j=0}^{N-i-1}
 \frac{w_i(j)p_i(j)^s}{q_{i+j}}.
\]
The recursion is chosen to reproduce the division equation
$q_ih_i(0)+h_{i+1}(0)=1$ for $i+1<N$.

\begin{lemma}[The moment model]\label{lem:moments}
\begin{enumerate}
\item The weights are positive and sum to $1$.
\item $h_i(0)=b_i/a$, that is, $b_i/a=\mathbb E\bigl[1/(q_i+X_i)\bigr]$ for $X_i\sim\mu_i$; and
$h_{N-1}(s)=0$ for $s\ge1$.
\item $h_i(s+1)=\sum_{t\le s}\binom st k_i^{s-t}h_{i+1}(t)$ for $i+1<N$.
\item $h_i(s)=(-q_i)^sb_i/a+z$ with $z\in\mathbb Z$.
\end{enumerate}
\end{lemma}

\begin{proof}
(1) Since $p_i(j+1)\ge1$, the sum defining $w_i(0)$ is at most $1$. It is strictly less: if
$N-i-1\ge2$, the term $j=1$ has $p_i(2)\ge2$; if $i=N-2$, the last gap satisfies
$q_{N-1}\ge q_{N-2}+2$, because $b_{N-1}\mid a$ gives
$q_{N-1}-q_{N-2}-1=q_{N-2}(b_{N-2}-b_{N-1})/b_{N-1}\ge1$. (2) Since $p_i(j+1)=q_{i+j+1}-q_i$, the definitions give
$q_ih_i(0)+h_{i+1}(0)=1$, which is the normalized division equation; and
$h_{N-1}(0)=1/q_{N-1}=b_{N-1}/a$; since $\mu_{N-1}=\delta_0$, also $h_{N-1}(s)=0$ for $s\ge1$.
(3) Use $p_i(j+1)=k_i+p_{i+1}(j)$ and the binomial theorem.
(4) For $i=N-1$ this holds with $z=-(-q_{N-1})^sb_{N-1}/a=-(-1)^sq_{N-1}^{s-1}$ for $s\ge1$, since
$q_{N-1}b_{N-1}=a$. Then induct backwards on $i$ using (3), $b_{i+1}/a=1-q_ib_i/a$ and
$k_i-q_{i+1}=-q_i$.
\end{proof}

The weights also explain the role of the Newton polynomials below. Write
$\mu_i(g)=\int g\,d\mu_i$. For $i+1<N$ and any polynomial $g$, the zero
atom contributes nothing to $\mu_i(xg)$, and
$p_i(j+1)w_i(j+1)=w_{i+1}(j)$ gives
\[
 \mu_i(xg(x))=\mu_{i+1}\bigl(g(k_i+x)\bigr).
\]
Put $s_j=q_{i+j}-q_i$ and $F_j(x)=\prod_{0\le l<j}(x-s_l)$, with
$F_0=1$. Iterating this identity gives the \emph{shift rule}
\begin{equation}\label{eq:moment-shift}
 \mu_i(F_j(x)g(x))=\mu_{i+j}\bigl(g(x+s_j)\bigr)
 \qquad(i+j<N).
\end{equation}
Indeed, each step removes the first factor of $F_j$ and translates the
remaining nodes by the next quotient gap. Thus multiplication by a Newton
polynomial advances along the orbit; this is why this basis is useful for
bounding its Hankel determinants.

\begin{theorem}[Hankel determinants]\label{thm:hankel}
For $1\le r\le N-i$, the number $I_r(i)=a\det\bigl(h_i(j+l)\bigr)_{0\le j,l<r}$ is a positive
integer. Moreover $I_1(i)=b_i$, $I_2(i)=d_i$ (if $i+1<N$), $I_3(i)=D_i$ (if $i+3<N$), and
$I_4(i)=L_i$ (if $i+5<N$), where $D_i$ and $L_i$ are the curvature and the six-step integer
of the blocks starting at $b_i$.
\end{theorem}

\begin{proof}
\emph{Integrality.} By Lemma~\ref{lem:moments}(4),
$h_i(s+1)+q_ih_i(s)\in\mathbb Z$. Multiply the Hankel matrix on the left by the unit lower
triangular matrix that adds $q_i$ times row $j-1$ to row $j$. This makes rows
$1,\dots,r-1$ integral and leaves the determinant unchanged. Expand along row $0$, whose entries
are $(-q_i)^lb_i/a+(\text{integer})$; then $a\det$ is an integer.

\emph{Positivity.} The Hankel matrix is the Gram matrix of $1,x,\dots,x^{r-1}$ in $L^2(\nu_i)$.
Its quadratic form is $\sum_jw_i(j)q_{i+j}^{-1}\bigl(\sum_lc_lp_i(j)^l\bigr)^2$. If this
vanishes, the polynomial $\sum_lc_lx^l$ vanishes at the $N-i\ge r$ distinct points $p_i(j)$, so
it is zero, by the invertibility of the Vandermonde matrix. So the matrix is positive definite.

\emph{The first four.} First, $I_1(i)=ah_i(0)=b_i$.
For $i+1<N$, write $\beta_j=b_{i+j}/a$. Then
$h_i(0)=\beta_0$, $h_i(1)=\beta_1$ and
$h_i(2)=k_i\beta_1+\beta_2$. Hence
\[
 I_2(i)=a\bigl(\beta_0(k_i\beta_1+\beta_2)-\beta_1^2\bigr)
       =a(\beta_0-\beta_1)=d_i,
\]
using $q_i\beta_0+\beta_1=q_{i+1}\beta_1+\beta_2=1$.
For the larger determinants, Lemma~\ref{lem:moments}(2),(3) makes the entries $h_i(s)$, $s\le2r-2$,
explicit polynomials in $\beta_j=b_{i+j}/a$ and $k_i,\dots,k_{i+2r-4}$; at the end of the orbit
this uses $h_{N-1}(s)=0$ for $s\ge1$, that is, $b_N=0$.
Expanding and reducing by $q_{i+j}\beta_j+\beta_{j+1}=1$ gives $D_i$
and $L_i$ in the stated ranges. The formal proof records polynomial
certificates for these two reductions.
\end{proof}

The ranges $i+3<N$ and $i+5<N$ cannot be relaxed to all $i$ at which $D_i$ or $L_i$ is defined.
For $a=926$ and $b=229$ the orbit is $229,10,6,2$, so $N=4$, and $I_3(1)=984$ while $D_1=988$.

\begin{remark}
In these index ranges, the positivity of $D$ and $L$, proved earlier by explicit identities, is an instance of the
positivity of Gram determinants. The integrality of $L$ (Lemma~\ref{lem:six}(1)) is an instance
of the integrality of $I_4$. A priori $I_r(i)$ involves the moments $h_i(s)$ for $s\le2r-2$,
and hence $b_i,\dots,b_{i+2r-2}$. In fact $I_3(i)$ depends only on $b_i,b_{i+1},b_{i+2}$ and
$q_i,q_{i+1},q_{i+2}$. Similarly $I_4(i)$ depends only on $b_i,\dots,b_{i+3}$ and
$q_i,\dots,q_{i+4}$: the later remainders cancel.
\end{remark}

\section{Patterns}\label{sec:patterns}

In this section we develop the patterns used in Section~\ref{sec:joint}, together with the
bounds on their size and number. We keep the notation of
Section~\ref{sec:hankel}: $b_0>b_1>\dots>b_{N-1}>0$ is the positive part of the orbit,
$N=P(a,b)$, and $q_i=\lfloor a/b_i\rfloor$.

\subsection{Heuristics}
The exponent $3/10$ of Theorem~\ref{thm:main} comes from a single balanced configuration:
$T\asymp a^{3/10}$ remainders in a window $[A,2A)$ with $A\asymp a^{3/5}$, quotients
$q\asymp a^{2/5}$ and typical increments $k\asymp a^{1/10}$. At this scale the four-term Hankel
invariant $I_4(i)=L(i)$ has size $\asymp aH^6/q^4\asymp 1$, where $H$ bounds the increments.
The idea is that a positive integer of bounded size carries information. The labels
$I_4(i),I_4(i+1),I_4(i+2)$, together with four further integers of the same size, form a
\emph{pattern} (made precise in Section~\ref{sec:minors}). States with a given pattern satisfy two linear equations and one quadratic
equation in five consecutive increments (Lemma~\ref{lem:pattern}), and the quotient $q_i$ is a
root of a quartic fixed by the pattern and the increments (Lemma~\ref{lem:multid}). So the
problem becomes one of counting patterns, and the states belonging to each pattern. Counting the
latter with these equations alone already improves on $3/10$ (see the intermediate bounds in
Section~\ref{sec:intro}). Section~\ref{sec:joint} instead keeps the fixed dividend $a$ in the
equations.

\subsection{The pattern matrix}
Fix a state $i$ with $i+6<N$ and write, for $m\ge0$,
\begin{gather*}
k_m=q_{i+m+1}-q_{i+m},\qquad d_m=b_{i+m}-b_{i+m+1},\\
s_m=k_0+\dots+k_{m-1},\qquad r_m=k_m(s_m+s_{m+1})+k_{m+1}.
\end{gather*}
All indices here are local to the chosen state, and the definitions apply
whenever the indicated orbit entries exist; the empty sum gives $s_0=0$.
So $s_m=q_{i+m}-q_i$, $r_0=k_0^2+k_1$, $r_1=k_1(2k_0+k_1)+k_2$, and so on. Consider the
$3\times4$ integer matrix
\[
\mathcal M_i=\begin{pmatrix} d_0&d_1&d_2&d_3\\ k_0&k_1&k_2&k_3\\ r_0&r_1&r_2&r_3\end{pmatrix},
\]
and let $U_m(i)=(-1)^{m+1}\det(\mathcal M_i\text{ with column $m$ deleted})$ for
$0\le m\le3$.

\begin{lemma}[Cofactor relations]\label{lem:pattern}
\begin{enumerate}
\item $U_3(i)=L(i)=I_4(i)$ and $U_0(i)=-L(i+1)=-I_4(i+1)$, where $L$ is the six-step integer of
Lemma~\ref{lem:six}.
\item $\sum_m U_m(i)\,d_m=\sum_m U_m(i)\,k_m=\sum_m U_m(i)\,r_m=0$.
\end{enumerate}
\end{lemma}

\begin{proof}
(2) is the Laplace expansion of a $4\times4$ determinant with a repeated row. For (1), the
identity $\det(\text{columns }0,1,2)=L(i)$ holds as a polynomial identity in
$d_0,d_1,d_2,k_0,\dots,k_3$, by direct expansion; and $L(i)=I_4(i)$ by
Theorem~\ref{thm:hankel}. Passing from $i$ to $i+1$ shifts the $d$- and $k$-rows by one
column, and, since $s_m(i+1)=s_{m+1}(i)-k_0$, it replaces $r_m$ by $r_{m+1}-2k_0k_{m+1}$. This
row operation does not change the determinant, so $\det(\text{columns }1,2,3)=L(i+1)$.
\end{proof}

For example, the positive part of the orbit with $(a,b)=(95,61)$ is
$61,34,27,14,11,7,4,3,2,1$. Its first pattern matrix and cofactor vector are
\[
 \mathcal M_0=\begin{pmatrix}27&7&13&3\\1&1&3&2\\2&6&23&29\end{pmatrix},
 \qquad U(0)=(-81,833,-296,68).
\]
Thus $L(0)=68$, $L(1)=81$, and each row of $\mathcal M_0$ is orthogonal
to $U(0)$.

\begin{remark}
The vector $C_m=U_m/U_3$ has a meaning: $\pi_i=\sum_{m\le3}C_m(F_{m+1}-F_m)$, with Newton
polynomials $F_m(x)=\prod_{l<m}(x-s_l)$, is the monic orthogonal polynomial of degree four for
the measure $\nu_i$ of Section~\ref{sec:hankel}. The three rows of $\mathcal M_i$ encode
$\nu_i(\pi_i)=0$ and $\mu_i(x\pi_i)=\mu_i(x^2\pi_i)=0$, and $\nu_i(xf)=\mu_i(f)-q_i\nu_i(f)$ turns these
into $\nu_i(x^j\pi_i)=0$ for $j\le3$. In particular $U_m(i)/I_4(i)$ are tail sums of the
Newton coefficients of $\pi_i$: writing $\pi_i=\sum_{l\le4}\gamma_lF_l$, one has $U_m(i)/I_4(i)=\gamma_{m+1}+\dots+\gamma_4$,
and $\gamma_0=-C_0=I_4(i+1)/I_4(i)$. This interpretation is used in Section~\ref{sec:joint}.
\end{remark}

The \emph{pattern integers} of $i$ are the seven integers
\[
\theta(i)=\bigl(L(i),\,L(i+1),\,L(i+2),\,U_1(i),\,U_2(i),\,U_1(i+1),\,U_2(i+1)\bigr),
\]
and we write $\theta(i)=(N_0,N_1,N_2,u_1,u_2,v_1,v_2)$. By Lemma~\ref{lem:pattern},
$\theta(i)$ determines $U(i)=(-N_1,u_1,u_2,N_0)$ and $U(i+1)=(-N_2,v_1,v_2,N_1)$ completely,
and the $k$-relations at $i$ and $i+1$ read
\[
N_1k_0=u_1k_1+u_2k_2+N_0k_3,\qquad N_1k_4=N_2k_1-v_1k_2-v_2k_3 .
\]

\subsection{Two identities}

\begin{lemma}[The normalization identity]\label{lem:multid}
Let $\pi_i$ be the monic orthogonal polynomial of degree four for $\nu_i$, as in the remark
above, and let $H_4=\det\bigl(\mu_i(x^{j+l})\bigr)_{j,l<4}$. Then $H_4$ is a polynomial in
$k_0,\dots,k_4$ with $43$ terms, independent of $i$, and
\[
I_4(i)\,\pi_i(-q_i)=a\,H_4(k_0,\dots,k_4).
\]
As a polynomial in $q$, $I_4(i)\pi_i(-q)$ has degree four and leading coefficient
$I_4(i)=L(i)\ge1$.
\end{lemma}

\begin{proof}
Since $F_{m+1}-F_m=F_m(x)(x-s_m-1)$, the remark gives
\[
I_4(i)\pi_i(x)=\sum_{m=0}^3U_m(i)\,F_m(x)\,(x-s_m-1).
\]
The division equations read
$b_{i+m+1}=a-(q_i+s_m)b_{i+m}$. Use them to eliminate $b_{i+1},\dots,b_{i+4}$ from both sides.
The difference becomes the zero polynomial in $a,q_i,b_i,k_0,\dots,k_4$; explicit certificates
are in the Lean proof. Conceptually, the identity follows from the determinant formula for
$\pi_i$, since $\nu_i(x^{s+1})+q_i\nu_i(x^s)=\mu_i(x^s)$. The leading coefficient is that of
the monic $\pi_i$.
\end{proof}

\begin{lemma}[Labels]\label{lem:labels}
Let $i+6<N$ and put
\[
G(k_0,\dots,k_4)=(k_0+1)\bigl(k_1s_2+s_2+k_1+k_2+1\bigr)
\bigl(k_2(k_1+k_2)s_3+e_3\bigr),
\]
where $e_3=s_3(k_1+k_2)+s_3k_2+(k_1+k_2)k_2+(k_0+2k_1+3k_2)(k_3+1)+k_3^2+2k_3+k_4+1$. Then
$L(i)\,q_i^4\le a\,G(k_0,\dots,k_4)$. If moreover $b_i<2A$, then
$L(i)\,a^3\le 256A^4\,G(k_0,\dots,k_4)$.
\end{lemma}

\begin{proof}
The Newton polynomials $F_0,\dots,F_3$ are monic of degrees $0,\dots,3$. So the matrix
$\bigl(\nu_i(F_jF_l)\bigr)$ is $\mathcal U\mathcal H\mathcal U^{\mathsf T}$, where $\mathcal H$ is the Hankel matrix of $\nu_i$ and $\mathcal U$ is unit
lower triangular. Hence it has determinant $I_4(i)/a$, and it is positive definite, as in the
proof of Theorem~\ref{thm:hankel}. By Hadamard's inequality,
\[
\frac{I_4(i)}a\le\prod_{j<4}\nu_i(F_j^2)\le\prod_{j<4}\frac{\mu_i(F_j^2)}{q_i},
\]
since $\nu_i=\mu_i/(q_i+x)\le\mu_i/q_i$ on the nonnegative support. By the shift
rule~\eqref{eq:moment-shift} and the moments $\mu_{i+m}(1)=\mu_{i+m}(x)=1$,
$\mu_{i+m}(x^2)=k_m+1$, $\mu_{i+m}(x^3)=k_m^2+2k_m+k_{m+1}+1$, the product equals $G$. For the second
claim, $a=q_ib_i+b_{i+1}<2A(q_i+1)\le4Aq_i$, so $L(i)a^4\le256A^4L(i)q_i^4\le 256A^4aG$.
\end{proof}

\subsection{Signs}
Throughout this subsection the pattern and gaps arise from the same state,
so the cofactor relations of Lemma~\ref{lem:pattern} remain in force.
Besides the polynomial $\sigma$ of Section~\ref{sec:six}, put
$\omega(x,y,z,w)=x^2z+(2x-1)yz+xz^2+xw$. Both $\sigma$ and $\omega$ are positive at
positive integers.

\begin{lemma}[Signs]\label{lem:signs}
If $N_0,N_1\ge1$ and $k_0,\dots,k_4\ge1$, then $u_1>0>u_2$.
\end{lemma}

\begin{proof}
Solving the $k$- and $r$-relations of Lemma~\ref{lem:pattern} for $(u_1,u_2)$ by Cramer's rule
gives the exact identities
\begin{align*}
u_1\sigma(k_1,k_2,k_3)&=N_1\omega(k_0,k_1,k_2,k_3)+N_0\sigma(k_2,k_3,k_4),\\
-u_2\sigma(k_1,k_2,k_3)&=N_1\sigma(k_0,k_1,k_2)+N_0\omega(k_1,k_2,k_3,k_4),
\end{align*}
since the Cramer determinant is $k_1r_2-k_2r_1=\sigma(k_1,k_2,k_3)$.
\end{proof}


\subsection{Good starts}

\begin{lemma}[Good starts]\label{lem:good}
Let $b_s,\dots,b_{s+T-1}$ be a window run in $[A,2A)$ with $T\ge36$, and put
$H=\lfloor 28a/(AT)\rfloor+1$. Call a start $i$ \emph{good} if $s\le i$, $i+9\le s+T$ and
$q_{j+1}-q_j\le H$ for $i\le j\le i+6$, that is, $k_0,\dots,k_6\le H$ at state $i$. Then at least
$T/2$ starts are good, and every good start satisfies $i+8<N$ and $b_i,b_{i+1},b_{i+2}<2A$.
\end{lemma}

\begin{proof}
Let $\mathcal B$ be the set of $j\in[s,s+T-1)$ with $q_{j+1}-q_j>H$. Then
$|\mathcal B|(H+1)\le q_{s+T-1}-q_s\le a/A$. Since $H\cdot AT>28a$, this gives
$28|\mathcal B|\le T$. Each $j\in\mathcal B$ spoils at most seven starts, so at least
$T-9-7|\mathcal B|\ge T/2$ starts are good, as $T\ge36$. The last assertion is immediate,
since the states $i,\dots,i+8$ lie in the window.
\end{proof}

\subsection{The full minor vector}\label{sec:minors}
The seven pattern integers are not independent: they are coordinates of a single vector of
$3\times3$ minors. If all ten minors are at most $X$ in absolute value, there are only
$X^{5+o(1)}$ such vectors (Proposition~\ref{prop:minorvec}), rather than the $X^7$ choices of seven
independent integers.
We therefore define the \emph{pattern} of a state $i$ to be the vector $p(i)$ of the ten
$3\times3$ minors of the matrix $\widehat{\mathcal M}_i$ below; it determines the pattern
integers $\theta(i)$.

Let $\widehat{\mathcal M}_i$ be the $3\times5$ matrix with rows $(d_0,\dots,d_4)$, $(k_0,\dots,k_4)$
and $(r_0,\dots,r_4)$, and let $p_{c_1c_2c_3}$ be its minor on columns $c_1<c_2<c_3$. Expanding
the cofactors above gives
\[
\theta(i)=(p_{012},\,p_{123},\,p_{234},\,p_{023},\,-p_{013},\,p_{134},\,-p_{124}).
\]

\begin{lemma}[Expansion]\label{lem:expand}
Let $q=q_i$ and $s_j=q_{i+j}-q$. If $i+j+3<N$ then
\[
\frac{b_{i+j}}a=\frac1q-\frac{1+s_j}{q^2}+\frac{s_j^2+2s_j+k_j+1}{q^3}-e_j,\qquad
0\le e_j\le\frac{\mu_{i+j}\bigl((x+s_j)^3\bigr)}{q^4}.
\]
\end{lemma}

\begin{proof}
Apply $\mu_{i+j}$ to $\frac1{q+y}=\frac1q-\frac y{q^2}+\frac{y^2}{q^3}-\frac{y^3}{q^3(q+y)}$ with
$y=s_j+x\ge0$, using $b_{i+j}/a=\nu_{i+j}(1)=\mu_{i+j}\bigl(1/(q+s_j+x)\bigr)$, $\mu(1)=\mu(x)=1$
and $\mu(x^2)=k+1$.
\end{proof}

\begin{proposition}[All minors]\label{prop:allminors}
Let $i+8<N$, $b_i<2A$ and $k_0,\dots,k_6\le H$ at state $i$. Then every minor of $\widehat{\mathcal M}_i$
satisfies $|p|\,a^3\le3609600\,A^4H^6$.
\end{proposition}

\begin{proof}
Subtracting the expansions of Lemma~\ref{lem:expand} for $j$ and $j+1$ gives
$d_j/a=k_j/q^2-(r_j+k_j)/q^3+e_{j+1}-e_j$. So replacing the first row by
$d'_j=d_j-a(q^{-2}-q^{-3})k_j+aq^{-3}r_j=a(e_{j+1}-e_j)$ changes no minor. Now
$|d'_j|\le235aH^3/q^4$, $|k_j|\le H$ and $|r_j|\le10H^2$, so $|p|\le14100\,aH^6/q^4$. Finally
$a\le4Aq$.
\end{proof}

\subsection{Counting minor vectors}

\begin{proposition}[Minor vectors]\label{prop:minorvec}
The number of vectors of minors of integer $3\times5$ matrices with all coordinates in $[-X,X]$ is
at most $1+20\sum_{\kappa\le X}729\,\kappa^4\tau(\kappa)\ll X^{5+o(1)}$.
\end{proposition}

\begin{proof}
For a nonzero vector $p$ choose a pivot $\pi$ with $\kappa=|p_\pi|$ maximal, and let $c_4,c_5$ be
the two columns outside $\pi$. Let $u_r$ ($v_r$) be the minor with the $r$-th pivot column
replaced by $c_4$ ($c_5$), with the sign of the column permutation. By Cramer's rule and the
Pl\"ucker relations, the remaining three minors satisfy $p_\pi\,p_{w_r}=\pm(u\times v)_r$
(thirty polynomial identities, checked symbolically). So $p$ is determined by
$(\pi,\operatorname{sgn}p_\pi,\kappa,u,v)$, where $|u_r|,|v_r|\le\kappa$ and $\kappa$ divides
every coordinate of $u\times v$.

Fix $\kappa\ge1$ and $u$, and let $g=\gcd(\kappa,u_1,u_2,u_3)$, $\ell=\kappa/g$, $u=gu_0$. By
B\'ezout there are $z\in\mathbb Z^3$ and $t\in\mathbb Z$ with $z\cdot u_0+t\ell=1$. If
$\kappa\mid u\times v$, then $\ell\mid u_0\times v$, and the identity
$v-(z\cdot v)u_0=t\ell\,v+z\times(v\times u_0)$ shows $v=\mu u_0+\ell y$ with $\mu=z\cdot v$ and
$y\in\mathbb Z^3$. Hence $v$ lies in one of $\ell$ classes modulo $\ell$, each meeting
$[-\kappa,\kappa]^3$ in at most $(3g)^3$ points: at most $27\kappa g^2$ choices. Summing over
$u$, grouped by $g=e\mid\kappa$, gives
$\sum_u g^2\le\sum_{e\mid\kappa}e^2(3\kappa/e)^3\le27\kappa^3\tau(\kappa)$, so at most
$729\kappa^4\tau(\kappa)$ pairs $(u,v)$. There are ten pivots and two signs.
\end{proof}

\section{Joint integer invariants and the optimized bound}\label{sec:joint}

We now prove Theorem~\ref{thm:joint}. Fix a window run in $[A,2A)$ with $T\ge36$
terms. Throughout this section constants may depend on fixed positive integers
$m,n$ (the powers of Section~\ref{sec:jointcount}) and on a loss $\eta>0$, but
never on the dividend, the orbit or the pattern. The argument has six steps.
\begin{enumerate}
\item For a fixed pattern, the nodes of every good start lie on at most
$24$ curves (Section~\ref{sec:jointcurve}).
\item On these curves the quotient $q_i$ and the curvature $I_3(i)$ are
rational functions of the nodes. Components on which the reconstruction
identities fail contribute $O(1)$ starts (Section~\ref{sec:jointrecon}).
\item The directions at infinity of the curves are controlled by the
Newton coefficients of the pattern. This makes nonexceptional
components have even degree (Section~\ref{sec:jointinfinity}), shows
that exceptional patterns are few, and excludes lines
(Section~\ref{sec:jointexc}).
\item The reconstructed functions, and hence two families of new integer
coordinates $W$ and $J$, have prescribed orders at infinity
(Sections~\ref{sec:jointorders} and~\ref{sec:jointcoords}).
\item Each nonexceptional conic contributes $a^{o(1)}$ starts
(Section~\ref{sec:jointconic}).
\item On the remaining components, adjoining $W_j^m$ and $J^n$ raises the
degree. The curve-counting theorem and a balancing of heights then give the
exponent (Sections~\ref{sec:jointgraph} and~\ref{sec:jointcount}).
\end{enumerate}

\emph{Notation.} In this section
\begin{itemize}
\item $\theta=(N_0,N_1,N_2,u_1,u_2,v_1,v_2)$ are the pattern integers of
Section~\ref{sec:patterns}, with $N_0=I_4(i)$, and bars denote their normalized
versions~\eqref{eq:jointbar};
\item $s_j$ are the nodes, $k_j$ the gaps, $P_4,R_4$ the two orthogonal
quartics, $\Psi$ the quartic norm and $f$ the pattern quadratic;
\item $D_{jl}=d_jk_l-d_lk_j$ are the $2\times2$ minors of drops and gaps, so that
$D_{01}=I_3(i)$ is the curvature of Section~\ref{sec:curv};
\item $Q$ and $E_0$ are the reconstructed quotient and curvature,
$W=(W_0,W_1,W_2)$ and $J$ the joint coordinates, and $S$, $Y$ their height
bounds;
\item $e$ is a multiplicity at infinity, $d$ the degree of a component,
$\delta$ the degree of a graph and $\delta_{\min}=10n+7m$.
\end{itemize}

\subsection{Good starts and the pattern curve}\label{sec:jointcurve}

Put
\[
 H=\left\lfloor\frac{28a}{AT}\right\rfloor+1,\qquad
 X=\left\lfloor\frac{3609600 A^4H^6}{a^3}\right\rfloor.
\]
Let the good starts be as in Lemma~\ref{lem:good}; at least $T/2$ starts are good. By
Proposition~\ref{prop:allminors}, the minor vector of a good start has all
coordinates at most $X$, and $X\ge1$ because its coordinate $p_{012}=L(i)$ is a
positive integer. The pattern of a good start (Section~\ref{sec:minors}) therefore has
height at most $X$, and it determines $\theta(i)$. By Proposition~\ref{prop:minorvec}, there are
$O_\eta(a^\eta X^5)$ patterns.

The \emph{nodes} of state $i$ are $s_j=q_{i+j}-q_i$. They are the support
points of the measure $\mu_i$ of Section~\ref{sec:hankel}, and the
interpolation nodes of the Newton polynomials of Section~\ref{sec:patterns}. The
three \emph{free nodes} $s_1,s_2,s_3$ of a good start are integers in $[0,3H]$.
Distinct good starts with the same pattern have distinct free-node triples
(Lemma~\ref{lem:jointcuts} below).

For the pattern integers $\theta$, put
\begin{equation}\label{eq:jointbar}
 (\bar N_1,\bar u_1,\bar u_2)=\frac{(N_1,u_1,u_2)}{N_0},\qquad
 (\bar N_2,\bar v_1,\bar v_2)=\frac{(N_2,v_1,v_2)}{N_1}.
\end{equation}
By Lemma~\ref{lem:signs}, $\bar N_1>0$, $\bar u_1>0$ and $\bar u_2<0$.
Treat the free nodes $s_1,s_2,s_3$ as coordinates and define
\[
 s_0=0,\quad
 s_4=(\bar N_1+\bar u_1)s_1-(\bar u_1-\bar u_2)s_2-(\bar u_2-1)s_3,\quad
 k_j=s_{j+1}-s_j\quad(0\le j\le3),
\]
\[
 k_4=\bar N_2k_1-\bar v_1k_2-\bar v_2k_3,\qquad s_5=s_4+k_4.
\]
At a good start these are the actual nodes and gaps, by the two
$k$-relations of Lemma~\ref{lem:pattern}; the first is equivalent to
$\bar N_1k_0=\bar u_1k_1+\bar u_2k_2+k_3$. Let
\[
 \gamma=(\bar N_1,\,-\bar N_1-\bar u_1,\,\bar u_1-\bar u_2,\,\bar u_2-1,\,1),
\]
and let $\gamma'$ be the same expression in $\bar N_2,\bar v_1,\bar v_2$. Put
\[
 P_4(x)=\sum_{j=0}^4\gamma_j\prod_{\ell=0}^{j-1}(x-s_\ell),
 \qquad
 R_4(x)=\sum_{j=0}^4\gamma'_j\prod_{\ell=1}^{j}(x-s_\ell).
\]
By the remark after Lemma~\ref{lem:pattern}, at a good start $P_4$ is the
monic orthogonal polynomial $\pi_i$ of degree four for $\nu_i$, and
$R_4(x)=\pi_{i+1}(x-k_0)$ is the corresponding polynomial of the next state,
written in the variable of the current state. Put $Z=xR_4$ and
$f=\sum_{j=0}^4\gamma_j(s_j^2+k_j)$, the \emph{pattern quadratic}, and let
\begin{align*}
 \Psi={}&\sum_{j=0}^4\gamma_j
 \bigl(s_j^3+3s_j^2+3s_j(1+k_j)+(1+k_j)^2\bigr)
 +\sum_{j=0}^3\gamma_jk_{j+1}\\
 &+\bar N_2r_1-\bar v_1r_2-\bar v_2r_3-k_4(2s_4+k_4),
\end{align*}
where $r_j=k_j(s_j+s_{j+1})+k_{j+1}$ for $j=1,2,3$. This is a cubic
polynomial in the free nodes, the \emph{quartic norm}.

Degrees throughout this section mean the leading coefficient of the eventual
affine Hilbert function, equivalently the degree of the projective closure.
``Integral'' means that the defining ideal is prime over $\mathbb Q$; it does
not assert geometric irreducibility or a rational parametrization.

\begin{lemma}[The pattern curve]\label{lem:jointcurve}
\begin{enumerate}
\item At a good start $i$, $\Psi=\nu_i(P_4^2)>0$, and the free nodes satisfy
\begin{equation}\label{eq:jointcurve}
 f=0,\qquad N_0\operatorname{Res}(P_4,Z)=a\Psi^4.
\end{equation}
\item For every pattern with $\bar N_1,\bar u_1>0>\bar u_2$, the two
equations~\eqref{eq:jointcurve} have degrees two and at most twelve and no
common factor. Their common zero set in $\mathbb A^3$ is the union of at most
twenty-four integral affine curves over $\mathbb Q$, of total degree at most
twenty-four.
\end{enumerate}
\end{lemma}

\begin{proof}
(1) Since $r_m=s_{m+1}^2-s_m^2+k_{m+1}$, summation by parts and the
$k$-relation at state $i$ show that $N_0f$ is the $r$-relation
$\sum_mU_m(i)r_m=0$ of Lemma~\ref{lem:pattern}. Next, the shift rule
$\mu_i\bigl(g\prod_{\ell<j}(x-s_\ell)\bigr)=\mu_{i+j}\bigl(g(x+s_j)\bigr)$ and
the moments listed in the proof of Lemma~\ref{lem:labels} give
$\mu_i(x^3P_4)=\sum_j\gamma_j\bigl(s_j^3+3s_j^2+3s_j(1+k_j)+(1+k_j)^2+k_{j+1}\bigr)$.
Eliminating $k_5$ with the $r$-relation at state $i+1$ turns this into $\Psi$.
Since $\nu_i(xg)=\mu_i(g)-q_i\nu_i(g)$ and $P_4$ is orthogonal to $x^3$,
$\nu_i(P_4^2)=\nu_i(x^4P_4)=\mu_i(x^3P_4)=\Psi$, and this is positive as in
the proof of Theorem~\ref{thm:hankel}.

For the resultant, Lemma~\ref{lem:moments}(3) gives
$\nu_i\bigl(xg(x)\bigr)=\nu_{i+1}\bigl(g(x+k_0)\bigr)$. Hence, for $j<4$,
$\nu_i(x^jZ)=\nu_{i+1}\bigl((x+k_0)^j\pi_{i+1}(x)\bigr)=0$. So the monic
quintic $Z$ differs from the monic orthogonal polynomial $\pi^{(5)}$ of degree
five for $\nu_i$ by a multiple of $P_4$, and $\operatorname{Res}(P_4,Z)=
\operatorname{Res}(P_4,\pi^{(5)})$. Let $n_j$ be the squared $\nu_i$-norm of
the monic orthogonal polynomial of degree $j$. The three-term recurrence,
whose coefficients are $\beta_j=n_j/n_{j-1}$, gives
\[
 \operatorname{Res}(P_4,\pi^{(5)})=\beta_1\beta_2^2\beta_3^3\beta_4^4
 =\frac{n_4^4}{n_0n_1n_2n_3}
 =\frac{\Psi^4}{\det\bigl(\nu_i(x^{r+s})\bigr)_{r,s<4}}
 =\frac{a\Psi^4}{I_4(i)}.
\]
Since $N_0=I_4(i)$, this is the second equation.

(2) The degree bounds are immediate, since the resultant has total degree at
most ten in the nodes. Note that $\gamma_1<0<\gamma_2$, $\gamma_3=\bar u_2-1<-1$ and
$\gamma_0,\gamma_4>0$. Since $s_0=0$ and $\sum_j\gamma_js_j=0$, the quadratic part of $f$ is
\[
 f_2=\gamma_1s_1^2+\gamma_2s_2^2+\gamma_3s_3^2+(\gamma_1s_1+\gamma_2s_2+\gamma_3s_3)^2,
\]
whose Hessian determinant is $8\gamma_1\gamma_2\gamma_3(1+\gamma_1+\gamma_2+\gamma_3)
=-8\gamma_0\gamma_1\gamma_2\gamma_3\ne0$. A ternary quadratic form of rank three is not a
product of linear forms, so $f_2$, and hence $f$, is irreducible. The degree-twelve part of
the second equation is $-ac^4$ with $c=\sum_j\gamma_js_j^3$. If $f$ divided the second
equation, the irreducible form $f_2$ would divide $c^4$, hence $c$. Divide $c$ by $f_2$ as
polynomials in $s_3$; the $s_3^2$ coefficient of $f_2$ is $\gamma_3(1+\gamma_3)\ne0$. The
coefficients of the remainder vanish only if
$\gamma_1\gamma_2\gamma_3^2(\gamma_3^2-1)=0$, which the signs above exclude. So the two
equations have no common factor, their common zero set is a curve, and by B\'ezout's theorem
its components over $\mathbb Q$ have total degree at most $2\cdot12$.
\end{proof}

\subsection{Reconstruction and retained components}\label{sec:jointrecon}

Two further identities distinguish the components on which the later rational
functions have the required behavior. We give their definitions before
stating the component bounds.

Form the $8\times9$ matrix of coefficients of
$P_4,xP_4,x^2P_4,x^3P_4,Z,xZ,x^2Z,x^3Z$ in the basis
$1,x,\dots,x^8$. Let $c_j$ be $(-1)^j$ times the minor obtained by
deleting column $j$, and define
\[
 Q=\frac{c_1-c_2}{c_1-c_0},\qquad
 E_0=a\,\frac{\det(c_{r+s})_{0\le r,s<3}(c_1-c_0)^3}
                  {(c_1^2-c_0c_2)^3}.
\]
At a good start the matrix annihilates the moment vector
$(h_i(0),\dots,h_i(8))$ and has rank eight. Indeed, the resultant in
\eqref{eq:jointcurve} is positive, so $P_4$ and $Z$ are coprime.
The identities
$q_ih_i(0)+h_i(1)=q_ih_i(1)+h_i(2)=1$, together with
$h_i(0)h_i(2)>h_i(1)^2$, then give
\[
 h_i(j)=\frac{c_j(c_1-c_0)}{c_1^2-c_0c_2}\quad(0\le j\le8),\qquad
 Q=q_i,\qquad E_0=I_3(i).
\]
In particular, all displayed denominators are nonzero at good starts.

To state the identities as equations on the curve, use the raw-moment
polynomials
\[
 M_0=1,\qquad
 M_{r+1}(k_0,k_1,\dots)
 =\sum_{t=0}^{r}\binom rt k_0^{r-t}M_t(k_1,k_2,\dots).
\]
Only $M_0,\dots,M_6$ are needed; they involve $k_0,\dots,k_4$.
At good starts, $M_r=\mu_i(x^r)$ for $0\le r\le6$, by the
weight recursion in Section~\ref{sec:hankel}. Set
\[
 H_3=\det(M_{r+s})_{0\le r,s<3},\qquad
 H_4=\det(M_{r+s})_{0\le r,s<4},\qquad
 H_{013}=\det(M_{r+s})_{\substack{r=0,1,2\\s=0,1,3}}.
\]
This $H_4$ is the polynomial of Lemma~\ref{lem:multid}. With
$p_3=\bar u_2-1-s_1-s_2-s_3$, the coefficient of $x^3$ in $P_4$, consider the
\emph{normalization} and \emph{resolvent identities}
\begin{equation}\label{eq:jointcuts}
 N_0P_4(-Q)=aH_4,\qquad
 E_0P_4(-Q)=a\bigl((Q-p_3)H_3-H_{013}\bigr).
\end{equation}
Both hold at good starts, where $Q=q_i$, $E_0=I_3(i)$ and $P_4=\pi_i$. The first is
Lemma~\ref{lem:multid}. For the second, write $v_j=h_i(j)$ and
$P_4(x)=x^4+p_3x^3+p_2x^2+p_1x+p_0$. Orthogonality of $P_4$ to $1,x,x^2$ gives
$p_0v_j+p_1v_{j+1}+p_2v_{j+2}+p_3v_{j+3}+v_{j+4}=0$ for $j<3$, and
$q_iv_j+v_{j+1}=\mu_i(x^j)=M_j$. Eliminating $v_4,v_5,v_6$ with the three relations, a direct
expansion gives the polynomial identity
\[
 \det(v_{r+s})_{r,s<3}\,P_4(-q_i)=(q_i-p_3)\det(M_{r+s})_{r,s<3}
 -\det(M_{r+s})_{\substack{r=0,1,2\\s=0,1,3}},
\]
and $a\det(v_{r+s})_{r,s<3}=I_3(i)$.
Among components containing a good start, we call a component
\emph{retained} if both equations, after clearing denominators, hold
identically on it.

\begin{lemma}[Removing components by polynomial equations]\label{lem:jointcuts}
The two identities in \eqref{eq:jointcuts}, after clearing denominators,
have degrees at most $166$ and $404$, respectively. On a degree-$d$
component, either an identity holds everywhere or it holds at
$166d$ or $404d$ good starts at most, respectively. Consequently each
component that is not retained contributes $O(1)$ good starts.
\end{lemma}
\begin{proof}
Write $Q=Q_{\mathrm n}/Q_{\mathrm d}$ and $E_0=E_{\mathrm n}/E_{\mathrm d}$ using the
cofactor expressions above. Each entry of the coefficient matrix has degree at most four in
the nodes, so each $8\times8$ cofactor has degree at most $32$, and in particular at most
$40$, the bound used in the formal proof. Hence
\[
 \deg Q_{\mathrm n},\deg Q_{\mathrm d}\le40,\qquad
 \deg E_{\mathrm n},\deg E_{\mathrm d}\le240.
\]
The latter bound follows from the $3\times3$ determinant and the cube
in the numerator, and the cube of a quadratic cofactor expression in the
denominator. The polynomial
$\widehat P=Q_{\mathrm d}^4P_4(-Q_{\mathrm n}/Q_{\mathrm d})$ has degree at most $163$: in Newton
form its longest product is
$Q_{\mathrm n}\prod_{j=1}^3(Q_{\mathrm n}+s_jQ_{\mathrm d})$, of degree at most $40+3\cdot41$,
because $s_0=0$. The cleared equations are
\begin{align*}
 \Phi_N&=N_0\widehat P-aH_4Q_{\mathrm d}^4,\\
 \Phi_R&=E_{\mathrm n}\widehat P
 -a\bigl((Q_{\mathrm n}-p_3Q_{\mathrm d})H_3-H_{013}Q_{\mathrm d}\bigr)Q_{\mathrm d}^3E_{\mathrm d}.
\end{align*}
Here $\deg H_3\le3$, $\deg H_4\le6$ and $\deg H_{013}\le4$.
Thus $\deg\Phi_N\le166$ and $\deg\Phi_R\le404$.

If a polynomial of degree $r$ does not vanish identically on an integral
curve of degree $d$, its zero set on that curve has at most $rd$ distinct
rational points, by the polynomial-section form of B\'ezout's bound.
Every good start satisfies both cleared equations. Moreover, the
free nodes and pattern recover $q_i=Q$, and $q_i$ strictly increases along
the fixed orbit. Thus distinct good starts of this pattern give distinct free-node triples.
Applying the point bound to whichever equation is not an identity proves
the claim. The absolute bound $d\le24$ makes the last estimate uniform.
\end{proof}

\subsection{Directions at infinity}\label{sec:jointinfinity}

Call a pattern \emph{exceptional} if two of the five Newton coefficients
$\gamma_j$ sum to zero. Since $\bar N_1>0$, $\bar u_1>0$ and $\bar u_2<0$,
every $\gamma_j$ is nonzero, and a pair sum can vanish only in one of the
four cases
\begin{equation}\label{eq:jointexceptional}
 \bar N_1+\bar u_2=1,\qquad \bar N_1+\bar u_2=0,\qquad \bar u_1=1,\qquad \bar N_1+\bar u_1=1.
\end{equation}
Equivalently, the integer pattern coordinates satisfy
$N_1+u_2=N_0$, $N_1+u_2=0$, $u_1=N_0$ or $N_1+u_1=N_0$.

\begin{lemma}[Leading nodes and determinant orders]\label{lem:jointleading}
Let $C$ be an integral component of positive degree of
\eqref{eq:jointcurve}. Let $P$ be a closed point above infinity on the
normalization of the projective closure of $C$, write
$\operatorname{ord}=\operatorname{ord}_P$, and let $e>0$ be the
multiplicity of the hyperplane at infinity at $P$.
For a local parameter $\varpi$ at $P$, let $z_j$ be the residue of
$\varpi^es_j$, $0\le j\le4$.
Then the $z_j$ are not all equal and satisfy
\[
 \sum_{j=0}^4\gamma_jz_j^r=0\qquad(0\le r\le3).
\]
If two of these five leading nodes coincide, the pattern is exceptional.
For a nonexceptional pattern, therefore,
\begin{align*}
 \operatorname{ord}(s_l-s_j)&=-e\quad(0\le j<l\le4),\\
 \operatorname{ord}\operatorname{Res}(P_4,Z)&=-10e,\\
 \operatorname{ord}H_3&=-3e,\qquad \operatorname{ord}H_4=-6e.
\end{align*}
Positive order denotes a zero and negative order a pole.
\end{lemma}

\begin{proof}
All nonzero rational constants have order zero. The multiplicity $e$ can
equivalently be written as
\[
 e=\max_{1\le j\le3}\{-\operatorname{ord}s_j\}>0.
\]
Thus a polynomial of total node degree at most $r$ has order at least
$-re$. The residues $z_j$ lie in the residue field of $P$; we do not
require $P$ to be rational.
We have $z_0=0$, and at least one of $z_1,z_2,z_3$ is nonzero.

\emph{The leading nodes are distinct.}
The definitions of the nodes give the moment equations of orders zero
and one. The highest-degree part of $f=0$ gives the equation of order two.
Finally, the highest-degree part of
$N_0\operatorname{Res}(P_4,Z)=a\Psi^4$ gives the equation of order three:
the resultant has total degree at most ten, whereas the cubic part of
$\Psi$ is $\sum_j\gamma_js_j^3$. Consequently
\begin{equation}\label{eq:joint-leading-moments}
 \sum_{j=0}^4\gamma_jz_j^r=0\qquad(0\le r\le3).
\end{equation}
For the last of these, multiply the fixed-dividend equation by
$\varpi^{12e}$ and take residues; its left side vanishes and its right side
is $a(\sum_j\gamma_jz_j^3)^4$.

If two nodes coincided, there would be at most four distinct values among
the $z_j$. Grouping the sums in \eqref{eq:joint-leading-moments} by those
values and using the corresponding Vandermonde matrix shows that the
sum of the $\gamma_j$ in each group is zero.
Since there are five indices and at least two groups, some group has at
most two elements. It cannot have one, because every $\gamma_j$ is nonzero;
it must therefore give a vanishing pair sum. This proves the collision
assertion. For a nonexceptional pattern it is impossible, so all five
nodes are distinct, and
\begin{equation}\label{eq:joint-node-orders}
 \operatorname{ord}(s_l-s_j)=-e\qquad(0\le j<l\le4).
\end{equation}

\emph{The determinant orders.}
Write $\mathcal R=\operatorname{Res}(P_4,Z)$.
The leading homogeneous parts needed below are
\begin{align}
 [\mathcal R]_{10}
   &=\gamma_0\gamma_1\gamma_2\gamma_3
        \prod_{0\le j<l\le4}(s_l-s_j), \label{eq:joint-resultant-leading}\\
 [H_3]_3&=k_0k_1(k_0+k_1),\notag\\
 [H_4]_6&=k_0k_1k_2(k_0+k_1)(k_1+k_2)(k_0+k_1+k_2).
 \label{eq:joint-hankel-leading}
\end{align}
Here $[F]_r$ denotes the homogeneous part of total node degree $r$.
For completeness, the degree-ten formula can be checked without expanding
the resultant. In the Newton basis beginning at $s_1,s_2,s_3$, its
$4\times4$ Bezout matrix for $P_4,R_4$ has entry degrees at most
$4-\max(i,j)$, for $0\le i,j<4$. Only the diagonal term of its determinant
can have degree ten. Its four leading diagonal entries are
\[
\begin{split}
 &-\gamma_1s_1(s_1-s_2)(s_1-s_3)(s_1-s_4),\\
 &-\gamma_2s_2(s_2-s_3)(s_2-s_4),\\
 &-\gamma_3s_3(s_3-s_4),\qquad -s_4.
\end{split}
\]
Their product gives the stated Vandermonde; the additional factor
$x$ in the second polynomial contributes
$P_4(0)=\gamma_0$.
The raw-moment recurrence gives
\[
 H_3=k_0k_1(k_0+k_1)+k_0k_2-k_1^2,
\]
and gives the displayed degree-six part of $H_4$, with remainder of
degree at most five. These are finite polynomial identities.

Every factor in these leading products has order $-e$ by
\eqref{eq:joint-node-orders}. The lower-degree terms have strictly larger
order, so they cannot cancel the leading terms. Thus
\begin{equation}\label{eq:joint-determinant-orders}
 \operatorname{ord}\mathcal R=-10e,\qquad
 \operatorname{ord}H_3=-3e,\qquad
 \operatorname{ord}H_4=-6e.
\end{equation}
\end{proof}

\begin{corollary}[Even component degrees]\label{cor:jointparity}
Every positive-degree nonexceptional component of \eqref{eq:jointcurve}
has even degree. At each point above infinity its hyperplane multiplicity
$e$ is even, and
$\operatorname{ord}(N_0\Psi)=-5e/2$.
\end{corollary}
\begin{proof}
Taking orders in the fixed-dividend equation and using
Lemma~\ref{lem:jointleading} gives
$4\operatorname{ord}\Psi=-10e$.
Since orders are integers and $N_0\ne0$, this proves both local assertions.
The degree of the curve is the sum of the multiplicities $e$, weighted
by the residue-field degrees, so it is even. The retained identities
are not needed here.
\end{proof}

\subsection{Exceptional patterns and line components}\label{sec:jointexc}

Exceptional patterns may have components of odd degree, including lines, but
there are fewer of them.

\begin{lemma}[Exceptional patterns and lines]\label{lem:jointcomponents}
There are $O(X^4\log(2X))=O_\eta(a^\eta X^4)$ exceptional patterns. In the range
\begin{equation}\label{eq:jointscale}
 124928 A^4H^5X<a^3,
\end{equation}
no line component contains a good start.
\end{lemma}
\begin{proof}
Each relation in \eqref{eq:jointexceptional} becomes a vanishing maximal
minor after a fixed invertible integer column change. The transformed
minor vector still has height $O(X)$. For example, if the original columns
are $C_0,\dots,C_4$, then
\[
 \det(C_1,C_2+C_0,C_3-C_0)=p_{123}-p_{013}-p_{012}
                              =N_1+u_2-N_0.
\]
The other three relations follow by the same column additions and
permutations.

A fixed zero minor reduces the Pl\"ucker relations to
three integer vectors $\mathbf x,\mathbf y,\mathbf z$ of height $O(X)$
with $\mathbf x\times\mathbf y=0$ and
$\mathbf x\cdot\mathbf z=\mathbf y\cdot\mathbf z=0$.
These vectors are signed rearrangements of the nine remaining minors,
so they determine the minor vector.
The case $\mathbf x=\mathbf y=0$ has $O(X^3)$ choices. Otherwise write
$\mathbf x=r\mathbf v$, $\mathbf y=s\mathbf v$ with $\mathbf v$ primitive.
In a dyadic range $R\le|\mathbf v|_\infty<2R$, there are $O(R^3)$
choices for $\mathbf v$ and $O((X/R)^2)$ for $(r,s)$.
The equation $\mathbf v\cdot\mathbf z=0$ has $O(X^2/R)$ solutions
in the box. Indeed, take a coordinate of absolute value
$n=|\mathbf v|_\infty$. Primitivity gives $n$ admissible pairs of the
other two coordinates modulo $n$, and the chosen coordinate is then
determined. Since $R\le n\ll X$, this gives $O(X^2/n)$ choices.
Each range therefore contributes $O(X^4)$, and summing over $O(\log(2X))$
ranges proves the pattern count. The crude bound $X\le10^{16}a^7$
absorbs its logarithm into $a^\eta$.

Now suppose that a line contained in the zero set of~\eqref{eq:jointcurve} passes
through the free-node triple $s$ of a good start, with direction $v$. Extend $v$ to the five
nodes linearly, as in Section~\ref{sec:jointcurve}, and write $s(t)=s+tv$. Both equations
of~\eqref{eq:jointcurve} hold identically in $t$.
\begin{enumerate}
\item The resultant has degree at most ten in $t$, so the $t^{12}$ coefficient of the second
equation, $-a\bigl(\sum_j\gamma_jv_j^3\bigr)^4$, vanishes. Then $\Psi(s(t))$ has degree at most
two, so $\Psi^4$ has degree at most eight, and the $t^{10}$ coefficient
$[\mathcal R]_{10}(v)$ of the resultant, given by~\eqref{eq:joint-resultant-leading}, vanishes
too: two coordinates of $v$ coincide.
\item The direction satisfies the four moment equations of Lemma~\ref{lem:jointleading}: those
of orders zero and one by the definition of the nodes, order two from the $t^2$ coefficient of
$f$, and order three from step~(1). The grouping argument in the proof of
Lemma~\ref{lem:jointleading} then shows that the pattern is exceptional, with exactly two
groups of equal coordinates, of sizes two and three. After scaling, and in the order
of~\eqref{eq:jointexceptional}, $v$ is one of
\[
 (0,1,1,0,1),\quad (0,1,1,0,0),\quad
 (0,0,1,1,0),\quad (0,1,0,0,1).
\]
\item For each of these four directions the resultant has degree at most six in $t$. This
follows from a row and column degree count in the Newton-basis Bezout matrix used in the proof
of Lemma~\ref{lem:jointleading}, in which coinciding leading nodes lower the degrees of the
corresponding entries. Hence
$\Psi^4$ has degree at most six, $\Psi(s(t))$ is affine in $t$, and its second difference
$\Psi(s+v)+\Psi(s-v)-2\Psi(s)$ vanishes.
\item Expanding this second difference, and using $f=0$ at $s\pm v$, gives respectively
\[
 -2\bar N_1s_3,\qquad 2(\bar N_1+\bar u_1)(k_1-1),\qquad 2(\bar u_2-1)(k_2-1),\qquad
 2\bigl((\bar u_1-\bar u_2)-(s_4-s_1)\bigr).
\]
The factors $\bar N_1$, $\bar N_1+\bar u_1$ and $\bar u_2-1$ are nonzero, so respectively
\[
 s_3=0,\qquad k_1=1,\qquad k_2=1,\qquad s_4-s_1=\bar u_1-\bar u_2.
\]
\end{enumerate}
The Lean modules \lean{FixedDividendDirectionExtraction}, \lean{FixedDividendLineResultant}
and \lean{FixedDividendLines} contain steps (2)--(4).
The first case is impossible at a good start, where $s_3\ge3$. The polynomial $G$ in
Lemma~\ref{lem:labels} satisfies, by expansion,
\[
 G\le H^5(12k_1+464),\qquad G\le H^5(12k_2+464).
\]
In the last obstruction, $k_1+k_2+k_3=\bar u_1-\bar u_2=(u_1-u_2)/N_0\le2X$.
Thus each surviving case gives $k_1\le2X$ or $k_2\le2X$.
The label bound $N_0q_i^4\le aG$, with $N_0\ge1$, now gives
\[
 q_i^4\le aH^5(24X+464).
\]
But $q_i\ge a/(4A)$ and $24X+464\le488X$ for $X\ge1$.
Condition~\eqref{eq:jointscale}, with $124928=4^4\cdot488$, gives
the strict reverse inequality.
\end{proof}

\subsection{Orders of the reconstructed functions}\label{sec:jointorders}

\begin{lemma}[Orders of the reconstructed functions]\label{lem:joint-infinity-orders}
Let $C$ be a retained nonexceptional component of positive degree
containing a good start.
At every point above infinity, with hyperplane multiplicity $e$, one has
\[
 \operatorname{ord}Q=-\frac{3e}{2},\qquad
 \operatorname{ord}E_0=\frac{3e}{2}.
\]
In particular, $E_0$ is a nonzero rational function vanishing at every
point above infinity.
\end{lemma}
\begin{proof}
The good start guarantees that the reconstruction denominators are
not identically zero on $C$. Thus $Q$ and $E_0$ are rational functions
on $C$, and retention makes both identities \eqref{eq:jointcuts}
identities in its function field.

\emph{The quotient order.}
Write
$P_4(x)=x^4+p_3x^3+p_2x^2+p_1x+p_0$.
Its Newton expression gives $\deg p_j\le4-j$, and hence
$\operatorname{ord}p_j\ge-(4-j)e$.
The first retained identity gives
\[
 \operatorname{ord}P_4(-Q)=\operatorname{ord}(aH_4/N_0)=-6e.
\]
If $\operatorname{ord}Q\ge-e$, every term of $P_4(-Q)$ would have order
at least $-4e$, a contradiction. Put $r=\operatorname{ord}Q<-e$.
For $0\le j<4$,
\[
 \operatorname{ord}\bigl(p_j(-Q)^j\bigr)-4r
 \ \ge\ (4-j)(-e-r)>0.
\]
Thus $Q^4$ is the unique term of lowest order. It follows that
$4r=-6e$, or $\operatorname{ord}Q=-3e/2$.

\emph{The reconstructed curvature order.}
The coefficient $p_3=\bar u_2-1-s_1-s_2-s_3$ has order at least $-e$.
The raw-moment recurrence also shows that $H_{013}$ has degree at most
four: its degree-four part is
$(2k_0+k_1)k_0k_1(k_0+k_1)$.
Consequently
\[
 \operatorname{ord}(Q-p_3)=-3e/2,\qquad
 \operatorname{ord}\bigl((Q-p_3)H_3\bigr)=-9e/2
       <-4e\le\operatorname{ord}H_{013}.
\]
There is again a unique term of lowest order. Dividing the second
retained identity by $P_4(-Q)$ therefore gives
\[
 \operatorname{ord}E_0=-9e/2-(-6e)=3e/2.
\]
The argument uses the discrete valuation on the normalization, so it
also covers branches above singular points of the original curve.
\end{proof}

\subsection{The joint coordinates}\label{sec:jointcoords}

Put $d_j=b_{i+j}-b_{i+j+1}$ and
\[
 D_{01}=d_0k_1-d_1k_0,\quad D_{02}=d_0k_2-d_2k_0,\quad
 D_{12}=d_1k_2-d_2k_1.
\]
The first of these, $D_{01}$, is the curvature $I_3(i)$. With $s_j=q_{i+j}-q_i$ and
$r_j=k_j(s_j+s_{j+1})+k_{j+1}$, the Hankel polynomial $H_3$ of Section~\ref{sec:jointrecon}
and its shift are
\begin{align*}
 H_3&=k_0r_1-k_1r_0=\sigma(k_0,k_1,k_2),\\
 H_3'&=k_1r_2-k_2r_1=\sigma(k_1,k_2,k_3),
\end{align*}
with $\sigma$ as in Section~\ref{sec:six}. They are positive at every state at which they are
defined. On a retained component, use the reconstructed function $E_0$ to define
\[
 E_{12}=\frac{H_3'E_0-N_0k_1}{H_3},\qquad
 E_{02}=\frac{k_0E_{12}+k_2E_0}{k_1},\qquad
 W=(E_{12},-E_{02},E_0),\qquad J=N_0\Psi.
\]
The determinant and cross-product identities give, identically,
\begin{equation}\label{eq:jointdot}
 W\cdot(k_0,k_1,k_2)=0,\qquad W\cdot(r_0,r_1,r_2)=N_0.
\end{equation}
At a good start these rational functions take the integer values
$W=(D_{12},-D_{02},D_{01})$. Changing
the rank-five Gram basis to $1,x,x^2,x^3,P_4$ gives
\[
 \det\bigl(\nu_i(x^{r+s})\bigr)_{r,s<5}=\det\bigl(\nu_i(x^{r+s})\bigr)_{r,s<4}\,\nu_i(P_4^2).
\]
Consequently $J=I_5(i)$ is a positive integer at good starts. The heights of these
integers are controlled by
\[
 S=\left\lfloor\frac{4352 A^3H^3}{a^2}\right\rfloor,\qquad
 Y=\left\lfloor\frac{1259991040 A^5H^{10}}{a^4}\right\rfloor.
\]

\begin{lemma}[Heights of the joint coordinates]\label{lem:jointheights}
At good starts, $|W_j|\le S$ and $1\le J\le Y$.
\end{lemma}
\begin{proof}
Write $q=q_i$, and recall $q\ge a/(4A)$ (proof of Lemma~\ref{lem:labels}). For $j\le3$, apply
$\mu_{i+j}$ to the identity $(q+y)^{-1}=q^{-1}-yq^{-2}+y^2/(q^2(q+y))$ with $y=s_j+x$, as in
Lemma~\ref{lem:expand}. Since $\mu_{i+j}((s_j+x)^2)=s_j^2+2s_j+k_j+1\le17H^2$, this gives
\[
 \frac{b_{i+j}}a=\frac1q-\frac{s_j+1}{q^2}+\epsilon_j,
 \qquad 0\le\epsilon_j\le17H^2/q^3.
\]
Subtracting consecutive expansions, $d_j/a=k_j/q^2+\delta_j$ with $|\delta_j|\le34H^2/q^3$.
In $D_{jl}=d_jk_l-d_lk_j$ the leading terms cancel, so $|D_{jl}|q^3\le68aH^3$, and
$q\ge a/(4A)$ gives $|D_{jl}|\le S$, including the floor because the minors are integers.

For $J=I_5(i)$, argue as in the proof of Lemma~\ref{lem:labels} with five Newton polynomials
instead of four: Hadamard's inequality and $\nu_i\le\mu_i/q$ give
$Jq^5\le a\prod_{j<5}\mu_i(F_j^2)$. By the shift rule,
$\mu_i(F_j^2)=\mu_{i+j}\bigl(F_j(x+s_j)\bigr)$ is a polynomial of degree $j$ in the gaps, and
the gap bound gives $\mu_i(F_j^2)\le c_jH^j$ with $(c_0,\dots,c_4)=(1,2,10,77,799)$. Hence
$Jq^5\le1230460\,aH^{10}$, and $q\ge a/(4A)$ gives $J\le Y$.
\end{proof}

\begin{corollary}[Orders of the joint coordinates]\label{cor:joint-curvature-orders}
In the setting of Lemma~\ref{lem:joint-infinity-orders}, the joint
coordinates satisfy
\begin{equation}\label{eq:jointorders}
 \operatorname{ord}W_0=\operatorname{ord}W_2=3e/2,\qquad
 \operatorname{ord}W_1\ge3e/2,\qquad
 \operatorname{ord}J=-5e/2.
\end{equation}
\end{corollary}
\begin{proof}
The cubic leading part of $H_3'$ is $k_1k_2(k_1+k_2)$.
By \eqref{eq:joint-node-orders}, it has order $-3e$, whereas the
quadratic remainder $k_1k_3-k_2^2$ has order at least $-2e$.
Thus $\operatorname{ord}H_3=\operatorname{ord}H_3'=-3e$ and
$\operatorname{ord}k_1=-e$.
In the numerator defining $E_{12}$ the two terms have different orders:
\[
 \operatorname{ord}(H_3'E_0)=-3e/2
       <-e=\operatorname{ord}(N_0k_1).
\]
Therefore $\operatorname{ord}E_{12}=3e/2$.
Each term in $k_0E_{12}+k_2E_0$ has order $e/2$, so their sum has
order at least $e/2$. Dividing by $k_1$ gives
$\operatorname{ord}E_{02}\ge3e/2$. Cancellation here is harmless:
equality for the first and third coordinates already gives
$\min_j\operatorname{ord}W_j=3e/2$.
The formula for $J=N_0\Psi$ is Corollary~\ref{cor:jointparity}.
\end{proof}

\subsection{Counting the nonexceptional conics}\label{sec:jointconic}

\begin{proposition}[Conics]\label{prop:jointconic}
A nonexceptional component of degree two contains the free-node triples of $O_\eta(a^\eta)$
good starts, with an implied constant independent of the pattern and the component.
Consequently all nonexceptional conics together contribute $O_\eta(a^\eta X^5)$ good starts.
\end{proposition}

The idea is short. A component that is not retained contributes $O(1)$ good starts
(Lemma~\ref{lem:jointcuts}). A retained conic through three free-node triples of good starts
lies in a plane, and Corollary~\ref{cor:jointparity} leaves it a single rational branch at
infinity, of multiplicity two. So it is a parabola, the image of a quadratic polynomial map
$t\mapsto(\xi_0(t),\xi_1(t),\xi_2(t))/\Delta_c$ whose parameter $t$ is an integer at every
integer point. By Lemma~\ref{lem:joint-infinity-orders}, the reconstructed curvature $E_0$ is
a rational function of $t$ with a zero of order three at infinity, and it takes the integer
value $I_3(i)$ at each good start. Its nonconstant denominator must then divide a fixed nonzero
resultant at each such $t$, which leaves $a^{o(1)}$ possibilities. Appendix~\ref{app:conic}
gives the details, including the uniform bounds on the coefficients.

\subsection{The extra degree and the powered graph}\label{sec:jointgraph}

The conic count leaves retained nonexceptional components of degree at
least four. Fix such a component $C$ of degree $d$. The additional
integer coordinates will increase its degree without increasing their
heights too much. The key is one extra zero over the affine part of $C$.

All divisors in this argument are on the normalization $\widetilde C$
of the projective closure of $C$, over $\mathbb Q$.
For a closed point $P$ of $\widetilde C$, write
$\deg P=[\kappa(P):\mathbb Q]$, where $\kappa(P)$ is the residue field; thus
$\deg(\sum_P a_P[P])=\sum_P a_P\deg P$.
Let $E_\infty=\sum_{P\mid\infty}e_P[P]$ be the pullback of the original
hyperplane at infinity. Then
$\sum_{P\mid\infty}e_P\deg P=d$.  We do not require the closed points, or the
branches over a rational point of $C$, to be rational.

Let $\mathcal F$ be the pole envelope of the curvature coordinates:
\[
 \mathcal F=\sum_P \max\{0,-\operatorname{ord}_P(W_0),
                   -\operatorname{ord}_P(W_1),
                   -\operatorname{ord}_P(W_2)\}[P],
\]
with $\operatorname{ord}_P(0)=+\infty$.
The orders in \eqref{eq:jointorders} show that $\mathcal F$ is supported over the
affine part of $C$, and hence is disjoint from $E_\infty$.

\begin{lemma}[An additional affine zero]\label{lem:jointaffinezero}
There is a rational linear combination $\phi$ of $W_0,W_1,W_2$ which is
not the zero function and has a positive-order zero at a closed point
$P_0$ of $\widetilde C$ outside $E_\infty$.
\end{lemma}

\begin{proof}
First, the $W_j$ cannot all be scalar multiples, in $\mathbb Q(C)$, of a
fixed nonzero vector $v\in\mathbb Q^3$.  Such a relation, together with
$W\cdot(k_0,k_1,k_2)=0$, would put $C$ in the plane
$v_0k_0+v_1k_1+v_2k_2=0$.  This is a genuine plane: the first three gaps
are an invertible linear change of the three node coordinates.
The pattern quadratic is irreducible, so its restriction to this plane
is nonzero and has degree at most two.  Every curve contained in this
plane section therefore has degree at most two, contrary to $d\ge4$.

Choose the free-node triple $z_*\in C(\mathbb Q)$ of a good start.  All the $W_j$
are regular there, because their common reconstruction denominator
does not vanish at a good start.  Write $w_j=W_j(z_*)$.
The identity $W\cdot(r_0,r_1,r_2)=N_0\ne0$ gives $w\ne0$.
Choose $h$ with $w_h\ne0$.  At least one function
$w_hW_j-w_jW_h$ is nonzero: otherwise
$W=(W_h/w_h)w$, contradicting the preceding paragraph.
Take such a function for $\phi$.  It has rational coefficients and
satisfies $\phi(z_*)=0$.

This zero survives normalization even if $z_*$ is singular.
Indeed, write $\phi=g/g_0$ near $z_*$, with $g$ not identically zero on $C$,
$g(z_*)=0$ and $g_0(z_*)\ne0$.
The finite normalization has a closed point $P_0$ above $z_*$.
In its discrete valuation ring the numerator is in the maximal ideal
and the denominator is a unit, so $\operatorname{ord}_{P_0}(\phi)>0$.
Since $z_*$ is affine, $P_0$ does not lie in $E_\infty$.
\end{proof}

\begin{lemma}[The curvature pole degree]\label{lem:jointpoledegree}
The pole envelope satisfies
\begin{equation}\label{eq:jointextra}
 \deg \mathcal F\ge \frac{3d}{2}+1.
\end{equation}
\end{lemma}

\begin{proof}
Use the function $\phi$ from Lemma~\ref{lem:jointaffinezero}.
At every point over infinity, each $W_j$ has order at least $3e_P/2$.
The same is true of a rational linear combination; cancellation can
only increase its zero order.  Consequently
\[
 (\phi)_0\ \ge\ \frac32 E_\infty+[P_0],
 \qquad
 (\phi)_\infty\ \le\ \mathcal F .
\]
The first divisor is integral because every $e_P$ is even.
The function $\phi$ is nonzero and has a zero, so its principal divisor
has equal zero and pole degrees.  Hence
\[
 \deg \mathcal F\ \ge\ \deg(\phi)_\infty
 =\deg(\phi)_0
 \ \ge\ \frac32\sum_{P\mid\infty}e_P\deg P+\deg P_0
 \ \ge\ \frac{3d}{2}+1.
\]
The last inequality uses only $\deg P_0\ge1$, not that $P_0$ is rational.
\end{proof}

For positive integers $m,n$, let $\Gamma_{m,n}\subseteq\mathbb A^7$ be
the closure of the graph obtained by adjoining
$W_0^m,W_1^m,W_2^m,J^n$ to the three original node coordinates.
The graph is taken on the common domain of these rational functions.
It is an integral curve, and projection to the first three coordinates
is birational: on that domain the four remaining coordinates are
uniquely determined.  Thus $\Gamma_{m,n}$ and $C$ have the same function
field.  Taking powers does not lose this property, because the original
coordinates remain in the graph.

\begin{lemma}[Degree of the powered graph]\label{lem:jointpowereddegree}
If $\delta=\deg\Gamma_{m,n}$, then
\begin{equation}\label{eq:jointdegree}
 2\delta\ge(5n+3m)d+2m,
 \qquad
 \delta\ge10n+7m=:\delta_{\min}.
\end{equation}
Moreover,
\[
 \delta\le\bigl(247(m+n)+1\bigr)d.
\]
For $m=n=1$ one has the sharper upper bound $\delta\le247d$ and the
lower bound $\delta\ge4d+1\ge17$.
\end{lemma}

\begin{proof}
\emph{Lower bound.}
At a point $P$ over the old infinity hyperplane, the original node
coordinates have pole orders at most $e_P$, while
$J^n$ has pole order $5ne_P/2$ and all $W_j^m$ vanish.
Since $n\ge1$, the pole envelope of the seven graph coordinates at
$P$ is exactly $5ne_P/2$.
Over the affine part, the node coordinates and $J$ are regular:
$J=N_0\Psi$ is polynomial in the nodes.  There the envelope is exactly
$m$ times the curvature envelope.  The complete envelope of the graph
coordinates is therefore
\[
 \mathcal F_{m,n}=\frac{5n}{2}E_\infty+m\mathcal F.
\]

There is a single rational linear combination $\lambda$ of the seven
coordinates whose pole divisor is $\mathcal F_{m,n}$.  To see this, fix a point $P$
in the finite support of $\mathcal F_{m,n}$ and multiply all seven coordinates by a
local uniformizer to the power $\mathcal F_{m,n}(P)$.  Their residues in $\kappa(P)$
are not all zero.  Cancellation of the leading pole is therefore a
proper $\mathbb Q$-linear condition on the seven coefficients.
The finitely many such proper subspaces do not cover $\mathbb Q^7$.
Choosing coefficients outside their union gives the asserted
simultaneous absence of cancellation, also when $\kappa(P)\ne\mathbb Q$.

As a linear function on $\Gamma_{m,n}$, $\lambda$ has degree at most
$\delta$.  For example, homogenize it on the projective closure of the
graph: its pencil with the homogenizing coordinate consists of
hyperplane sections of degree $\delta$.  Removing common zeros can
only decrease the degree of the resulting map to $\mathbb P^1$.
Birational projection identifies its function field with $\mathbb Q(C)$,
so this degree is $\deg(\lambda)_\infty=\deg \mathcal F_{m,n}$.  It follows that
\[
 \delta\ \ge\ \frac{5n}{2}d+m\deg \mathcal F
 \ \ge\ \frac{5n+3m}{2}d+m.
\]
This is the first inequality in \eqref{eq:jointdegree}; $d\ge4$ gives
the second.  The auxiliary linear combination is used only to measure
degree: no bound on its coefficients enters the integer-point box.

\emph{Upper bound.} This is a routine count of polynomial degrees, given in
Appendix~\ref{app:degree}.
\end{proof}

Every good start whose free-node triple lies on $C$ gives an integer point on this graph.
By Lemma~\ref{lem:jointheights}, its coordinates have absolute values
at most
\[
 \max(3H,S^m,Y^n).
\]
The original node coordinates also show that distinct such node
triples remain distinct after adjoining the powers.

\subsection{The window bound and the final exponents}\label{sec:jointcount}

We use the coefficient-uniform affine-curve theorem of
Castryck--Cluckers--Dittmann--Nguyen~\cite[Theorem~3]{CCDN}, also stated as
\cite[Corollary~2.2]{CDHNV}: an integral affine curve $C$ over $\mathbb Q$ of degree $\delta$
in $\mathbb A^r$, with $r\ge2$ fixed, has
\begin{equation}\label{eq:jointexternal}
 N_{\rm aff}(C,B)\ll_r \delta^3 B^{1/\delta}(\log B+\delta),\qquad B\ge1,
\end{equation}
integer points in the box $[-B,B]^r$, with an implied constant
independent of the coefficients of $C$. Its formal proof is discussed in
Section~\ref{sec:lean}.

\begin{proposition}[Window bound]\label{prop:jointwindow}
Fix positive integers $m,n$ and $\eta>0$, and put $\mathcal B=\max(H,S^m,Y^n)$.
For a window run in $[A,2A)$ with $T\ge36$ satisfying~\eqref{eq:jointscale},
\begin{equation}\label{eq:jointcount}
 T\ll_{m,n,\eta}a^\eta\left(
 X^5+X^4H^{1/2}+X^5\mathcal B^{1/\delta_{\min}}\right).
\end{equation}
\end{proposition}

\begin{proof}
At least $T/2$ starts are good, and for each pattern the free-node triples of
its good starts are distinct (Lemma~\ref{lem:jointcuts}). By
Lemma~\ref{lem:jointcurve}, these triples lie on at most twenty-four integral
curves of degree at most twenty-four, and we bound the good starts on each
curve, allowing overlaps. All the heights $H$, $S$, $Y$ and $X$ are at most a fixed
power of $a$, so the logarithm in~\eqref{eq:jointexternal} is $O_\eta(a^\eta)$.
\begin{itemize}
\item \emph{Exceptional patterns.} There are $O_\eta(a^\eta X^4)$ of them
(Lemma~\ref{lem:jointcomponents}). Under~\eqref{eq:jointscale} no line
contains a good start, so every relevant curve has degree $2\le\delta\le24$, and
\eqref{eq:jointexternal} in $\mathbb A^3$ with $B=3H$ gives $O_\eta(a^\eta H^{1/2})$
good starts per pattern. This gives the term $X^4H^{1/2}$.
\item \emph{Nonexceptional patterns: components that are not retained, and
conics.} By Lemma~\ref{lem:jointcuts} and Section~\ref{sec:jointconic},
these contribute $O_\eta(a^\eta)$ per pattern, hence $O_\eta(a^\eta X^5)$ in all.
\item \emph{Remaining nonexceptional components.} Lines do not
occur, by Lemma~\ref{lem:jointcomponents}, and odd degrees do not occur, by
Corollary~\ref{cor:jointparity}, so these components have degree $d\ge4$.
Their good starts give distinct integer points of $\Gamma_{m,n}\subseteq\mathbb A^7$ in a box of
side $\max(3H,S^m,Y^n)\le3\mathcal B$. By Lemma~\ref{lem:jointpowereddegree}, the degree
$\delta$ of $\Gamma_{m,n}$ satisfies $\delta_{\min}\le\delta\le24(247(m+n)+1)$, so
\eqref{eq:jointexternal} gives $O_{m,n,\eta}(a^\eta\mathcal B^{1/\delta_{\min}})$
good starts per component. This gives the term $X^5\mathcal B^{1/\delta_{\min}}$.
\end{itemize}
The added coordinates $W_j$ and $J$ are functions of the free nodes on each
component. They therefore introduce no further freely counted coordinates, and
the pattern count is not enlarged.
\end{proof}

\begin{corollary}\label{cor:jointexponent}
For fixed positive integers $m,n$ and every $\varepsilon>0$,
\[
 P(a,b)\ll_{m,n,\varepsilon}a^{E(m,n)+\varepsilon},\qquad
 E(m,n)=\max\left(\frac{15\delta_{\min}+m}{51\delta_{\min}+3m},\
           \frac{15\delta_{\min}+6n}{51\delta_{\min}+20n}\right),
\]
where $\delta_{\min}=10n+7m$. In particular $E(1,1)=128/435$ and
$E(50,51)=2201/7480$.
\end{corollary}

\begin{proof}
We bound the number $T$ of terms in each dyadic window, eliminating the
window scale $A$. If $T<36$ there is nothing to prove. If $AT>28a$, then
Proposition~\ref{prop:six} gives $T^4\ll a$. Otherwise $HAT\le56a$, and
Proposition~\ref{prop:drop} gives $T^2\le8\Delta A$. Substituting $H\le56a/(AT)$ and then
$A\ge T^2/(8\Delta)$ into the definitions of $X$, $S$ and $Y$ gives
\begin{equation}\label{eq:jointscales}
 H\ll\frac{\Delta a}{T^3},\qquad X\ll\frac{\Delta^2a^3}{T^{10}},\qquad
 S\ll\frac{a}{T^3},\qquad Y\ll\frac{\Delta^5a^6}{T^{20}}.
\end{equation}
Every power of $\Delta$ is $O(a^\eta)$, and we absorb such factors into $a^{O(\eta)}$.
\begin{itemize}
\item If~\eqref{eq:jointscale} fails, then $124928A^4H^5X\ge a^3$. With
$Xa^3\le3609600A^4H^6$ this gives $A^8H^{11}\gg a^6$, and $H\le56a/(AT)$ turns it into
$A^3T^{11}\ll a^5$, so $T^{17}\ll\Delta^3a^5$.
\item Otherwise Proposition~\ref{prop:jointwindow} applies. By~\eqref{eq:jointscales}, its
first term gives $T\ll a^{O(\eta)}a^{15}/T^{50}$, so $T^{51}\ll a^{15+O(\eta)}$. Its second term
gives $T\ll a^{O(\eta)}a^{12+1/2}/T^{40+3/2}$, so $T^{85}\ll a^{25+O(\eta)}$.
\item For the graph term, $T\ll a^{O(\eta)}X^5\mathcal B^{1/\delta_{\min}}$ gives
$T^{51\delta_{\min}}\ll a^{O(\eta)}a^{15\delta_{\min}}\mathcal B$. Since $H\ge1$,
$\max(H,S^m)\le\max(H,S)^m\ll a^{O(\eta)}(a/T^3)^m$, so if $\mathcal B=\max(H,S^m)$ then
$T^{51\delta_{\min}+3m}\ll a^{15\delta_{\min}+m+O(\eta)}$. If $\mathcal B=Y^n$ then
$T^{51\delta_{\min}+20n}\ll a^{15\delta_{\min}+6n+O(\eta)}$.
\end{itemize}
The first three exponents all equal $5/17$, both fractions in $E(m,n)$ exceed $5/17$, and
$T^4\ll a$ is stronger still. Choosing $\eta$ small in terms of $\varepsilon$, $m$ and $n$, every
window has $T\ll a^{E(m,n)+\varepsilon/2}$, and summing over the $O(\log a)$ windows proves
the bound. For $m=n=1$ the two fractions are $256/870=128/435$ and $261/887$, and the first is
larger. For $(m,n)=(50,51)$ the larger fraction is $2201/7480$.
\end{proof}

\begin{proof}[Proof of Theorem~\ref{thm:joint}]
Set $r=n/m$. The two fractions in $E(m,n)$ become
\[
 E_S(r)=\frac{150r+106}{510r+360},\qquad
 E_Y(r)=\frac{156r+105}{530r+357}.
\]
The first decreases and the second increases on $r>0$. Their crossing satisfies
$30r^2-10r-21=0$, so $r_*=(5+\sqrt{655})/30$ and their common value is
$\alpha_*=(1310-\sqrt{655})/4365$.
Positive rational ratios approximate $r_*$; choose the fixed positive integers
$m,n$ to make $E(m,n)<\alpha_*+\varepsilon/2$, and apply
Corollary~\ref{cor:jointexponent} with loss $\varepsilon/2$.
The powers and all constants depend only on $\varepsilon$, not on $a,b$.
\end{proof}

\begin{remark}
In terms of the scales, at the critical point $T\approx a^{\alpha_*}$ one has
$A\approx a^{2\alpha_*}$, $H\approx S\approx a^{1-3\alpha_*}\approx a^{0.117}$,
$X\approx a^{3-10\alpha_*}\approx a^{0.057}$ and
$Y\approx a^{6-20\alpha_*}\approx a^{0.115}$. The pattern count $X^5\approx a^{0.287}$
alone would give the exponent $5/17$. The graph term multiplies it by
$\mathcal B^{1/\delta_{\min}}$, and at $n/m\approx r_*$ the two box radii $S^m$ and $Y^n$
balance, at $\mathcal B\approx a^{0.117m}$. This gives the extra factor
$a^{0.117/(10r_*+7)}\approx a^{0.0068}$.
\end{remark}

\section{Exterior lattice coordinates}\label{sec:lattice}

We now prove Theorem~\ref{thm:lattice}. We keep the notation and the conventions on constants
of Section~\ref{sec:joint}, and write $h=\delta_{\min}=10n+7m$ throughout. Only the count of
good starts on retained nonexceptional components of degree at least four changes.
Proposition~\ref{prop:latticewindow} replaces the term $X^5\mathcal B^{1/h}$ of
Proposition~\ref{prop:jointwindow} by
\[
 X^{5-1/(3h)}H^{1/h}+X^{5-2m/(3h)}S'^{\,m/h}+X^5Y^{n/h},
\]
where $S'$ differs from $S$ only in its constant (Lemma~\ref{lem:allcurv}). In effect $H$ is
replaced by $H/X^{1/3}$ and $S$ by $S/X^{2/3}$, while $Y$ is unchanged. The node term stays
below the critical exponent, so of the two fractions balanced in the proof of
Theorem~\ref{thm:joint} only the curvature fraction improves, and the crossing moves from
$\alpha_*=0.2942513\ldots$ to $\gamma_*=0.2942294\ldots$. Every pattern has
$\lambda(p)\ll X^{1/3}$ (Lemma~\ref{lem:index}), and the weighted count of
Lemma~\ref{lem:weights} shows that this saving is available on average: the patterns with an
unusually short vector are few enough to pay for their larger boxes.

Section~\ref{sec:latticepattern} attaches to each pattern $p$ a rank-three lattice
$\Lambda_p\subseteq\mathbb Z^3$ that contains the restrictions of the rows of
$\widehat{\mathcal M}_i$ to three pivot columns, and gives it a reduced basis. Section~\ref{sec:latticeshort} shows
that few patterns have a short vector in $\Lambda_p$. Section~\ref{sec:latticeexterior} writes
the nodes and the curvature coordinates in the reduced basis.
Section~\ref{sec:latticegraph} checks that this linear change of coordinates keeps the degree
bounds of Lemma~\ref{lem:jointpowereddegree}, and Section~\ref{sec:latticecount} sums over
patterns and optimizes the exponent.

\subsection{The lattice of a pattern}\label{sec:latticepattern}
Let $p\ne0$ be the minor vector of a $3\times5$ integer matrix $\mathcal M$, with pivot $\pi$
as in the proof of Proposition~\ref{prop:minorvec}; so $|p_\pi|$ is maximal among the
coordinates of $p$. Put $\kappa=|p_\pi|\ge1$, let $j_1<j_2<j_3$ be the pivot columns and
$j_4<j_5$ the other two, and call $x_\pi=(x_{j_1},x_{j_2},x_{j_3})$ the \emph{pivot part} of a
row $x$ of $\mathcal M$. Let $u,v\in\mathbb Z^3$ be as in that proof: $u_r$ ($v_r$) is the minor
with the $r$-th pivot column replaced by column $j_4$ ($j_5$), with the sign of the column
permutation. So the coordinates of $u$ and $v$ are, up to sign, coordinates of $p$, and
$p_\pi x_{j_4}=u\cdot x_\pi$ and $p_\pi x_{j_5}=v\cdot x_\pi$ for every row $x$
(Lemma~\ref{lem:lift}). Let $\Lambda_p\subseteq\mathbb Z^3$ be the subgroup of $x$ with
$\kappa\mid u\cdot x$ and $\kappa\mid v\cdot x$. It contains $\kappa\mathbb Z^3$. Write $|x|$
for the Euclidean norm and $|x|_\infty$ for the sup norm, and let the \emph{minimum}
$\lambda(p)$ be the least sup norm of a nonzero vector of $\Lambda_p$. Then
$|x|\ge\lambda(p)$ for every nonzero $x\in\Lambda_p$.

\begin{lemma}[Lifting]\label{lem:lift}
For a column $l$, let $L_l\in\mathbb Z^3$ be $p_\pi\mathbf e_r$ if $l=j_r$ is a pivot column,
$u$ if $l=j_4$ and $v$ if $l=j_5$, where $\mathbf e_1,\mathbf e_2,\mathbf e_3$ is the standard
basis. For every row $x$ of $\mathcal M$ and every column $l$, $p_\pi x_l=L_l\cdot x_\pi$.
Moreover $\det(L_0,L_1,L_2)=p_\pi^2p_{012}$,
and $|L_l|_\infty\le\kappa$.
\end{lemma}

\begin{proof}
The rows $L_l$ are the rows of $\mathcal M^{\mathsf T}\operatorname{adj}(\mathcal M_\pi^{\mathsf
T})$, where $\mathcal M_\pi$ is the pivot $3\times3$ block. This is Cramer's rule. Both
identities are polynomial identities in the fifteen entries, checked symbolically for each of
the ten pivots. The bound holds because the coordinates of $u$ and $v$ are, up to sign,
coordinates of $p$, and $|p_\pi|$ is maximal.
\end{proof}

So the pivot parts $x_\pi$ of the rows of $\mathcal M$ lie in $\Lambda_p$.

\begin{lemma}[Index]\label{lem:index}
If $K$ is a nonnegative integer and $(K+1)^3>\kappa$, then $\lambda(p)\le K$.
\end{lemma}

\begin{proof}
Write $g=\gcd(\kappa,u)$, $\ell=\kappa/g$ and $u=gu_0$. The proof of
Proposition~\ref{prop:minorvec} gives $v=\mu u_0+\ell y$ with $\mu\in\mathbb Z$ and
$y\in\mathbb Z^3$. Put
$\psi(x)=\bigl(u_0\cdot x\bmod \ell,\ (\mu\lfloor u_0\cdot x/\ell\rfloor+y\cdot x)\bmod g\bigr)$,
which takes at most $\ell g=\kappa$ values. The box $[0,K]^3$ has $(K+1)^3>\kappa$ points, so
by pigeonhole there are points $x\ne x'$ of $[0,K]^3$ with $\psi(x)=\psi(x')$. Then
$u_0\cdot(x-x')=\ell(f-f')$, where $f,f'$ are the two floors, so $u\cdot(x-x')=\kappa(f-f')$,
and $v\cdot(x-x')=\ell\bigl(\mu(f-f')+y\cdot(x-x')\bigr)$ with the bracket divisible by $g$.
Hence $x-x'$ is a nonzero vector of $\Lambda_p$ with sup norm at most $K$.
\end{proof}

The next lemma gives a reduced basis in dimension three. Its absolute
constant will not affect the exponents; the coordinate bound is what we use.

\begin{lemma}[Reduction]\label{lem:reduce}
Let $\Lambda\subseteq\mathbb Z^3$ be a subgroup containing $\kappa\mathbb Z^3$, $\kappa\ge1$.
There is a basis $\beta_1,\beta_2,\beta_3$ of $\Lambda$ with $V=\det(\beta_1,\beta_2,\beta_3)$,
$0<V\le\kappa^3$ and
\[
27\,|\beta_1|^2|\beta_2|^2|\beta_3|^2\le64\,V^2 .
\]
Hence every $x\in\Lambda$ is $\sum\xi_i\beta_i$ with $\xi\in\mathbb Z^3$ and
$27|\beta_i|^2\xi_i^2\le64|x|^2$.
\end{lemma}

The proof is given in Appendix~\ref{app:reduction}.

\subsection{Few patterns have a short vector}\label{sec:latticeshort}

\begin{lemma}[Divisibility count]\label{lem:dotcount}
For $\ell\ge1$ and $x\in\mathbb Z^3$, at most $27\ell^2\gcd(\ell,x)$ vectors
$y\in[-\ell,\ell]^3$ have $\ell\mid y\cdot x$. Here $\gcd(\ell,x)=\gcd(\ell,x_1,x_2,x_3)$.
\end{lemma}

\begin{proof}
Each residue class modulo $\ell$ meets $[-\ell,\ell]^3$ in at most $3^3$ points, so it suffices
to show that the set $\mathcal K$ of solutions in $[0,\ell)^3$ has at most $\ell^2e$ elements, where
$e=\gcd(\ell,x)$. By B\'ezout choose $z\in\mathbb Z^3$ with $z\cdot x\equiv e\pmod\ell$. The map
$(y,j)\mapsto(y+jz)\bmod\ell$ is injective on $\mathcal K\times[0,\ell/e)$: equal images give
$\ell\mid(j-j')e$, so $j=j'$, and then $y=y'$. Hence $|\mathcal K|\,\ell/e\le\ell^3$.
\end{proof}

\begin{proposition}[Short vectors]\label{prop:short}
Let $\mathcal S$ be a set of nonzero minor vectors with coordinates in $[-X,X]$, and let
$\tau(j)\le\tau_0$ for $1\le j\le X$. Then for every $\rho\ge1$,
\[
\#\{p\in\mathcal S:\lambda(p)\le\rho\}\le393660\,\rho^3X^4\tau_0^2 .
\]
\end{proposition}

\begin{proof}
If $\lambda(p)\le\rho$, some nonzero $x\in[-\rho,\rho]^3$ lies in $\Lambda_p$, so
$\kappa\mid u\cdot x$, where $\kappa=|p_\pi|\le X$; we drop the condition $\kappa\mid v\cdot x$. For fixed $x$, the encoding $p\mapsto(\pi,\operatorname{sgn}p_\pi,\kappa,u,v)$
of the proof of Proposition~\ref{prop:minorvec} gives at most $20\sum_{\kappa\le X}N(\kappa,x)$
such $p$, where $N(\kappa,x)$ counts pairs $(u,v)$ in $[-\kappa,\kappa]^3$ with
$\kappa\mid u\times v$ and $\kappa\mid u\cdot x$. By the proof of
Proposition~\ref{prop:minorvec}, for each $u$ there are at most $27\kappa\gcd(\kappa,u)^2$
vectors $v$. Group the $u$ by $g=\gcd(\kappa,u)$ and write $u=gu_0$ with $(\kappa/g)\mid u_0\cdot x$.
Lemma~\ref{lem:dotcount} with $\ell=\kappa/g$, and $\gcd(\kappa/g,x)\le\gcd(\kappa,x)$, give
\[
 \sum_ug^2\le\sum_{g\mid\kappa}g^2\cdot27\Bigl(\frac\kappa g\Bigr)^2\gcd(\kappa,x)
 =27\kappa^2\tau(\kappa)\gcd(\kappa,x),
\]
so $N(\kappa,x)\le729\kappa^3\tau(\kappa)\gcd(\kappa,x)$. Next,
$\gcd(\kappa,x)=\sum_{e\mid\gcd(\kappa,x)}\varphi(e)\le\sum_{e\mid\kappa,\,e\mid x}e$, and
swapping the sums,
\[
 \sum_{0<|x|_\infty\le\rho}\gcd(\kappa,x)
 \le\sum_{e\mid\kappa}e\,\#\{x\ne0:|x|_\infty\le\rho,\ e\mid x\}
 \le27\rho^3\tau(\kappa).
\]
Indeed, a divisor $e>\rho$ of $\kappa$ divides no nonzero $x$ in the box, and for $e\le\rho$
the multiples of $e$ number at most $(2\rho/e+1)^3$, with
$e(2\rho/e+1)^3\le27\rho^3$. Sum over $x$ and then over $\kappa\le X$,
using $\sum_{\kappa\le X}\kappa^3\le X^4$. The product of constants is
$20\cdot729\cdot27=393660$.
\end{proof}

\begin{lemma}[Fractional weights]\label{lem:weights}
Let $\mathcal S$ be a set of nonzero minor vectors with coordinates in $[-X,X]$, $X\ge1$, and
let $\tau_0\ge1$ satisfy $\tau(j)\le\tau_0$ for $1\le j\le X$. Then for $0\le t\le1$,
\[
 \sum_{p\in\mathcal S}\lambda(p)^{-t}\le3163861\,\tau_0^2X^{5-t/3}.
\]
\end{lemma}

\begin{proof}
Write $C_0=393660$ for the constant of Proposition~\ref{prop:short} and
$C_1=14581+8C_0=3163861$.
Choose an integer $K$ with $X\le K^3\le8X$. By Lemma~\ref{lem:index}, $1\le\lambda(p)\le K$
for $p\in\mathcal S$. Writing $1/\lambda=1/K+\sum_{\lambda\le j<K}1/(j(j+1))$ and using
Proposition~\ref{prop:short} with $\rho=j$,
\[
 K\sum_{p\in\mathcal S}\frac1{\lambda(p)}
 \le\#\mathcal S+C_0KX^4\tau_0^2\sum_{1\le j<K}\frac{j^2}{j+1}
 \le\#\mathcal S+C_0X^4\tau_0^2K^3 .
\]
The proof of Proposition~\ref{prop:minorvec} gives
$\#\mathcal S\le14581\,\tau_0X^5$. Since $K^3\le8X$ and $\tau_0\ge1$,
the right side is at most $C_1\tau_0^2X^5$. Dividing by $K\ge X^{1/3}$
proves the case $t=1$. For $0<t<1$, H\"older's inequality
with exponents $1/(1-t)$ and $1/t$ gives
$\sum\lambda^{-t}\le(\#\mathcal S)^{1-t}\bigl(\sum\lambda^{-1}\bigr)^t$, which is the claim;
the cases $t=0$ and $t=1$ are the two bounds above.
\end{proof}

\subsection{Exterior coordinates}\label{sec:latticeexterior}

\begin{lemma}[Exterior coordinates]\label{lem:exterior}
Let $\beta_1,\beta_2,\beta_3$ be a basis of $\Lambda\subseteq\mathbb Z^3$ as in
Lemma~\ref{lem:reduce}, with $|\beta_i|\ge\lambda>0$. Let $x,x'\in\Lambda$ have coordinates
$\xi,\xi'\in\mathbb Z^3$, and put $w=x\times x'$ and $c=\xi\times\xi'\in\mathbb Z^3$. Then
\[
 w=c_1\,\beta_2\times\beta_3+c_2\,\beta_3\times\beta_1+c_3\,\beta_1\times\beta_2,\qquad
 \beta_i\cdot w=Vc_i,\qquad 27\lambda^4c_i^2\le64|w|^2 .
\]
\end{lemma}

\begin{proof}
The first identity follows by expanding $(\sum\xi_i\beta_i)\times(\sum\xi'_i\beta_i)$
bilinearly, and taking the dot product with $\beta_i$ gives the second. By Cauchy--Schwarz and
Lemma~\ref{lem:reduce},
\[
 27|\beta_1|^2|\beta_2|^2|\beta_3|^2c_i^2\le64V^2c_i^2=64(\beta_i\cdot w)^2
 \le64|\beta_i|^2|w|^2 .
\]
Cancelling $|\beta_i|^2$ and using $|\beta_j|\ge\lambda$ for the other two vectors gives the
claim.
\end{proof}

The lemma replaces the bounds $|\xi|\ll|x|/\lambda$ and $|\xi'|\ll|x'|/\lambda$, which would give
$|c|\ll|x||x'|/\lambda^2$, by $|c|\ll|x\times x'|/\lambda^2$. This matters because the drop row
is large while the curvature minors are small.

The pivot columns of a pattern may include columns $3$ and $4$, so the cross product of the
pivot parts involves minors $D_{jl}$ with $l\ge3$. We therefore bound all ten minors, not only
the three of Lemma~\ref{lem:jointheights}.

\begin{lemma}[All curvature minors]\label{lem:allcurv}
At a good start, all ten minors $D_{jl}=d_jk_l-d_lk_j$ with $0\le j<l\le4$ satisfy
\[
 |D_{jl}|\le S':=\left\lfloor\frac{9472A^3H^3}{a^2}\right\rfloor .
\]
\end{lemma}

\begin{proof}
Argue as in Lemma~\ref{lem:jointheights}, using the expansions for $j\le5$, hence $d_j$ for
$j\le4$. The gaps $k_0,\dots,k_5$ are at most
$H$, so $s_j\le5H$ and
$\mu_{i+j}((s_j+x)^2)=s_j^2+2s_j+k_j+1\le(5^2+2\cdot5+2)H^2=37H^2$. This gives
$d_j/a=k_j/q^2+\delta_j$ with $|\delta_j|\le2\cdot37H^2/q^3$ for $j\le4$, hence
$|D_{jl}|q^3\le4\cdot37\,aH^3=148\,aH^3$, and $q\ge a/(4A)$ gives
$|D_{jl}|\le4^3\cdot148A^3H^3/a^2=9472A^3H^3/a^2$.
\end{proof}

Now fix a pattern $p$ that occurs at a good start, and let $\beta_1,\beta_2,\beta_3$ be a
basis of $\Lambda_p$ as in Lemma~\ref{lem:reduce}. Write $\lambda=\lambda(p)$. At a good start
$i$ with pattern $p$, the pivot parts $d_\pi$ and $k_\pi$ of the drop and gap rows of
$\widehat{\mathcal M}_i$ lie in $\Lambda_p$, with coordinates $\xi$ and $\zeta$. Put
\[
 c=\xi\times\zeta,\qquad w=d_\pi\times k_\pi.
\]
The entries of $w$ are, up to sign, the minors $D_{jl}$ with $j,l$ pivot columns. The next lemma
shows that $c=G_pW$ for a constant matrix $G_p$. We use $c$ rather than $G_pW$ because $c$ is
visibly integral, while $G_p$ has rational entries.

\begin{lemma}[Coordinates in the reduced basis]\label{lem:latticecoords}
There are matrices $B_p,G_p\in\mathrm{GL}_3(\mathbb Q)$, depending only on $p$, such that at
every good start with pattern $p$
\[
 (s_1,s_2,s_3)=B_p\zeta,\qquad c=G_pW,\qquad
 |\zeta|_\infty\le\frac{3H}{\lambda},\qquad |c|_\infty\le\frac{8S'}{3\lambda^2}.
\]
\end{lemma}

\begin{proof}
By Lemma~\ref{lem:lift}, $p_\pi(k_0,k_1,k_2)=\mathcal Lk_\pi$, where $\mathcal L$ has rows
$L_0,L_1,L_2$ and $\det\mathcal L=p_\pi^2p_{012}\ne0$, since $p_{012}=N_0=I_4(i)\ge1$ at a good
start. With $k_\pi=(\beta_1\,\beta_2\,\beta_3)\zeta$ and the cumulative sums $s_j=k_0+\dots+k_{j-1}$,
this gives $B_p$. Since $|k_\pi|^2\le3H^2$ and $|\beta_i|\ge\lambda$, Lemma~\ref{lem:reduce} gives
$|\zeta_i|\le(8/3)H/\lambda\le3H/\lambda$.

For the curvature, $W=(D_{12},-D_{02},D_{01})=(d_0,d_1,d_2)\times(k_0,k_1,k_2)$, and
$p_\pi(d_0,d_1,d_2)=\mathcal Ld_\pi$. Let $\operatorname{Cof}(\mathcal L)$ have rows
$L_1\times L_2$, $L_2\times L_0$, $L_0\times L_1$; its determinant is
$(\det\mathcal L)^2\ne0$. The Binet--Cauchy identity
$(L\cdot x)(L'\cdot y)-(L'\cdot x)(L\cdot y)=(L\times L')\cdot(x\times y)$, applied to pairs of
rows of $\mathcal L$, gives $(\mathcal Lx)\times(\mathcal Ly)=\operatorname{Cof}(\mathcal L)(x\times
y)$, so $p_\pi^2W=\operatorname{Cof}(\mathcal L)\,w$. Lemma~\ref{lem:exterior} gives
$Vc=(\beta_1\,\beta_2\,\beta_3)^{\mathsf T}w$. Together,
$c=p_\pi^2V^{-1}(\beta_1\,\beta_2\,\beta_3)^{\mathsf T}\operatorname{Cof}(\mathcal L)^{-1}W$,
which defines $G_p$. Finally $|w|^2\le3S'^2$ by Lemma~\ref{lem:allcurv}, and
Lemma~\ref{lem:exterior} gives $3\lambda^2|c_i|\le8S'$.
\end{proof}

\subsection{The lattice graph}\label{sec:latticegraph}
On a retained nonexceptional component $C$ of degree $d\ge4$ of a pattern $p$, regard
$c=G_pW$ as a vector of rational functions on $C$, and let $\Gamma'_{m,n}$ be the image of the
graph of
\[
 \bigl(s_1,s_2,s_3,c_1^m,c_2^m,c_3^m,J^n\bigr)
\]
under the linear map that replaces the nodes $s$ by $\zeta=B_p^{-1}s$.

\begin{lemma}[Degree of the lattice graph]\label{lem:latticedegree}
The degree $\delta'$ of $\Gamma'_{m,n}$ satisfies
$h\le\delta'\le\bigl(247(m+n)+1\bigr)d$.
\end{lemma}

\begin{proof}
An invertible linear change of the node coordinates is a linear automorphism of
$\mathbb A^7$, so it preserves the affine Hilbert function, and we may keep the nodes $s$. The
proof of Lemma~\ref{lem:jointpowereddegree}, with Appendix~\ref{app:degree}, uses the curvature
coordinates only through the following facts.
\begin{enumerate}
\item At every point $P$ over infinity, $\operatorname{ord}_PW_j\ge3e_P/2>0$ for each $j$
(Corollary~\ref{cor:joint-curvature-orders}).
\item The pole envelope $\mathcal F$ is supported over the affine part of $C$, and
$\deg\mathcal F\ge3d/2+1$ (Lemma~\ref{lem:jointpoledegree}).
\item Some rational linear combination of the seven graph coordinates has pole divisor
$\frac{5n}2E_\infty+m\mathcal F$; its construction uses only this envelope.
\item The $W_j$ have a common denominator $\Omega$ of degree at most $244$, with numerators of
degree at most $244$ (Appendix~\ref{app:degree}).
\end{enumerate}
Each of these is unchanged when $W$ is replaced by $c=G_pW$ for a constant invertible matrix
$G_p$. Each $c_j$ is a constant combination of the $W_l$, and each $W_l$ is a constant
combination of the $c_j$, so for every point $P$
\[
 \min_j\operatorname{ord}_Pc_j=\min_l\operatorname{ord}_P W_l .
\]
This gives (1) for $c$, and it shows that the pole envelope of $c$ is $\mathcal F$, which gives
(2) and (3). The numerators of $c$ over $\Omega$ are constant combinations of those of $W$,
which gives (4). The proof of Lemma~\ref{lem:jointpowereddegree} therefore applies unchanged.
\end{proof}

\subsection{The window bound and the optimized exponent}\label{sec:latticecount}

\begin{proposition}[Lattice window bound]\label{prop:latticewindow}
Fix positive integers $m,n$ and $\eta>0$. For a window run in $[A,2A)$ with $T\ge36$
satisfying~\eqref{eq:jointscale},
\[
 T\ll_{m,n,\eta}a^\eta\left(X^5+X^4H^{1/2}+X^{5-1/(3h)}H^{1/h}
 +X^{5-2m/(3h)}S'^{\,m/h}+X^5Y^{n/h}\right).
\]
\end{proposition}

\begin{proof}
The proof of Proposition~\ref{prop:jointwindow} applies, except for the remaining components of
the nonexceptional patterns. For such a component of a pattern $p$, the good starts give
integer points $(\zeta,c_1^m,c_2^m,c_3^m,J^n)$ of $\Gamma'_{m,n}$, and distinct node triples
give distinct points because $\zeta$ determines $s$. By Lemma~\ref{lem:latticecoords} they lie
in the box of side
\[
 \mathcal B_p=\max\Bigl(1,\frac{3H}{\lambda(p)},\Bigl(\frac{8S'}{3\lambda(p)^2}\Bigr)^m,Y^n\Bigr).
\]
By Lemma~\ref{lem:latticedegree} and~\eqref{eq:jointexternal}, there are
$O_{m,n,\eta}(a^\eta\mathcal B_p^{1/h})$ of them, and
\[
 \mathcal B_p^{1/h}\le1+3H^{1/h}\lambda(p)^{-1/h}+8S'^{\,m/h}\lambda(p)^{-2m/h}+Y^{n/h}.
\]
Sum over the $O_\eta(a^\eta X^5)$ patterns. Since $h\ge2m$, both $1/h$ and $2m/h$ lie in
$[0,1]$, and Lemma~\ref{lem:weights}, with $\tau_0=a^{o(1)}$, gives
$\sum_p\lambda(p)^{-1/h}\ll a^\eta X^{5-1/(3h)}$ and
$\sum_p\lambda(p)^{-2m/h}\ll a^\eta X^{5-2m/(3h)}$.
\end{proof}

\begin{corollary}\label{cor:latticeexponent}
For fixed positive integers $m,n$ and every $\varepsilon>0$,
\[
 P(a,b)\ll_{m,n,\varepsilon}a^{E'(m,n)+\varepsilon},\qquad
 E'(m,n)=\max\left(\frac5{17},\ \frac{45h}{153h-1},\ \frac{45h-3m}{153h-11m},\
 \frac{15h+6n}{51h+20n}\right),
\]
where $h=10n+7m$. In particular $E'(3,2)=153/520$.
\end{corollary}

\begin{proof}
Argue as in the proof of Corollary~\ref{cor:jointexponent}, with $S'\ll a/T^3$ in place of $S$
in~\eqref{eq:jointscales}. The first two terms and the windows in which~\eqref{eq:jointscale}
fails give the exponent $5/17$. For the third term, $T^{3h}\ll a^{O(\eta)}X^{15h-1}H^3$, and
\eqref{eq:jointscales} gives $T^{153h-1}\ll a^{45h+O(\eta)}$. For the fourth,
$T^{3h}\ll a^{O(\eta)}X^{15h-2m}S'^{\,3m}$ gives $T^{153h-11m}\ll a^{45h-3m+O(\eta)}$. For the
fifth, $T^h\ll a^{O(\eta)}X^{5h}Y^n$ gives $T^{51h+20n}\ll a^{15h+6n+O(\eta)}$. For
$(m,n)=(3,2)$, $h=41$ and the last three fractions are $1845/6272$, $153/520$ and $627/2131$.
\end{proof}

\begin{proof}[Proof of Theorem~\ref{thm:lattice}]
Set $r=n/m$. The last two fractions in $E'(m,n)$ become
\[
 E'_S(r)=\frac{450r+312}{1530r+1060},\qquad
 E_Y(r)=\frac{156r+105}{530r+357},
\]
where $E_Y$ is the same function as in the proof of Theorem~\ref{thm:joint}. The first decreases
and the second increases on $r>0$. Their crossing satisfies $15r^2=7$, so $r_*=\sqrt{105}/15$, and
their common value is $\gamma_*=471/1559-6\sqrt{105}/7795$. The fraction $45h/(153h-1)$ decreases
in $h$, and it is below $\gamma_*$ once $h>\gamma_*/(153\gamma_*-45)\approx17.2$; also
$5/17<\gamma_*$. Choose a positive rational $r$ close to $r_*$ with
$\max(E'_S(r),E_Y(r))<\gamma_*+\varepsilon/2$, write $r=n_0/m_0$, and take
$(m,n)=(6m_0,6n_0)$; any scaling with $h\ge18$ works, and the Lean proof uses $6$. Then
$h=6(10n_0+7m_0)\ge102$, and the fraction $45h/(153h-1)$ is at most $4590/15605<\gamma_*$. So
$E'(m,n)<\gamma_*+\varepsilon/2$, and Corollary~\ref{cor:latticeexponent} with loss
$\varepsilon/2$ gives the theorem.
\end{proof}

\section{Exact polar projection and the branch sieve}\label{sec:polar}

The curvature vector has an affine-linear projective direction. This gives a
better graph degree once two possible obstructions are controlled: cancellation
at its unique base point, and components of degree less than twenty-four.
A local intersection calculation handles the first obstruction; a sieve for
branches at infinity handles the second. We then count the additional
exceptional patterns and obtain Theorem~\ref{thm:polar}.

\subsection{An exact identity at orbit states}

For this section abbreviate the normalized parameters of~\eqref{eq:jointbar} by
\[
 (\rho,u,v)=(\bar N_1,\bar u_1,\bar u_2),\qquad
 (\rho',u',v')=(\bar N_2,\bar v_1,\bar v_2).
\]
Thus $\rho,u>0>v$ and
$\gamma=(\rho,-\rho-u,u-v,v-1,1)$; primes indicate the next pattern.
We retain the notation $N_0$, $\Psi$, $P_4$, $R_4$ and $D_{jl}$ from
Section~\ref{sec:joint}. Put $x=(k_0,k_1,k_2)^T$ and
\[
 G=\begin{pmatrix}
 \rho(\rho+1)&-\rho(u-1)&-\rho(v-1)\\
 -\rho(u-1)&u(u-1)&u(v-1)\\
 -\rho(v-1)&u(v-1)&v(v-1)
 \end{pmatrix},\qquad
 \ell=\begin{pmatrix}
 \rho(v-v')\\ \rho'-\rho-u(v-v')\\ u-u'-v(v-v')
 \end{pmatrix}.
\]
In these coordinates $f=x^TGx+\ell^Tx$, and
$\det G=\rho(\rho+u)(u-v)(v-1)\ne0$. Write
$W=(D_{12},-D_{02},D_{01})^T=d\times x$.

\begin{proposition}[Polar identity]\label{prop:polaridentity}
At every orbit state with initial index $i$ satisfying $i+7<P(a,b)$,
\[
 \Psi W=N_0(Gx+\ell),\qquad JW=N_0^2(Gx+\ell).
\]
In particular, $[W]=[Gx+\ell]$.
\end{proposition}
\begin{proof}
Set $V=Gx+\ell$. The quadratic equation gives $x\cdot V=0$.
Multiplying the five division equations by $\gamma_j$ and summing, the
relations $\sum\gamma_j=\sum\gamma_js_j=\sum\gamma_jb_j=0$ leave
\[
 \sum_{j=0}^4\gamma_js_jt_j+\sum_{j=0}^4\gamma_jd_j=0,
 \qquad t_j=b_0-b_j.
\]
The cumulative drop and gap identities identify this expression with
$d\cdot V$. For $r_j=k_j(s_j+s_{j+1})+k_{j+1}$, expansion also gives
\[
 V\cdot(r_0,r_1,r_2)-\Psi=-(s_4+1+v-v')f.
\]
Thus $V\cdot r=\Psi$, whereas $W\cdot r=N_0$.
Since $N_0>0$, the vectors $d,x$ are independent. Their common orthogonal
line contains both $V$ and $W$, so $\Psi W=N_0V$.
Now use $J=N_0\Psi$ and $\Psi>0$.
\end{proof}

The three homogeneous forms $Gx+\ell w$ have a unique common zero,
the finite point $x_*=-G^{-1}\ell$. It lies on $f=\Psi=0$.
It also lies on the fixed-dividend curve for every $a$. Indeed, its nodes are
\[
 s_0^*=0,\qquad
 s_j^*=\frac{\gamma'_{j-1}-\gamma_{j-1}}{\gamma_j}\ (1\le j\le3),
 \qquad s_4^*=v'-v.
\]
Dividing $Z(\xi)=\xi R_4(\xi)$ by $P_4(\xi)$ gives quotient $\xi+h$,
where $h=v'-v-s_4$, and a remainder with Newton coefficients
\[
 -h\gamma_0,\qquad
 \gamma'_{j-1}-\gamma_{j-1}-\gamma_j(s_j+h)\quad(1\le j\le3).
\]
They all vanish at $x_*$, so $Z_*=\xi P_{4,*}$ and
$N_0\operatorname{Res}(P_{4,*},Z_*)-a\Psi(x_*)^4=0$.

We add the three denominator-cleared polar identities to the retained
cuts of Section~\ref{sec:jointrecon}. Each has degree at most $250$.
A cut that is not identical on a degree-$d$ component removes all but
$250d$ of its points, by B\'ezout's theorem. Hence these additional cuts cost
$O(1)$ points per pattern. Throughout this section, a \emph{retained}
component satisfies the earlier cuts and these three new ones; on it,
\[
 JW=N_0^2(Gx+\ell)
\]
holds in the function field.

\subsection{The base point and the graph degree}

\begin{lemma}[The common-base divisor]\label{lem:polarbase}
There is a fixed rational polynomial $\Theta$ in the ten minor coordinates,
nonzero on the decomposable cone, with the following property. For every
pattern with $\Theta\ne0$ and every integral component $C$ of the pattern
curve, the common-base divisor of the homogeneous forms $Gx+\ell w$
on the normalization of $C$
has degree at most four. The polynomial $\Theta$ is independent of $a$.
\end{lemma}
\begin{proof}
Work first where $\ell_0\ne0$. Since $\nabla f(x_*)=-\ell$, the quadric is
smooth at $x_*$. Take $s=x_1-x_{*,1}$ and $t=x_2-x_{*,2}$ as local
coordinates and solve $f=0$ for $x_0$ by formal power series.
Let $R_1(\xi;s,t)$ be the linear part of the Newton remainder just computed,
and let $\Psi_1(s,t)$ be the linear part of $\Psi$ on this tangent plane.
The curve equation has order at least four and quartic initial form
\begin{equation}\label{eq:polarinitial}
 N_0N_4(s,t)-a\Psi_1(s,t)^4,\qquad
 N_4(s,t)=\operatorname{Res}_{\xi}(P_{4,*},R_1).
\end{equation}
To see that only first derivatives enter, compute the resultant as a
$4\times4$ B\'ezout determinant. Each row vanishes at $x_*$, so its term of
degree four uses exactly the linear part of every row. Higher terms in the
formal solution of $f=0$ cannot change~\eqref{eq:polarinitial}.

Consider the rational function of the pattern parameters
\[
 M=N_4(1,0)\Psi_1(0,1)^4-N_4(0,1)\Psi_1(1,0)^4.
\]
If $M\ne0$ and $N_0\ne0$, the two axis coefficients in
\eqref{eq:polarinitial} cannot both vanish, for any $a$.
Appendix~\ref{app:polarcertificates} gives the tangent frame, quartic and
first derivatives at
\[
 (\rho,u,v,\rho',u',v')=(2,5,-3,3,6,-5),\qquad
 x_*=(-1/7,-3/28,-1/2),
\]
where
\[
 M=\frac{26291112152499}{60236288}\ne0.
\]
These parameters arise from the maximal minors of the integer matrix
\[
 \begin{pmatrix}1&0&0&2&10\\0&1&0&-5&-22\\0&0&1&3&9\end{pmatrix}.
\]
Clear denominators in $M$ and the chart conditions, and include the necessary
nonzero denominator factors. Their product defines $\Theta$.
The displayed matrix witnesses its nonvanishing on the decomposable cone.

We now justify the passage from a quartic initial term to a divisor bound.
Let $S$ be the local ring of the smooth quadric at $x_*$ and let $F$ be the
curve equation. Choose $h=s$ or $h=t$ so that the complementary axis in
\eqref{eq:polarinitial} has nonzero coefficient. The completion of
$S/(h)$ is a power-series ring in one variable, in which $F$ has order
exactly four. Consequently
$\dim_{\mathbb Q}S/(F,h)=4$.
If $C$ passes through $x_*$, its local ring $R=S/I_C$ is a quotient with
$F\in I_C$, so
\[
 \dim_{\mathbb Q}R/hR\le4.
\]
In particular, $h$ is nonzero on this one-dimensional integral local ring.

Let $\overline R$ be its finite normalization and $Q=\overline R/R$.
The vector space $Q$ is finite-dimensional over $\mathbb Q$.
Apply multiplication by $h$ to
$0\to R\to\overline R\to Q\to0$. Multiplication is injective on the
first two rings, and its kernel and cokernel on $Q$ have equal dimensions.
The resulting exact sequence therefore gives
\[
 \dim_{\mathbb Q}\overline R/h\overline R
 =\dim_{\mathbb Q}R/hR
 =\sum_{P\mid x_*}[\kappa(P):\mathbb Q]\operatorname{ord}_P h.
\]
The final equality follows by factoring the ideal $(h)$ in the semilocal
Dedekind ring $\overline R$. Since $G$ is invertible, the coordinates of
$Gx+\ell=G(x-x_*)$ generate the centered maximal ideal. Their minimum order
at each $P$ is at most $\operatorname{ord}_P h$. The weighted sum of these
minimum orders is therefore at most four. A component missing $x_*$ has
no common-base contribution.
\end{proof}

Suppose now that a retained nonexceptional component has degree $24$ and
$\Theta\ne0$. At each point at infinity,
$\operatorname{ord}_P J=-5e_P/2$, by Corollary~\ref{cor:jointparity}, and
$\sum e_P[\kappa(P):\mathbb Q]=24$. Thus $J$ has pole divisor, and hence
zero divisor, of degree $60$; all its zeros are affine.
In the identity $W=N_0^2(Gx+\ell)/J$, cancellation at the affine zeros
of $J$ can cost at most the common-base degree. Lemma~\ref{lem:polarbase}
therefore gives an affine curvature pole envelope of degree at least
$60-4=56$. These poles are disjoint from the poles of $J$ at infinity,
so the pole formula of Lemma~\ref{lem:jointpowereddegree} gives
\begin{equation}\label{eq:polargraphdegree}
 \deg(\text{powered joint graph})\ge h:=60n+56m.
\end{equation}
An invertible constant change of curvature coordinates preserves the minimum
valuation, so this also holds for the lattice coordinates of
Section~\ref{sec:lattice}. It remains to show that components of degree
less than $24$ occur for relatively few patterns.

\subsection{An explicit label for a branch at infinity}

Let $P$ be a normalized branch at infinity of a nonexceptional component.
The five leading nodes are distinct by Lemma~\ref{lem:jointleading}.
Scale them to $(0,1,t,y,z)$, and put
\[
 z=-\gamma_1-\gamma_2t-\gamma_3y.
\]
The second and third moment equations become
\[
 Q_2(t,y)=\gamma_1+\gamma_2t^2+\gamma_3y^2+z^2=0,\qquad
 Q_3(t,y)=\gamma_1+\gamma_2t^3+\gamma_3y^3+z^3=0.
\]
Define the universal sextic $E_\gamma(t)$ by the polynomial identity
\begin{equation}\label{eq:branchsextic}
 \operatorname{Res}_y(Q_2,Q_3)=\gamma_3^2(\gamma_3+1)E_\gamma(t).
\end{equation}
The quotient on the right is a polynomial over
$\mathbb Z[\gamma_1,\gamma_2,\gamma_3]$, of degree six in $t$.
At an orbit pattern $\gamma_3=v-1<-1$, so
$\gamma_3^2(\gamma_3+1)\ne0$. Cancelling this factor in the resultant
and linear subresultant identities gives $E_\gamma(t)=0$ and
$d(t)y+n(t)=0$, where
\begin{align*}
 d(t)&=(\gamma_3-1)\bigl(\gamma_2(\gamma_2+\gamma_3+1)t^2
            +2\gamma_1\gamma_2t+\gamma_1(\gamma_1+\gamma_3+1)\bigr),\\
 n(t)&=\gamma_2(\gamma_2+1)(\gamma_2+\gamma_3+1)t^3
       +\gamma_1\gamma_2(3\gamma_2+\gamma_3+2)t^2\\
     &\quad+\gamma_1\gamma_2(3\gamma_1+\gamma_3+2)t
       +\gamma_1(\gamma_1+1)(\gamma_1+\gamma_3+1).
\end{align*}
These are polynomial identities, so the cleared relation remains valid when
a specialized denominator vanishes.

Write
\[
 \Delta=t y(t-1)(y-1)(y-t),\qquad V_5=z(z-1)(z-t)(z-y).
\]
By~\eqref{eq:joint-resultant-leading}, the degree-ten leading resultant at
these nodes is $R_0=\gamma_0\gamma_1\gamma_2\gamma_3\Delta V_5$.
The four moment equations give the barycentric identity
\begin{equation}\label{eq:branchbarycentric}
 R_0\Delta=V_5^4.
\end{equation}
For example, the identities
$\gamma_j\prod_{l\ne j}(z_j-z_l)=V_5$ for $0\le j<4$
follow by evaluating the first four moment relations; their product is
\eqref{eq:branchbarycentric}.
Let $w=1/s_1$ on this chart. The residue
$\eta_P=\operatorname{res}_P(w^5\Psi^2)$ is nonzero, because
$\operatorname{ord}_P\Psi=-5e_P/2$.
Taking leading terms in the curve equation gives
$N_0R_0=a\eta_P^2$. Thus the actual residue field contains
\begin{equation}\label{eq:branchsign}
 \beta=V_5^2/\eta_P,\qquad N_0\beta^2=a\Delta.
\end{equation}
This quadratic relation supplies information beyond the six possible
locations $t$.

For a slope $c$ define
\[
 D_\gamma(t)=t(-n(t))(t-1)(-n(t)-d(t))(-n(t)-td(t))
\]
and the label polynomial
\begin{equation}\label{eq:branchlabel}
 L_{a,N_0,c,\gamma}(Z)=
 \operatorname{Res}_t\bigl(E_\gamma(t),
       N_0d(t)^3(Z-t)^2-c^2aD_\gamma(t)\bigr).
\end{equation}
Here the resultant means the Sylvester determinant with degree bounds $6$
and $11$, padding missing coefficients with zero. This convention makes
\eqref{eq:branchlabel} a polynomial family under every specialization.
The relation $d(t)y+n(t)=0$ gives $D_\gamma(t)=d(t)^3\Delta$, so
\eqref{eq:branchsign} implies
\[
 L_{a,N_0,c,\gamma}(t+c\beta)=0
\]
in every actual residue field, including when $d(t)=0$.

\begin{lemma}[The degree-twelve label]\label{lem:branchlabel}
Fix $a\in\mathbb Q^\times$, and regard $N_0,\gamma_1,\gamma_2,\gamma_3$
as independent parameters. For some integer $c$ with $0\le c\le144$,
$L_{a,N_0,c,\gamma}$ is irreducible of degree twelve over
$\mathbb Q(N_0,\gamma_1,\gamma_2,\gamma_3)$.
For a specialized nonexceptional pattern, if one such label polynomial is
irreducible of degree twelve over $\mathbb Q$, every integral component of
its pattern curve has degree twenty-four.
\end{lemma}
\begin{proof}
At $(\gamma_1,\gamma_2,\gamma_3)=(-7,8,-4)$, the sextic in
\eqref{eq:branchsextic} equals $24E_0(t)$, where
\[
 E_0(t)=2400t^6-20160t^5+69776t^4-126448t^3
                  +125636t^2-64680t+13475.
\]
Appendix~\ref{app:polarcertificates} proves $E_0$ irreducible modulo $71$.
This specialization preserves degree six. Gauss's lemma therefore shows
that $E_\gamma$ is irreducible over
$K=\mathbb Q(\gamma_1,\gamma_2,\gamma_3)$: a factorization into two factors of positive $t$-degree over the
parameter ring would specialize to two factors of positive degree.
Choose a root $t$ of $E_\gamma$, so $[K(t):K]=6$.
In this field $d(t)\ne0$ since its degree is less than six, and
$y=-n(t)/d(t)$. The same degree comparison shows $\Delta\ne0$.

Adjoin the independent parameter $N_0$. The element $a\Delta/N_0$ has
odd valuation at $N_0=0$, so it is not a square in $K(t)(N_0)$.
Thus $(t,\beta)$ generates an extension of degree twelve over $K(N_0)$.
Among its twelve embeddings, equality of two values $t+c\beta$ excludes
at most one slope per pair. Consequently one of $0,\dots,144$ gives a
primitive element. Its minimal polynomial divides~\eqref{eq:branchlabel},
whose degree is at most twelve; hence that label polynomial is irreducible
of degree twelve. This argument holds for each fixed nonzero $a$; the
chosen slope may depend on $a$.

For an actual pattern, the residue field at any infinity branch contains
the root $t+c\beta$. Irreducibility of the specialized degree-twelve label
forces $[\kappa(P):\mathbb Q]\ge12$.
Corollary~\ref{cor:jointparity} gives $e_P\ge2$, so
\[
 \deg C=\sum_{P\text{ at infinity}}e_P[\kappa(P):\mathbb Q]\ge24.
\]
Lemma~\ref{lem:jointcurve} gives the reverse inequality.
\end{proof}

Call a pattern \emph{branch-good} if some rational slope $c$ makes its label
irreducible of degree twelve; otherwise call it \emph{branch-bad}.
Only the bounded set of integer slopes from Lemma~\ref{lem:branchlabel}
will be needed to estimate the latter set.

\subsection{The arithmetic sieve}

We use the following monic variant of effective Hilbert irreducibility.
Its proof follows the approach of~\cite[Theorem~1.10]{CDHNV}; the logarithmic
powers below are those of the monic variant proved here.

\begin{theorem}[Monic specialization estimate]\label{thm:monichilbert}
For fixed integers $r,d\ge1$ there is a constant $C_{r,d}$ with the following
property. Let $F\in\mathbb Z[T_1,\dots,T_r][Y]$ be monic in $Y$ and
irreducible over $\mathbb Q(T_1,\dots,T_r)$. Suppose
$\deg_YF\le d$, each coefficient of $F$ has total parameter degree at most
$d$, and every integer coefficient has absolute value at most $H\ge1$.
For every integer $B\ge1$,
\[
 \#\{t\in[-B,B]^r\cap\mathbb Z^r:
              F(t,Y)\text{ is reducible over }\mathbb Q\}
 \le C_{r,d}(1+\log(H+1))^4 B^{r-1/2}(1+\log B)^{4r+1}.
\]
\end{theorem}
\begin{proof}
First consider an irreducible $A(T,Z)\in\mathbb Z[T,Z]$ of bounded degree,
height $H_0$, and $Z$-degree $e\ge2$. Apart from the bounded number of zeros
of its leading coefficient, an integer root at $|T|\le B$ has
$|Z|\ll_d H_0B^d$, by the root bound. Apply the shear $Z\mapsto Z+T^k$,
with $k\gg_d1+\log(H_0+1)$. The transformed irreducible curve has degree
between $ke$ and $kd+d$, and the transformed points have height
$O_d(H_0B^k)$, taking $k>d$. The coefficient-uniform curve
bound~\eqref{eq:jointexternal} consequently gives
\[
 \#\{t\in[-B,B]\cap\mathbb Z:A(t,Z)\text{ has an integer root}\}
 \ll_d (1+\log(H_0+1))^4B^{1/e}(1+\log B).
\]
Indeed, $H_0^{1/(ke)}$ is bounded, and the degree and logarithmic factors
cost $O_d(k^4(1+\log B))$. The case $B=1$ is immediate.

For a monic $F(T,Y)$, a proper factor of a specialization is monic with
integer coefficients, by Gauss's lemma. Its coefficients are elementary
symmetric functions of a proper subset of the roots. For each generic
proper subset, at least one such coefficient is outside $\mathbb Q(T)$,
or the generic polynomial would factor. Its irreducible coefficient
equation has root degree at least two. Universal symmetric resolvents
bound the number and degrees of these equations in terms of $d$; integer
factor-height bounds bound their heights by $C_dH^{C_d}$.
Specialize in the integral closure of $\mathbb Z[T]$ in a splitting field.
Every specialized proper factor then supplies an integer root of one of
these equations, even when specialized roots coincide. The preceding root
estimate and a finite union give the one-parameter bound.

For one parameter with remaining polynomial variables, use Kronecker
substitution: encode the remaining variables and $Y$ by successive powers
of one variable, with base larger than their degree bounds and with $Y$
in the highest position. The encoded polynomial is still monic, with
degree bounded in terms of $r,d$. Factor it over $\mathbb Q(T_1)$.
Outside the one-parameter exceptional set, its irreducible factors remain
irreducible after specialization. A factorization of the original
specialization must therefore come from one of finitely many products of
these factors. Decode each candidate pair. If its product were the original
family identically, it would contradict generic irreducibility; otherwise
one nonzero coefficient polynomial in $T_1$ must vanish. There are only
$O_{r,d}(1)$ such extra parameter values. This proves the same square-root
bound for failure of irreducibility as a family in the remaining variables.

Finally induct on $r$. Exceptional values of $T_1$ contribute their count
times $(2B+1)^{r-1}$. At each remaining value, apply the induction hypothesis
to the specialized family, whose coefficient height is at most
$(d+1)HB^d$. Since
\[
 1+\log((d+1)HB^d+1)
 \ll_d(1+\log(H+1))(1+\log B),
\]
each induction step costs at most four further logarithmic powers.
The stated exponent $4r+1$ bounds all these losses.
\end{proof}

To apply the theorem with $a$ fixed, use the four integer minors of the
first four columns,
\[
 \mathbf v=(v_0,v_1,v_2,v_3)=(p_{012},p_{013},p_{023},p_{123}).
\]
For $v_0=N_0\ne0$ their relation to the independent parameters above is
\begin{equation}\label{eq:branchvertex}
 \gamma_1=-\frac{v_3+v_2}{v_0},\qquad
 \gamma_2=\frac{v_2+v_1}{v_0},\qquad
 \gamma_3=-\frac{v_1+v_0}{v_0}.
\end{equation}
This change is birational. Substitute it in~\eqref{eq:branchlabel} and
clear denominators by a fixed power of $v_0$. To make the resulting
degree-twelve polynomial monic without changing its generic irreducibility,
use the following polynomial operation. For $F(Y)=\sum_{i=0}^q b_iY^i$ put
\[
 \mathcal M_q(F)(Y)=Y^q+\sum_{i<q}b_i b_q^{q-1-i}Y^i.
\]
When $b_q\ne0$, this is $b_q^{q-1}F(Y/b_q)$.
When $b_q=0$ and $q\ge2$, it is $Y^{q-1}(Y+b_{q-1})$, hence reducible.
Thus irreducibility of a specialized monic normalization implies both
preservation of degree and irreducibility of the original label.

For each fixed positive integer $a$, choose the slope supplied by
Lemma~\ref{lem:branchlabel} and set $q=12$.
We obtain a generically irreducible monic integer family in $\mathbf v$,
with uniformly bounded degree and coefficient height at most $C(a+1)^D$,
for absolute constants $C,D$. The same is true after dilation
$\mathbf v=g\mathbf z$, with height at most $C(a+1)^Dg^D$.
The slope may depend on $a$, but the bounds are uniform because there are
only $145$ choices. Theorem~\ref{thm:monichilbert}, together with the
$O(B^3)$ zero-scale slice $v_0=0$, gives
\begin{equation}\label{eq:badvertices}
 \#\{\mathbf z\in[-B,B]^4\cap\mathbb Z^4:
                  g\mathbf z\text{ is branch-bad}\}
 \ll_\eta a^\eta B^{7/2}
\end{equation}
whenever $1\le g,B\ll a^7$. Vertices with $v_0=0$ are included among the bad vertices. This range
contains all the dilations and boxes used for orbit patterns below.

\begin{lemma}[Branch-bad full patterns]\label{lem:branchpatterns}
Among the selected orbit patterns of coordinate height at most $X$, the
number that are branch-bad is $O_\eta(a^\eta X^{9/2})$.
\end{lemma}
\begin{proof}
Partition the first-four-column vertices into dyadic bands with
$\max_j|v_j|\asymp M$, and choose a maximal coordinate as signed pivot $q$.
The other three vertex coordinates form $u$; put
$g=\gcd(q,u_1,u_2,u_3)$. Completing the fifth column amounts to choosing
$v\in[-X,X]^3$ with $q\mid u\times v$. The congruence count is
\[
 O(X^3g^2/M^2).
\]
In particular, these completion coordinates remain in the full $X$-box,
even when $M$ is small. Apply~\eqref{eq:badvertices} after dilation by $g$.
There are $O_\eta(a^\eta(M/g)^{7/2})$ bad vertices in this common-factor
class; imposing the exact gcd only reduces this count. Their total
contribution is at most
\[
 O_\eta\left(a^\eta X^3M^{3/2}\sum_{g\ge1}g^{-3/2}\right).
\]
The series converges, and summing the bands $M\le X$ proves the claim.
The finitely many pivot choices only change the constant.
\end{proof}

\subsection{Counting the geometric exceptional locus}

The new hypersurface $\Theta=0$ requires a saving over the ordinary
$O(X^5)$ pattern count. The following estimate includes all common factors.

\begin{proposition}[Hypersurface pattern count]\label{prop:polarpatterns}
Let $\Theta$ be a rational polynomial in the ten minor coordinates, with
$\Theta(\operatorname{plk}M)\ne0$ for some rational $3\times5$ matrix $M$.
For every $\eta>0$, the number of integral Pl\"ucker vectors of coordinate
height at most $X\ge1$ satisfying $\Theta=0$ is
\[
 O_{\Theta,\eta}(X^{14/3+\eta}).
\]
The count includes nonprimitive vectors and the zero vector.
\end{proposition}
\begin{proof}
Choose a minor of maximal absolute value $q>0$ and record its sign.
The six one-column replacements form $u,v\in[-q,q]^3$, and the other
coordinates are signed entries of $(u\times v)/q$.
There are twenty signed pivot charts. In each chart, clearing denominators
in $\Theta$ gives a nonzero polynomial $B(q,u,v)$.
Indeed, the product of the pivot minor and $\Theta(\operatorname{plk}M)$
is a nonzero polynomial in the matrix entries, so it is nonzero at some
rational matrix; every pivot chart meets the nonvanishing locus.

Put $g=\gcd(q,u_1,u_2,u_3)$, $n=q/g$ and $u_0=u/g$.
Since $\gcd(n,u_0)=1$, the congruence $u_0\times v\equiv0\pmod n$
is equivalent to
\[
 v\in\Lambda:=\mathbb Zu_0+n\mathbb Z^3.
\]
Let $\mu(u_0)$ be the least nonzero sup norm in the normal lattice
$\Lambda_n^\perp=\{z\in\mathbb Z^3:n\mid z\cdot u_0\}$.
If $x,y\in\Lambda$ are independent, then $x\times y=nz$ for a nonzero
$z\in\Lambda_n^\perp$. Hence $|x||y|\gg n\mu(u_0)$.
For a reduced basis of $\Lambda$ with lengths $\ell_1,\ell_2,\ell_3$,
this gives $\ell_i\ell_j\gg n\mu(u_0)$ for each pair.
The basis lengths are $O(n)$: apply the coordinate bound of
Lemma~\ref{lem:reduce} to $n$ times each standard basis vector, which lies
in $\Lambda$. For each basis coordinate, at least one of these three
integer coordinates is nonzero.
By the coordinate bounds of Lemma~\ref{lem:reduce}, vectors with
$|v|_\infty\le q$ have basis coordinates in a box of side lengths
$O(q/\ell_i)$; these lengths are bounded below because $q\ge n$.
For a nonzero polynomial fiber $B(q,u,\cdot)$, the polynomial grid-zero
bound therefore gives
\begin{equation}\label{eq:polarfiber}
 \#\{v\in\Lambda\cap[-q,q]^3:B(q,u,v)=0\}
 \ll_B\sum_{i<j}\frac{q^2}{\ell_i\ell_j}
 \ll_B\frac{q^2}{n\mu(u_0)}.
\end{equation}

Here is the reciprocal estimate that produces the exponent $14/3$:
\begin{equation}\label{eq:normalreciprocal}
 \sum_{u_0\in[-n,n]^3\cap\mathbb Z^3}\frac1{\mu(u_0)}
 \ll\tau(n)n^{8/3}.
\end{equation}
It holds for all $u_0$, without a primitivity restriction.
Pigeonholing $z\mapsto u_0\cdot z\pmod n$ gives $\mu(u_0)\ll n^{1/3}$.
For $R\ge1$, Lemma~\ref{lem:dotcount} and a divisor sum give
\begin{align*}
 \#\{u_0:\mu(u_0)\le R\}
 &\le\sum_{0<|z|_\infty\le R}\#\{u_0\in[-n,n]^3:n\mid u_0\cdot z\}\\
 &\ll n^2\sum_{0<|z|_\infty\le R}\gcd(n,z_1,z_2,z_3)
 \ll\tau(n)n^2R^3.
\end{align*}
For the last bound, expand the gcd as a divisor sum and note that a
nonzero $z$ of size at most $R$ has no common divisor greater than $R$.
Partial summation up to $K\asymp n^{1/3}$ now bounds the reciprocal sum by
$O(n^3/K+\tau(n)n^2K^2)$, proving~\eqref{eq:normalreciprocal}.

For fixed $q,g$, equations~\eqref{eq:polarfiber} and
\eqref{eq:normalreciprocal} bound the nonzero fibers by
$O_B(\tau(n)q^2n^{5/3})$. Sum over $g\mid q$ and $q\le X$ to obtain
\[
 O_B\bigl(\tau_{\max}(X)^2X^{14/3}\bigr),\qquad
 \tau_{\max}(X)=\max_{1\le j\le X}\tau(j).
\]

Finally, an identically zero fiber forces one fixed nonzero coefficient
polynomial $B_0(q,u)$ to vanish. For fixed $g$, the polynomial
$B_0(gq',gu')$ is still nonzero. The four-variable grid-zero bound supplies
$O_B((X/g)^3)$ possible scaled vertices. Each has
$O(qg^2)\le O(Xg^2)$ completions in its $q$-box, so these fibers total
\[
 O_B\left(X^4\sum_{g\le X}\frac1g\right)
 =O_B(X^4\log(2X)).
\]
The divisor estimate absorbs $\tau_{\max}(X)^2$ into $X^\eta$.
Summing the charts and adding the zero vector completes the proof.
\end{proof}

\subsection{The window estimate and the polar exponent}

Retain the window notation $[A,2A)$ of the earlier sections. For $T\ge36$
consecutive positive terms, put
\[
 H=\left\lfloor\frac{28a}{AT}\right\rfloor+1,\quad
 X=\left\lfloor\frac{3609600A^4H^6}{a^3}\right\rfloor,\quad
 S'=\left\lfloor\frac{9472A^3H^3}{a^2}\right\rfloor,\quad
 Y=\left\lfloor\frac{1259991040A^5H^{10}}{a^4}\right\rfloor.
\]
Only windows satisfying~\eqref{eq:jointscale} need the new estimate.
As before, call $R:\mathbb N\to\mathbb N$ slowly growing if
$R(a)\ll_\eta a^\eta$ for every $\eta>0$.

\begin{proposition}[Polar window estimate]\label{prop:polarwindow}
For fixed positive integers $m,n$, set $h=60n+56m$.
There is a slowly growing $R$, depending only on $m,n$, such that every
window above satisfying~\eqref{eq:jointscale} obeys
\begin{equation}\label{eq:polarsix}
\begin{split}
 T\le R(a)\bigl(&X^5+X^4H^{1/2}
 +X^{14/3}\max(H,S',Y)^{1/17}\\
 &+X^{5-1/(3h)}H^{1/h}
 +X^{5-2m/(3h)}(S')^{m/h}+X^5Y^{n/h}\bigr).
\end{split}
\end{equation}
\end{proposition}
\begin{proof}
The earlier exceptional patterns, conics and bounded cuts contribute
$X^5+X^4H^{1/2}$, as in Proposition~\ref{prop:latticewindow}.
Among the remaining patterns, Lemma~\ref{lem:branchpatterns} and
Proposition~\ref{prop:polarpatterns} bound those that are branch-bad or
satisfy $\Theta=0$ by $O_\eta(a^\eta X^{14/3})$.
Here $X\ll a^7$, so the extra $X^\eta$ loss is absorbed into $a^\eta$.
Use the earlier degree-$17$ unpowered graph for these patterns, giving the
third term of~\eqref{eq:polarsix}.

For each remaining retained component, Lemma~\ref{lem:branchlabel}
gives degree $24$, and~\eqref{eq:polargraphdegree} gives graph degree
at least $h$. Its upper degree is bounded in terms of $m,n$ by
Lemma~\ref{lem:jointpowereddegree}. The coefficient-uniform curve bound
and Lemma~\ref{lem:latticecoords} therefore give, for a pattern with shortest
lattice length $\lambda$, a contribution bounded by a slowly growing
factor times
\[
 (H/\lambda)^{1/h}+(S'/\lambda^2)^{m/h}+Y^{n/h}.
\]
Summing with Lemma~\ref{lem:weights}, at exponents $1/h$ and $2m/h$,
produces the last three terms. All omitted constants and logarithmic
factors depend only on the fixed powers and fit into $R$.
\end{proof}

\begin{theorem}[Polar bound]\label{thm:polarconditional}
For every fixed pair of positive integers $m,n$ and every $\varepsilon>0$,
\[
 P(a,b)\ll_{m,n,\varepsilon} a^{E_{\rm pol}(m,n)+\varepsilon}
 \quad(0<b<a),
\]
where, with $h=60n+56m$,
\[
 E_{\rm pol}(m,n)=\max\left\{\frac5{17},\frac{45h}{153h-1},
 \frac{45h-3m}{153h-11m},\frac{15h+6n}{51h+20n}\right\}.
\]
For $(m,n)=(3,2)$ this exponent is $4317/14677$.
Optimizing over the fixed powers gives
\[
 P(a,b)\ll_\varepsilon a^{\beta_*+\varepsilon},\qquad
 \beta_*=\frac{41301-3\sqrt{82849}}{137480}
        =0.2941336553\ldots.
\]
\end{theorem}
\begin{proof}
Apply the window argument of Section~\ref{sec:lattice} to
\eqref{eq:polarsix}. The three terms from the exceptional contribution
$X^{14/3}\max(H,S',Y)^{1/17}$ give exponents
$717/2440$, $717/2440$ and $732/2491$, all below $5/17$.
The other terms give the displayed maximum. The six-step and drop bounds
handle windows outside~\eqref{eq:jointscale}; dyadic summation and slow
growth absorb the remaining losses.

For $(m,n)=(3,2)$, $h=288$ and the last three fractions are
$12960/44063$, $4317/14677$ and $1083/3682$.
For optimization put $r=n/m$. The curvature and rank-five fractions become
\[
 E_S(r)=\frac{2700r+2517}{9180r+8557},\qquad
 E_Y(r)=\frac{906r+840}{3080r+2856}.
\]
The first decreases and the second increases. They cross at
$r_*=(\sqrt{82849}-47)/360$, with common value $\beta_*$.
Approximate $r_*$ by positive rationals and scale the corresponding
numerator and denominator until $45h/(153h-1)<\beta_*$.
The arbitrarily small exponent loss absorbs the approximation, proving
Theorem~\ref{thm:polar}.
\end{proof}

\section{Rises in the quotient increments}\label{sec:rises}

Throughout this section $k_i=q_{i+1}-q_i\ge1$, $d_i=b_i-b_{i+1}\ge1$, and a
\emph{rise} is an index $0\le i\le P(a,b)-3$ with $k_i<k_{i+1}$.

\begin{lemma}[Six quotients in arithmetic progression]\label{lem:ap}
Let $q,h\ge1$ and $x_0>x_1>\dots>x_6\ge0$ satisfy $x_j(q+jh)+x_{j+1}=a$ for $0\le j\le5$. Then
$q^4\le 324\,h^3a$.
\end{lemma}

\begin{proof}
Let $f_j=x_j-x_{j+1}$, and let $E=(f_1-f_2)h$ and $F=(f_2-f_3)h$ be the curvatures of the blocks
starting at $x_1$ and $x_2$ (all their increments equal $h$). With $U=f_0-2f_1+f_2$, the
division equations $\mathcal E_j\colon x_j(q+jh)+x_{j+1}-a=0$ give the exact identity
\[
qhU=(3h-1)E+F,\qquad\text{namely}\qquad qhU-(3h-1)E-F=h(\mathcal E_0-3\mathcal E_1+3\mathcal E_2-\mathcal E_3).
\]
This is the case $u=v=w=z=h$ of Lemma~\ref{lem:six}, where $L=2h^3U$. By
Lemma~\ref{lem:curv}, $E,F\ge1$, so the integer $U$ is at least $1$.

Now bound $E$ and $F$. Subtracting consecutive division equations gives
$q_jf_j\le hx_{j+1}$, where $q_j=q+jh$. Since $x_{j+1}q_{j+1}\le a$, this yields
$q^2f_j\le ha$ for $1\le j\le4$. By Lemma~\ref{lem:size},
$Ex_3\le(h+1)(f_1+f_2+f_3)^2\le 9h^2(h+1)a^2/q^4$. Also $a<x_3(q+3h+1)$, so
$Eq^4\le 9h^2(h+1)(q+3h+1)a$, and likewise $Fq^4\le 9h^2(h+1)(q+4h+1)a$. Hence
\[
q^5h\le q^4\bigl((3h-1)E+F\bigr)\le 27h^3(h+1)(q+4h+1)a .
\]
If $q\ge h$ this gives $q^4\le 324h^3a$. If $q<h$ then $q^4<h^4\le h^3a$, since $h\le q_1\le a$.
\end{proof}

\begin{lemma}[Drops decrease between rises]\label{lem:norise}
Whenever $b_{i+2}>0$, $d_{i+1}k_i<d_ik_{i+1}$. In particular, if $k_{i+1}\le k_i$ then
$d_{i+1}<d_i$.
\end{lemma}

\begin{proof}
This is $D=d_ik_{i+1}-d_{i+1}k_i\ge1$ (Lemma~\ref{lem:curv}). If $k_{i+1}\le k_i$ and
$d_{i+1}\ge d_i$, then $d_ik_{i+1}\le d_{i+1}k_i$, which is a contradiction.
\end{proof}

\begin{lemma}[Bounded multiplicity]\label{lem:mult}
Let $c_j\ge1$ be integers, $j\in J$, such that no value occurs more than $M$ times. Then
$|J|^2\le 2M\sum_{j\in J}c_j$.
\end{lemma}

\begin{proof}
Induct on $|J|$. Let $c_{j_0}$ be the largest value. The values lie in $[1,c_{j_0}]$, so
$|J|\le Mc_{j_0}$. Removing $j_0$ and using $|J|^2=(|J|-1)^2+2|J|-1$ completes the induction.
\end{proof}

\begin{proposition}[Rises in a window]\label{prop:rises}
Let $x_0,\dots,x_{T-1}$ be a window run in $[A,2A)$ with $T\ge2$, and let $\rho$ be the number
of rises among its $T-1$ increments. Then:
\begin{enumerate}
\item $(T-1)^2\le 2(\rho+1)(A-1)$;
\item $T-1\le\lfloor a/A\rfloor$;
\item if $6A\le a$, then $T-1=\ell+s$ with $\ell^2\le 8(\rho+1)\lfloor a/A\rfloor$ and
$s^3\le 26244\,A$.
\end{enumerate}
\end{proposition}

\begin{proof}
Split the increments into \emph{runs}: maximal stretches containing no rise. There are at most
$\rho+1$ runs.

(1) By Lemma~\ref{lem:norise}, the drops strictly decrease along a run, so each value of $d_j$
occurs at most once per run, hence at most $\rho+1$ times. Lemma~\ref{lem:mult} and
$\sum_jd_j=x_0-x_{T-1}\le A-1$ give the claim.

(2) The increments are $\ge1$, and they sum to $q_{T-1}-q_0\le\lfloor a/A\rfloor$.

(3) Call an increment $h$ \emph{low} if $26244h^3A^4<a^3$. Five consecutive equal increments
$h$ inside the run give six quotients in arithmetic progression. Lemma~\ref{lem:ap} then
gives $q^4\le324h^3a$. On the other hand, $a<3Aq$, which follows from $a=qx_0+x_1$,
$x_1<x_0<2A$ and $6A\le a$. Together these give $a^3\le26244h^3A^4$, so $h$ is not low.
Within a run the increments are nonincreasing, so equal values are consecutive. Hence each
low value occurs at most $4$ times per run, and at most $4(\rho+1)$ times in all.
Lemma~\ref{lem:mult} bounds the number $\ell$ of low increments by
$\ell^2\le8(\rho+1)\sum_jk_j\le8(\rho+1)\lfloor a/A\rfloor$. Let $s$ be the number of high
increments. If $s=0$, its bound is immediate. Otherwise let $h_*$ be their
smallest value. Then $sh_*\le\lfloor a/A\rfloor$ and
$a^3\le26244h_*^3A^4$, so $s^3a^3\le26244A^4(a/A)^3\le26244Aa^3$.
\end{proof}

\begin{lemma}[Geometric sums]\label{lem:geom}
Let $\alpha,\beta,s,t>0$ and $K\ge0$. Then
\[
\sum_{k<K}\min\bigl(\alpha2^{-sk},\ \beta2^{tk}\bigr)\le
\Bigl(\frac1{1-2^{-s}}+\frac{2^t}{2^t-1}\Bigr)\alpha^{t/(s+t)}\beta^{s/(s+t)} .
\]
\end{lemma}

\begin{proof}
Let $M=\alpha^{t/(s+t)}\beta^{s/(s+t)}$. Then $(\alpha2^{-sk})^t(\beta2^{tk})^s=M^{s+t}$, so for each
$k$ either $\beta2^{tk}\le M$ or $\alpha2^{-sk}<M$. The terms $\beta2^{tk}\le M$ grow with ratio
$2^t$, so they sum to at most $M2^t/(2^t-1)$. The terms $\alpha2^{-sk}<M$ decay with ratio $2^{-s}$,
so they sum to at most $M/(1-2^{-s})$.
\end{proof}

\begin{proof}[Proof of Theorem~\ref{thm:rises}]
The rises inside a window are rises of the orbit, so $\rho+1\le R$. Combining the three parts
of Proposition~\ref{prop:rises} gives for the
window $[2^k,2^{k+1})$ (when $6A>a$, part (2) gives $T-1\le a/A<6\le30A^{1/3}$)
\[
T_k\le 1+\sqrt R\,\min\bigl(\sqrt{8a}\,2^{-k/2},\ \sqrt2\,2^{k/2}\bigr)
+\min\bigl(a\,2^{-k},\ 30\cdot2^{k/3}\bigr).
\]
Lemma~\ref{lem:geom} gives
$\sum_k\min(\sqrt{8a}\,2^{-k/2},\sqrt2\,2^{k/2})\ll(\sqrt{8a}\sqrt2)^{1/2}=2a^{1/4}$. It also gives
$\sum_k\min(a2^{-k},30\cdot2^{k/3})\ll a^{1/4}$. Since $\log_2a+1\ll a^{1/4}$, summing over $k$
gives the theorem.
\end{proof}

\begin{remark}
By Theorem~\ref{thm:rises}, an orbit with $P(a,b)\asymp a^{3/10}$ must have $R\gg a^{1/10}$.
Proposition~\ref{prop:rises}(3) says more: inside the critical window $A\asymp a^{3/5}$, a
run of $T\asymp a^{3/10}$ remainders must contain $\gg a^{1/5}$ rises.
\end{remark}

\section{The Lean formalization}\label{sec:lean}

The source repository is
\url{https://github.com/jasondavies/pierce-length-bound}
(private during manuscript preparation).

The proof baseline is commit
\lean{27e6ac9ed3e49378064126fbe3c80b9b1f91ffd3}.
The development uses Lean~4 \lean{v4.35.0} and Mathlib \lean{v4.35.0}.
The unconditional results of this paper are proved without additional
mathematical hypotheses, including the curve-counting input and the
effective Hilbert specialization theorem required by
Theorem~\ref{thm:polar}.

The strongest bound is exported as
\begin{quote}\ttfamily\small
theorem P\_le\_polar\_optimal :\\
\hspace*{2em}$\forall$ $\varepsilon$ > (0 : $\mathbb R$), $\exists$ C > 0, $\forall$ a b : $\mathbb N$,\\
\hspace*{2em}0 < b $\to$ b < a $\to$ (P a b : $\mathbb R$) $\le$\\
\hspace*{2em}C * (a : $\mathbb R$) \^{} (((41301 : $\mathbb R$) - 3 * Real.sqrt 82849) / 137480 + $\varepsilon$)
\end{quote}
The quantifier order makes uniformity explicit: $C$ may depend on
$\varepsilon$, but not on $a$ or $b$. The function \lean{P a b} is defined
by well-founded recursion,
\[
 P(a,0)=0,\qquad P(a,b)=1+P(a,a\bmod b)\quad(b>0),
\]
and counts the positive terms of the zero-indexed orbit, as in the paper.
The rational bound $4317/14677$ is exported separately as
\lean{P\_le\_4317\_14677}. Both results are assembled in
\lean{PolarProjectionMain} from proved arithmetic, geometric and
counting estimates. The repository README gives the correspondence between
the other results and their Lean declarations.

The affine-curve theorem~\eqref{eq:jointexternal} is imported from the separate
\href{https://github.com/jasondavies/ccdn-linear-growth/tree/9ade6200d053ceef5f6ccc4154ae30c6056be150}{\lean{ccdn-linear-growth} library at revision \lean{9ade620}}.
That development proves the plane estimate, degree-preserving projection
and box rescaling, using the same affine Hilbert degree as this paper.
The present repository proves the required monic Hilbert irreducibility
estimate and its application to the actual branch-label family. Thus neither
counting input is an assumed proposition in the final polar bound.

The separate logarithmic refinement \lean{P\_le\_loglog} retains two explicit
hypotheses: the short-sum estimate for the Erd\H{o}s--Hooley function
from~\cite{NT}, over intervals $(x,x+y]$ with $y\ge x^{1/4}$, and the
mean-value estimate from~\cite{BT}. Their statements include the constants
and thresholds; their proofs are not formalized here.

The proof library contains no \lean{sorry} and uses no \lean{native\_decide}.
Its axiom audit permits only \lean{propext}, \lean{Classical.choice} and
\lean{Quot.sound}, the standard axioms used by Mathlib. Fifteen public
statements are compared with independent Mathlib-only specifications, and
their proofs are replayed by Lean, NanoDa and con-ron. The specification
file \lean{Challenge.lean} deliberately contains theorem holes; it is
separate from the proved solution. Build instructions, the full theorem
map, and the verification commands are recorded in the README.

The divisor estimate uses \lean{DivisorPowerBound.lean} from
\lean{openai/math}~\cite{OAIMath} under Apache-2.0, with unchanged proof
bodies and adapted imports and module visibility. The repository records
the source revision and the provenance of its other dependencies.

\section{Limits of the method}\label{sec:limits}

The polar refinement brings the exponent close to $5/17$, but does not
remove that threshold. The current local estimates still contain the
$X^5$ pattern contribution from bounded cuts and conics, the earlier
exceptional-pattern contribution, and the windows where the line-exclusion
condition~\eqref{eq:jointscale} fails. Each leads to exponent $5/17$, independently
of the coordinate powers $m,n$. The new sieve and geometric exceptional
locus fit below this threshold.

The improved graph degree gives
\[
 \beta_*-\frac5{17}=0.00001600826\ldots.
\]
Increasing the graph degree further can reduce this gap. Going below
$5/17$ requires an improvement in the other counting contributions, or a
different way to combine the information from successive orbit states.
This is a limitation of the present estimates, not a lower bound for the
Pierce iteration.

The earlier $3/10$ argument has a similar barrier. Even with no divisor
loss, its two window bounds
\[
 T\ll a^{1/2}A^{-1/3},\qquad T\ll A^{1/2}
\]
balance at $A\asymp a^{3/5}$ and permit $T\asymp a^{3/10}$. Better
constants or short-interval divisor estimates alone do not change this
exponent. The higher invariants and fixed-dividend equations supply
restrictions on configurations that this elementary balance misses.

The linear and quadratic pattern relations alone still admit families of
solutions. For example, take the Newton coefficients
\[
 (\alpha,-3\alpha-1,3\alpha+3,-\alpha-3,1)
\]
at every state, with coefficient polynomial $(z-1)^3(z-\alpha)$,
and constant increments $k_j=t$. All the linear and
quadratic relations hold for every $t$. The remainder equations exclude
such a free family; this is why retaining the fixed dividend is essential
to the curve construction.

Nor do the curvatures necessarily decrease. For $a=13242$, the orbit of
$255$ begins
\[
 255,\ 237,\ 207,\ 201,\ 177,\ 144,\ 138,\ 132,\ldots
\]
with quotients $51,55,63,65,74,91,95,100,\ldots$ and first five curvatures
$24,12,6,111,30$. The corresponding blocks all lie within a factor of two.
The rise bound gives one constraint on the order of the increments, but
the present argument does not capture all compatibility conditions along
an entire orbit. Shallit's conjectured $O((\log a)^2)$ bound remains open.

\appendix

\section{Conic components}\label{app:conic}

In this appendix we prove Proposition~\ref{prop:jointconic}. Let $C$ be a nonexceptional
component of degree two. A component that is not retained
already contributes at most $404\cdot2=808$ starts, by
Lemma~\ref{lem:jointcuts}. We may therefore assume that the component
$C$ is retained. If it contains at most two free-node triples of good starts,
there is nothing to prove. In the remaining case we first construct a
parameter whose value is an integer at every such triple.

\begin{lemma}[A bounded integral parameter]\label{lem:jointconicparameter}
Let $C$ be a retained nonexceptional degree-two component containing the
free-node triples of at least three good starts. After permuting the three
free-node coordinates to $(x,y,z)$, there are
an integer $\Delta_c\ne0$ and polynomials $\xi_0,\xi_1,\xi_2\in\mathbb Z[t]$ such
that
\[
 \deg \xi_j\le2,\qquad
 |\Delta_c|,\ \max_{j,k}|[t^k]\xi_j|\le(aX)^{O(1)}.
\]
The map
\[
 t\longmapsto
 \bigl(\xi_0(t),\xi_1(t),\xi_2(t)\bigr)/\Delta_c
\]
parametrizes $C$ over $\mathbb Q$. Its inverse takes an integer value at
every integer point of $C$, and is injective on those points.
\end{lemma}
\begin{proof}
Every free node of a good start has absolute value at most $a$. Three distinct
such triples admit an integer plane
\[
 \lambda_zz+\lambda_xx+\lambda_yy+\lambda_1=0,\qquad
 |\lambda_z|,|\lambda_x|,|\lambda_y|\le8a^2,\qquad |\lambda_1|\le24a^3.
\]
For example, take a cross product of two differences to obtain the
normal; if that cross product vanishes, a nonzero perpendicular vector
with the same bounds suffices. Put $K_{\mathrm p}=24a^3$. A plane not containing an integral curve
of degree two meets it in at most two distinct points. Thus this plane
contains $C$. Choose the eliminated coordinate $z$ with $\lambda_z\ne0$.

The six rational pattern coordinates~\eqref{eq:jointbar} have common denominator
$N_0N_1\ne0$ and respective integer numerators
\[
 N_1^2,\ u_1N_1,\ u_2N_1,\ N_2N_0,\ v_1N_0,\ v_2N_0,
\]
each of absolute value at most $X^2$. Expanding the quadratic $f$, its
homogeneous part is a diagonal quadratic plus the square of a linear
form. Its three diagonal numerators have denominator $N_0N_1$, and
its three linear numerators have denominator $(N_0N_1)^2$; all six
numerators have absolute value at most $12X^4$.
Substitute $z=-(\lambda_xx+\lambda_yy+\lambda_1)/\lambda_z$ and multiply by
$(N_0N_1\lambda_z)^2$. The resulting plane equation is
\begin{equation}\label{eq:jointconicplane}
 \alpha x^2+\beta xy+\gamma_c y^2+\ell_xx+\ell_yy+\ell_0=0,
\end{equation}
with integer coefficients of absolute value at most
$K=10(12X^4)^2K_{\mathrm p}^2$. Indeed, each
coefficient is a sum of products of two quantities bounded by $12X^4$
and two plane coefficients bounded by $K_{\mathrm p}$; expansion gives the stated
factor $10$. The restriction of $f$ to the plane is nonzero: otherwise
the plane would be a component of the irreducible pattern quadric.
It has degree two because it contains $C$, and hence defines $C$ in
that plane.

The pullback of the hyperplane at infinity to the normalization of $C$
has degree two. Corollary~\ref{cor:jointparity} says
that each of its positive multiplicities is even. Counting residue
degrees, there is therefore just one branch over infinity, of
multiplicity two and residue degree one. Consequently the leading
binary quadratic in \eqref{eq:jointconicplane} is a nonzero rational
scalar times the square of a rational linear form. Indeed, the
normalization map is finite and surjective, so the projective conic
meets the line at infinity only at the image of this rational point.
Its leading binary quadratic therefore has a single rational zero,
of multiplicity two. Interchanging $x$ and $y$ if necessary gives
\[
 \alpha\ne0,\qquad \beta^2-4\alpha\gamma_c=0.
\]
Set
\[
 \kappa_c=2\beta\ell_x-4\alpha\ell_y,\qquad
 \epsilon_c=\ell_x^2-4\alpha\ell_0,\qquad
 t=2\alpha x+\beta y+\ell_x.
\]
Completing the square in \eqref{eq:jointconicplane} gives
\begin{equation}\label{eq:jointconicparameter}
 t^2=\kappa_c y+\epsilon_c.
\end{equation}
Here $\kappa_c\ne0$. If it were zero, a rational point of $C$ with
parameter $t_0$ would give $\epsilon_c=t_0^2$, and the equation would
factor as $(t-t_0)(t+t_0)=0$. Integrality would then place $C$ in an
affine line, contradicting its degree.

We can therefore solve explicitly:
\[
 y=\frac{t^2-\epsilon_c}{\kappa_c},\qquad
 x=\frac{t-\beta y-\ell_x}{2\alpha},\qquad
 z=-\frac{\lambda_xx+\lambda_yy+\lambda_1}{\lambda_z}.
\]
Substitution verifies the conic and plane equations for every rational
$t$, and the displayed linear expression in $x,y$ recovers $t$.
In particular, $t$ is an integer whenever $(x,y,z)$ is an integer point.
Distinct points have distinct parameters.
Taking $\Delta_c=2\alpha\kappa_c\lambda_z$ clears all three denominators.
The inequalities
\[
 |\kappa_c|\le6K^2,\qquad |\epsilon_c|\le5K^2,\qquad
 |\beta\epsilon_c-\ell_x\kappa_c|\le11K^3
\]
give $|\Delta_c|\le12K_{\mathrm p}K^3$ and the coefficient bounds $33K_{\mathrm p}K^3$
for the three resulting quadratic numerators. For good starts one
also has $|t|\le4Ka$, although the counting argument below does not
need a parameter interval.
\end{proof}

\begin{lemma}[Properness and coefficient height]\label{lem:jointconicproper}
In the parameter of Lemma~\ref{lem:jointconicparameter}, write
\[
 E_0(t)=\frac{F_E(t)}{G_E(t)},\qquad F_E,G_E\in\mathbb Z[t],
 \qquad \gcd_{\mathbb Q[t]}(F_E,G_E)=1.
\]
The polynomials may be chosen so that
\[
 \deg G_E-\deg F_E=3,\qquad
 \max(\deg F_E,\deg G_E)\le D_0,\qquad
 \max_k\bigl(|[t^k]F_E|,|[t^k]G_E|\bigr)\le C_0a^{k_0},
\]
for absolute constants $D_0,C_0,k_0$. In particular, $G_E$ is
nonconstant. At the parameter $t$ of every good start, $G_E(t)\ne0$ and
$F_E(t)/G_E(t)=I_3(i)$ is an integer.
\end{lemma}
\begin{proof}
The explicit parametrization identifies $C$ with the affine line and
its smooth projective model with $\mathbb P^1$. Its point at infinity
has hyperplane multiplicity two. Lemma~\ref{lem:joint-infinity-orders} therefore gives
\[
 \operatorname{ord}_{t=\infty}E_0=3.
\]
The function is nonzero, since it is positive at a good start.
For a reduced rational function on $\mathbb P^1$, its order at infinity
is $\deg G_E-\deg F_E$, proving properness.

We next bound the coefficients uniformly in the component, including
after cancellation. Retain $K_{\mathrm p}=24a^3$ and
$K=10(12X^4)^2K_{\mathrm p}^2$ from the preceding proof, and define
\[
 K_{\mathrm h}=1+a+X+33K_{\mathrm p}K^3.
\]
The six rational pattern coordinates, the three coordinate polynomials
$\xi_j/\Delta_c$, and the fixed dividend $a$ have the common nonzero
denominator $N_0N_1\Delta_c$. Their ten numerator polynomials have degree
at most two and coefficient height at most $K_{\mathrm h}^4$; the common denominator
has absolute value at most $K_{\mathrm h}^4$ as well. For instance, the coordinate
numerators are $N_0N_1\xi_j$, while the numerator for $a$ is
$aN_0N_1\Delta_c$.

The cofactor formula defining $E_0$ is a quotient of two fixed
polynomials in these ten inputs, with fixed rational coefficients.
Clear those fixed coefficients once, then substitute the ten
polynomials and clear a fixed power of their common denominator.
This gives integer polynomials $\tilde n,\tilde d$, both nonzero, of bounded degree
and coefficient height at most $CK_{\mathrm h}^r$, where the degree bound, $C$
and $r$ are absolute. This assertion follows directly from polynomial
multiplication: both the number of factors and the number of summands
are fixed by the displayed cofactor formula.

Choose a primitive $\Pi\in\mathbb Z[t]$ representing
$\gcd_{\mathbb Q[t]}(\tilde n,\tilde d)$. By Gauss's lemma, $\Pi$ divides both $\tilde n$ and
$\tilde d$ in $\mathbb Z[t]$. Thus $F_E=\tilde n/\Pi$ and $G_E=\tilde d/\Pi$ are integer polynomials
with the same ratio, coprime over $\mathbb Q$, and with
degrees no larger than those of $\tilde n,\tilde d$. For a nonzero integer polynomial of
degree at most $D_0$ and coefficient height $\mathcal H$, every integer
polynomial factor has coefficient height at most
$2^{D_0}(D_0+1)\mathcal H$. To see this, the Mahler measure of an integer
cofactor is at least one, the original polynomial's Mahler measure is
at most $(D_0+1)\mathcal H$, and each coefficient of a factor is bounded by
$2^{D_0}$ times its Mahler measure. Thus cancellation preserves a
fixed polynomial bound in $K_{\mathrm h}$.

Finally, $K_{\mathrm p}=24a^3$ and $X\le10^{16}a^7$ show that $K_{\mathrm h}$ is bounded by a fixed power of $a$.
Absorbing the fixed exponents and constants gives the asserted
$C_0a^{k_0}$ bound. The original reconstruction denominator is nonzero
at every good start. Its cleared polynomial $\tilde d$ is therefore
nonzero at the corresponding parameter, and so is its factor $G_E$.
The reconstruction identity then gives $F_E(t)/G_E(t)=I_3(i)\in\mathbb Z$.
\end{proof}

\begin{lemma}[Integer values of a rational function]\label{lem:jointconicdivisors}
Let $F_E,G_E\in\mathbb Z[t]$ be coprime over $\mathbb Q$, of degrees at most
$D_0$, with $G_E$ nonconstant and coefficient heights at most
$C_0a^{k_0}$. For every $\eta>0$, the number of integers $t$ such that
$G_E(t)\ne0$ and $F_E(t)/G_E(t)\in\mathbb Z$ is $O_\eta(a^\eta)$.
The implied constant depends only on $\eta,D_0,C_0,k_0$.
\end{lemma}
\begin{proof}
The integer resultant $R=\operatorname{Res}(F_E,G_E)$ is nonzero.
The adjugate of the Sylvester matrix gives $U,V\in\mathbb Z[t]$ with
$UF_E+VG_E=R$. At any parameter in question, $G_E(t)\mid F_E(t)$, and hence
$G_E(t)\mid R$. There are $2\tau(|R|)$ possible signed values of $G_E(t)$.
For each such value $v$, the nonzero polynomial $G_E-v$ has at most
$\deg G_E\le D_0$ roots. Thus the number of parameters is at most
\[
 2D_0\tau(|R|).
\]
The Sylvester matrix has size at most $2D_0$, and its entries are
bounded by $C_0a^{k_0}$. Expanding its determinant gives
\[
 1\le |R|\le(2D_0)!\bigl(C_0a^{k_0}\bigr)^{2D_0}.
\]
The divisor estimate $\tau(n)\ll_\varepsilon n^\varepsilon$, with
$\varepsilon$ chosen in terms of $\eta$ and the fixed degree and height
exponents, proves the claim.
\end{proof}

\begin{proof}[Proof of Proposition~\ref{prop:jointconic}]
Apply the last lemma to the reconstructed invariant $E_0$. The
parameter is injective on free-node triples, and reconstruction
of the strictly increasing $q_i$ makes those triples injective in the
good starts. Each retained nonexceptional conic therefore
contributes $O_\eta(a^\eta)$ starts. The components with at most two
free-node triples of good starts and the nonretained components were already bounded
absolutely. Apply the conic estimate with $\eta/2$ and the pattern
estimate $O_{\eta/2}(a^{\eta/2}X^5)$. Since there are at most twenty-four
components per pattern, all nonexceptional conics together contribute
$O_\eta(a^\eta X^5)$.
\end{proof}

\section{The upper bound for graph degrees}\label{app:degree}

\begin{proof}[Proof of the upper bound in Lemma~\ref{lem:jointpowereddegree}]
Here only total polynomial degrees matter, not the heights of the
pattern-dependent coefficients.
Write the cofactor formula for the reconstructed invariant as
$E_0=E_{\mathrm n}/E_{\mathrm d}$, as in the proof of Lemma~\ref{lem:jointcuts}.
Each entry of the moment matrix has degree at most four in the nodes.
Expanding its $8\times8$ cofactors gives, in particular, the uniform
bound $\deg c_j\le40$, so
$\deg E_{\mathrm n},\deg E_{\mathrm d}\le240$.
Since $H_3,H_3'$ have degree at most three and the gaps have degree at
most one, the three curvature coordinates have a common denominator
\[
 \Omega=E_{\mathrm d} H_3k_1,\qquad \deg\Omega\le244,
\]
and numerators $\mathcal N_0,\mathcal N_1,\mathcal N_2$ of degree at
most $244$.  The numerator for $J$ with this denominator is
$\mathcal N_3=J\Omega$, of degree at most $247$.
The denominator is nonzero on $C$, since it is nonzero at a good start.

Put $\bar m=m+n$.  Represent the powered coordinates with denominator
$\Omega^{\bar m}$ and numerators
\[
 \mathcal N_j^m\Omega^{\bar m-m}\quad(0\le j\le2),
 \qquad
 \mathcal N_3^n\Omega^{\bar m-n}.
\]
All these numerators have degree at most $247\bar m$.
The original node coordinates, written with the same denominator,
have numerators of degree at most $244\bar m+1$.
It follows that multiplying a polynomial of degree at most $t$ in
the graph coordinates by $\Omega^{\bar mt}$ gives a polynomial in the
original nodes of degree at most $(247\bar m+1)t$.
This operation is injective on the corresponding spaces of functions,
because $\Omega$ is nonzero in $\mathbb Q(C)$.
Comparing their eventual affine Hilbert functions gives
\[
 \delta t+O(1)\le d(247\bar m+1)t+O(1),
\]
and hence the asserted upper bound.
In the unpowered case use $\Omega$ directly: the four new
numerators have degree at most $247$, and the original coordinates
have numerators of degree at most $245$.  The same comparison gives
$\delta\le247d$.
\end{proof}

\section{A reduced basis in dimension three}\label{app:reduction}

The constant $64/27=(4/3)(16/9)$ in Lemma~\ref{lem:reduce} comes from two
successive size reductions. We give the elementary argument for completeness.

\begin{proof}[Proof of Lemma~\ref{lem:reduce}]
Let $V$ be the least nonzero $|\det|$ of a triple in $\Lambda$; the triple
$(\kappa\mathbf e_1,\kappa\mathbf e_2,\kappa\mathbf e_3)$ gives $V\le\kappa^3$. A triple
$(\beta_1,\beta_2,\beta_3)$ with $|\det|=V$ is a basis: if not, some $x\in\Lambda$ has a Cramer
coordinate that is not an integer, say the first, $\det(x,\beta_2,\beta_3)/\det(\beta_1,\beta_2,
\beta_3)$, and subtracting the integer part of that coordinate times $\beta_1$ from $x$ gives a
triple with smaller nonzero determinant. Among all bases (all have $|\det|=V$) minimize
$g_1=|\beta_1|^2$, and then the Gram determinant
$g_2=g_1|\beta_2|^2-\langle\beta_1,\beta_2\rangle^2$. Put $g_3=V^2$.

Choose an integer $z_1$ with $|2(\langle\beta_1,\beta_2\rangle+z_1g_1)|\le g_1$. Then
$(\beta_2+z_1\beta_1,\beta_1,\beta_3)$ is a basis, so $g_1\le|\beta_2+z_1\beta_1|^2$, and
$g_1|\beta_2+z_1\beta_1|^2=g_2+(\langle\beta_1,\beta_2\rangle+z_1g_1)^2$. Hence
$3g_1^2\le4g_2$.

Let $\langle x,x'\rangle_1=g_1\langle x,x'\rangle-\langle\beta_1,x\rangle\langle\beta_1,x'\rangle$,
so that $\langle\beta_2,\beta_2\rangle_1=g_2$, and Sylvester's identity gives
$\langle\beta_2,\beta_2\rangle_1\langle\beta_3,\beta_3\rangle_1-\langle\beta_2,\beta_3\rangle_1^2
=g_1g_3$. Choose an integer $z_2$ with $|2(\langle\beta_2,\beta_3\rangle_1+z_2g_2)|\le g_2$ and
put $y=\beta_3+z_2\beta_2$. Then $(\beta_1,y,\beta_2)$ is a basis with first vector $\beta_1$,
so $g_2\le g_1|y|^2-\langle\beta_1,y\rangle^2=\langle y,y\rangle_1$, while
$g_2\langle y,y\rangle_1=(\langle\beta_2,\beta_3\rangle_1+z_2g_2)^2+g_1g_3$. Hence
$3g_2^2\le4g_1g_3$.

Finally choose an integer $z_3$ with $|2(\langle\beta_1,y\rangle+z_3g_1)|\le g_1$, and use the
basis $\beta_1$, $\beta_2'=\beta_2+z_1\beta_1$, $\beta_3'=y+z_3\beta_1$. The identities above
give $3g_1|\beta_2'|^2\le4g_2$ and
$4g_1g_2|\beta_3'|^2\le g_2^2+4g_1g_3+g_1^2g_2$. With the two inequalities this gives
$9g_2|\beta_3'|^2\le16g_3$. Since $|\beta_1|^2=g_1$, the product of the bounds
$|\beta_2'|^2\le\frac43g_2/g_1$ and $|\beta_3'|^2\le\frac{16}9g_3/g_2$ is the claim for the
basis $\beta_1,\beta_2',\beta_3'$; rename $\beta_2',\beta_3'$ as $\beta_2,\beta_3$.
Replacing $\beta_1$ by $-\beta_1$ if necessary makes the determinant positive. For the
coordinates, $\xi_1V=\det(x,\beta_2,\beta_3)$, and Hadamard's inequality gives
$\det(x,\beta_2,\beta_3)^2\le|x|^2|\beta_2|^2|\beta_3|^2$, so
$27|\beta_1|^2\xi_1^2V^2\le64V^2|x|^2$; likewise for $\xi_2$ and $\xi_3$.
\end{proof}

\section{Certificates for the polar bound}\label{app:polarcertificates}

\subsection{The quartic initial form at the base point}

At $(\rho,u,v,\rho',u',v')=(2,5,-3,3,6,-5)$ the polar base point is
$x_*=(-1/7,-3/28,-1/2)$ and $\nabla f(x_*)=(-4,9,-5)$.
Use tangent coordinates $(s,t)$ with first-order displacement
\[
 \delta x=((9s-5t)/4,s,t).
\]
Thus the two axes used in Lemma~\ref{lem:polarbase} have tangent vectors
$(9/4,1,0)$ and $(-5/4,0,1)$. Direct substitution in the Newton quartic,
its remainder and the norm polynomial gives
\begin{align*}
 P_{4,*}(\xi)&=\xi^4-\frac{20}{7}\xi^3+\frac{757}{112}\xi^2
                         -\frac{669}{112}\xi+2,\\
 R_1(\xi;s,t)&=\frac{77s+7t}{14}
       +\frac{-56s-124t}{14}\xi
       +\frac{-45s+157t}{14}\xi^2
       +\frac{28s-28t}{14}\xi^3,\\
 \Psi_1(s,t)&=\frac{3s+85t}{28}.
\end{align*}
The resultant is
\[
 N_4(s,t)=\frac{
 32233977s^4-95047956s^3t+147873750s^2t^2
             -111199572st^3+31702169t^4}{6272}.
\]
In particular,
\[
 N_4(1,0)\Psi_1(0,1)^4-N_4(0,1)\Psi_1(1,0)^4
 =\frac{26291112152499}{60236288}.
\]
These identities specify the frame and first derivatives needed to reproduce
the local certificate. The realizing matrix in Lemma~\ref{lem:polarbase}
certifies a nonzero polynomial on the decomposable cone; no orbit with these
parameters is required.

\subsection{The irreducible sextic}

The specialization $E_\gamma=24E_0$ at
$(\gamma_1,\gamma_2,\gamma_3)=(-7,8,-4)$ reduces modulo $71$ to a nonzero
scalar multiple of
\[
 q(T)=T^6+20T^5-14T^4+15T^3-28T^2+5T-4.
\]
More precisely, $E_0=57q$ in $\mathbb F_{71}[T]$.
The following certificates prove that $q$ is irreducible. Repeated squaring
modulo $q$ gives
\begin{align*}
 T^{71^2}&\equiv r_2(T):=36T^5+48T^4+49T^3+53T^2+36T+42,\\
 T^{71^3}&\equiv r_3(T):=37T^5+22T^4+60T^3+13T^2+66T+35.
\end{align*}
For checking the iteration one may first use
\[
 T^{71}\equiv14T^5+33T^4+31T^3+34T^2+48T\pmod q
\]
and raise each remainder to the $71$st power. Define
\begin{align*}
 U_2&=20T^4+34T^3+31T^2+13T+24,\\
 V_2&=31T^5+61T^4+42T^3+61T^2+39T+4,\\
 U_3&=25T^4+11T^3+53T^2+59T+48,\\
 V_3&=7T^5+30T^4+65T^3+64T^2+25T+40.
\end{align*}
Direct multiplication in $\mathbb F_{71}[T]$ verifies
\[
 U_2q+V_2(r_2-T)=1,\qquad U_3q+V_3(r_3-T)=1.
\]
Thus $q$ has no irreducible factor of degree one or two, since such a factor
would divide $T^{71^2}-T$, and none of degree three, since it would divide
$T^{71^3}-T$. Every reducible sextic has an irreducible factor of degree at
most three, so $q$ is irreducible. Reduction modulo $71$ preserves degree,
and Gauss's lemma proves $E_0$ irreducible over $\mathbb Q$.

\section*{Acknowledgements}
The argument was developed with the help of GPT-6-Astra. This includes the use of
orthogonal polynomials and adjacent labels behind Section~\ref{sec:patterns}, the first proof of a
bound $a^{3/10-\delta}$ through resultants, the use of Lemma~\ref{lem:multid} to bound the
number of states per pattern, the sign identities of Lemma~\ref{lem:signs} and the slice count behind
the $11/37$ bound, the observation that the pattern integers form a Pl\"ucker vector, and the
proposal to weight minor vectors by the shortest vectors of their lattices, behind the $45/152$
bound. Subsequent AI-assisted
work developed and formalized the joint-invariant graphs, their normalization
and valuation arguments, the powered-coordinate optimization, and the external
affine-curve counting theorem. OpenAI Codex assisted with this formalization
and with automated statement, dependency and submission checks. These checks
are not independent human peer review. The unconditional results, including
the polar refinement and its arithmetic sieve, and the stated conditional
implications were verified in Lean~4 against Mathlib. We thank the Mathlib community, and the authors of \lean{openai/math}
for the formalized divisor bound.

\begin{thebibliography}{99}

\bibitem{BV}
O.~Baraskar and I.~Vukusic,
\emph{Bounds for sets of remainders},
\href{https://arxiv.org/abs/2508.20853}{arXiv:2508.20853} (2025).

\bibitem{BT}
R.~de~la~Bret\`eche and G.~Tenenbaum,
\emph{Note on the mean value of the Erd\H{o}s--Hooley Delta-function},
Acta Arith.\ \textbf{219} (2025), no.~4, 379--394.

\bibitem{CCDN}
W.~Castryck, R.~Cluckers, P.~Dittmann and K.~H.~Nguyen,
\emph{The dimension growth conjecture, polynomial in the degree and without logarithmic factors},
Algebra Number Theory \textbf{14} (2020), no.~8, 2261--2294,
\href{https://doi.org/10.2140/ant.2020.14.2261}{doi:10.2140/ant.2020.14.2261}.

\bibitem{CDHNV}
R.~Cluckers, P.~D\`ebes, Y.~I.~Hendel, K.~H.~Nguyen and F.~Vermeulen,
\emph{Improvements on dimension growth results and effective Hilbert's irreducibility theorem},
Forum Math.\ Sigma \textbf{13} (2025), e153,
\href{https://doi.org/10.1017/fms.2025.10096}{doi:10.1017/fms.2025.10096}.

\bibitem{CP}
Z.~Chase and M.~Pandey,
\emph{On the length of Pierce expansions},
\href{https://arxiv.org/abs/2211.08374}{arXiv:2211.08374} (2022).

\bibitem{ES}
P.~Erd\H{o}s and J.~O.~Shallit,
\emph{New bounds on the length of finite Pierce and Engel series},
S\'em.\ Th\'eor.\ Nombres Bordeaux (2) \textbf{3} (1991), 43--53.

\bibitem{Lean}
L.~de~Moura and S.~Ullrich,
\emph{The Lean~4 theorem prover and programming language},
in: Automated Deduction -- CADE 28, Lecture Notes in Comput.\ Sci.\ \textbf{12699},
Springer, 2021, 625--635.

\bibitem{Mathlib}
The mathlib Community,
\emph{The Lean mathematical library},
in: Proceedings of the 9th ACM SIGPLAN International Conference on Certified Programs and
Proofs (CPP 2020), 367--381.

\bibitem{NT}
M.~Nair and G.~Tenenbaum,
\emph{Short sums of certain arithmetic functions},
Acta Math.\ \textbf{180} (1998), 119--144.

\bibitem{OAIMath}
OpenAI,
\texttt{openai/math}, \url{https://github.com/openai/math}.

\bibitem{OPG}
J.~O.~Shallit,
\emph{A discrete iteration related to Pierce expansions},
Open Problem Garden,
\url{https://www.openproblemgarden.org/op/a_discrete_iteration_related_to_pierce_expansions}.

\bibitem{Pierce}
T.~A.~Pierce,
\emph{On an algorithm and its use in approximating roots of algebraic equations},
Amer.\ Math.\ Monthly \textbf{36} (1929), 523--525.

\bibitem{Shallit86}
J.~O.~Shallit,
\emph{Metric theory of Pierce expansions},
Fibonacci Quart.\ \textbf{24} (1986), 22--40.

\bibitem{ShallitTalks}
J.~O.~Shallit,
\emph{Talks}, \url{https://cs.uwaterloo.ca/~shallit/talks.html}
(accessed October 2026).

\end{thebibliography}

\end{document}
